\documentclass{article}

\PassOptionsToPackage{numbers,compress}{natbib}

% \usepackage[main, final]{neurips_2026}
\usepackage{neurips_2026}

\usepackage[utf8]{inputenc}
\usepackage[T1]{fontenc}
\usepackage{hyperref}
\usepackage{url}
\usepackage{xcolor}
\usepackage{booktabs}
\usepackage{amsfonts}
\usepackage{nicefrac}
\usepackage{multirow}
\usepackage{microtype}

\usepackage{graphicx}
\usepackage{float}
\usepackage{comment}
\usepackage{enumitem}
\usepackage{subfigure}
\usepackage{dsfont}

\usepackage{amsmath,amssymb,amsthm,mathtools,mathrsfs}
\usepackage{bbm}
\usepackage{bm}

\usepackage[capitalize,noabbrev]{cleveref}
\usepackage[textsize=tiny]{todonotes}
\usepackage[ruled,vlined]{algorithm2e}

\theoremstyle{plain}
\newtheorem{theorem}{Theorem}[section]
\newtheorem{proposition}[theorem]{Proposition}
\newtheorem{lemma}[theorem]{Lemma}
\newtheorem{remark}[theorem]{Remark}
\newtheorem{corollary}[theorem]{Corollary}
\theoremstyle{definition}
\newtheorem{definition}[theorem]{Definition}
\newtheorem{assumption}[theorem]{Assumption}


\title{Watermarkable Multi-Draft Speculative Sampling \\
via Poisson Processes with Drafter Invariance}


\author{
  Yanxiao Liu*
  % \thanks{Use footnote for providing further information
  %   about author (webpage, alternative address)--\emph{not} for acknowledging
  %   funding agencies.} 
    \\ 
  Imperial College London\\
  London, United Kingdom \\
  \texttt{y.liu2@imperial.ac.uk} \\
  % examples of more authors
  \And
  Sicheng Wan* \\
  University of Washington \\
  Seattle, Washington, USA \\
  \texttt{sicheng4@uw.edu} \\
  \And
  Zhan Gao \\
  Imperial College London\\
  London, United Kingdom \\
 \texttt{?????@imperial.ac.uk} 
  \AND
  Deniz G\"und\"uz \\
  Imperial College London\\
  London, United Kingdom \\
  \texttt{d.gunduz@imperial.ac.uk} \\
}


\begin{document}


\maketitle


\begin{abstract}
Large language models (LLMs) based on autoregressive generation have achieved state-of-the-art performance across a wide range of tasks, motivating two important aspects of deployment: inference efficiency and output provenance, which can be %are 
tackled by speculative sampling and watermarking, respectively.
However, recent works have shown that, combining these two goals is highly nontrivial and can even be impossible in some cases. 
In this work, we circumvent this fundamental tradeoff by designing a novel multi-draft speculative sampling algorithm based on Poisson processes. 
The proposed %Our %sampling 
algorithm has strong efficiency on its own and, more interestingly, is naturally \emph{watermarkable}: we can embed an unbiased watermark without degrading speculative acceptance. 
Moreover, our algorithm %method 
is based on an exact list-coupling-without-communication scheme, which yields a \emph{drafter invariance} property that benefits both sampling and watermarking.
It is the first multi-draft, drafter-invariant speculative sampling scheme that maintains both watermark strength and sampling efficiency, and we experimentally verify its strong performance in both aspects.
\end{abstract}


\section{Introduction}
\label{sec::intro}
\vspace{-2pt}

Large language models (LLMs) have rapidly become a central component of modern generative AI, achieving strong performance in dialogue, reasoning, code generation, and open-ended content creation. Their probabilistic nature based on autoregressive sampling enables them to generate diverse and fluent responses~\citep{brown2020language, chowdhery2023palm, touvron2023llama, ouyang2022training, achiam2023gpt}. 
However, as LLM-generated text becomes increasingly fluent and difficult to distinguish from human-written content, concerns about accountability, provenance, misuse, and intellectual property have attracted increasing attention~\citep{wu2025survey}. 
For instance, ChatGPT-generated scientific abstracts were shown to be difficult for human reviewers to distinguish from human-generated ones, yet they may contain fabricated data~\citep{gao2023comparing, shumailov2024ai}. % weidinger2022taxonomy
Hence, watermarking has become an important, and potentially necessary, tool for tracing the provenance of LLM-generated text by embedding statistical signals into the sampling process while preserving output quality~\citep{kirchenbauer2023watermark, aaronson2023watermarking, kuditipudi2023robust}. 


Apart from provenance, the efficiency of autoregressive generation is another bottleneck in the system pipeline. 
Speculative decoding, also known as speculative sampling, has emerged as a promising approach to accelerate content generation without changing the output distribution of the target model~\citep{chen2023accelerating, leviathan2023fast}. 
However, speculative sampling is not independent of watermarking.
\textcolor{teal}{On the contrary, recent works~\citep{hu2024inevitable, he2026improving} show that combining the two is highly nontrivial: in~\citep{hu2024inevitable}, a ``no-go'' theorem was proved, showing that one cannot embed a watermark without degrading the performance of speculative sampling, thereby establishing a fundamental tradeoff between watermark strength and sampling efficiency.
This tradeoff was improved in~\citep{he2026improving} by considering alternative watermark metrics. }


\textcolor{teal}{In this work, we ask the following question: can we somehow sidestep the no-go theorem~\citep{hu2024inevitable}, i.e., design a watermarking mechanism on top of speculative sampling without degrading sampling efficiency, thereby improving the tradeoff between watermark strength and sampling efficiency?}


We answer this question in the \textbf{affirmative}: we design a novel multi-draft speculative sampling algorithm that achieves strong sampling performance, is naturally watermarkable, and, importantly, has watermark strength that does not degrade as the acceptance rate increases. 
Our algorithm relies on Poisson processes~\citep{li2021unified, maddison2016poisson} and yields a \emph{list-coupling without communication} scheme. This design %, which 
endows our sampling algorithm with a \emph{drafter-invariance} property~\citep{daliri2025coupling, rowan2025list}, which becomes even more desirable and essential in the context of watermarking. 
In summary, our contributions are: 
\begin{itemize}
    \item \textbf{Multi-draft speculative sampling}: We design a novel multi-draft speculative sampling algorithm that generates and verifies multiple draft continuations through a list-coupling scheme. 
    Our algorithm utilizes Poisson processes~\citep{li2021unified,maddison2016poisson}, which naturally induce coupling without communication and yield drafter invariance in LLM decoding~\citep{daliri2025coupling,rowan2025list}. 
    Compared with~\citep{rowan2025list}, we achieve better sampling acceptance rate on various tasks, while providing a stronger form of drafter invariance that is useful for watermarking. 
        
    \item \textbf{Efficiency-preserving watermarking}: Except desirable sampling efficiency, our algorithm is naturally \emph{watermarkable} with %by 
    the use of Poisson processes, which can utilize %use 
    a key to embed recoverable statistical signal while maintaining the distribution \emph{unbiased}. 
    By building the watermark directly into sampling through target-side keyed randomness, we maintain watermark strength while preserving sampling efficiency, thereby \emph{circumventing} the inevitable tradeoff in~\citep{hu2024inevitable} and achieving higher watermark strength.
    Compared with~\citep{hu2024inevitable, he2026improving}, our multi-draft  algorithm achieve better sampling efficiency and watermark strength. 
\end{itemize}

\textbf{Notations}. 
We use uppercase and lowercase letters to denote random variables and realizations, respectively. 
The logarithm is on base $2$. 
We use $x_{i:j}$ to denote the sequence $x_i, x_{i + 1}, \ldots, x_j$, where $x_i$ is the $i$-th entry, and $x_{:j} := x_{1:j}$. 
$[N]$ denotes $\{1,\ldots,N\}$. 
We assume that all mentioned random variables lie in a Polish space with Borel $\sigma$-algebra and that all functions are measurable. 
The Lebesgue measure over $\mathbb{R}$ is denoted as $\lambda$. 
For probability measures $P, Q$, $P\ll Q$ denotes that $P$ is absolutely continuous with respect to $Q$, and $\mathrm{d}P(\cdot)/ \mathrm{d}Q$ denotes the Radon-Nikodym derivative. 
% The total variation distance between $P$ and $Q$ is denoted as $\Vert P - Q\Vert_\mathrm{TV} := \sup_{A\subseteq \mathcal{X}\text{ measurable}} \big| P(A) - Q(A)\big|$.  



\section{Related Works}
\label{sec::related}


\paragraph{LLM Decoding and Speculative Sampling.}

Standard autoregressive decoding methods, including top-$k$ sampling~\citep{radford2019language, fan2018hierarchical}, nucleus sampling~\citep{holtzman2020curious}, and the permute-and-flip decoder~\citep{zhao2024permute}, generate tokens sequentially; and therefore, suffer from the inherent latency of one target-model call per token~\citep{stern2018blockwise}. 
Speculative sampling~\citep{chen2023accelerating, leviathan2023fast} addresses this bottleneck by using a smaller draft model to propose several future tokens, while the larger target model verifies the proposed continuations in parallel, and a rejection-correction step preserves the exact target-model distribution. 
The efficiency gain depends on how well the draft distribution matches the target distribution: more accurate proposals lead to higher acceptance rates and longer verified blocks. 
Recent works further improve this paradigm through optimal-transport-based token selection~\citep{sun2023spectr}, tree-based inference and verification~\citep{miao2024specinfer}, and multi-draft architectures~\citep{khisti2024multi, rowan2025list}. 



\paragraph{Watermarking.}


LLM watermarking aims to embed detectable statistical signals into generated text while preserving fluency and utility. 
Early logit-bias methods pseudorandomly select green-list tokens and detect the resulting bias by hypothesis testing~\citep{kirchenbauer2023watermark}, while later works improve robustness, accessibility, and probability balance~\citep{wu2023resilient, giboulot2024watermax, park2026watermod}. 
A parallel line studies distortion-free or unbiased watermarking, where the sampler maintains the original distribution while correlating the output with secret randomness~\citep{kuditipudi2023robust, hu2023unbiased, christ2024undetectable}. 
These works connect watermarking to information-theoretic~\citep{moulin2003information} and statistical questions, including detection power, false-positive control, entropy dependence, optimal tests, and couplings~\citep{li2025statistical, he2025empirical, tsur2025optimized, tsur2025heavywater}. 
Other directions include robust or multi-bit attribution~\citep{zhao2023provable, boroujeny2024multi, qu2025provably}, and efficient identification of watermarked segments~\citep{zhao2025efficiently}. 
In recent works~\citep{hu2024inevitable, he2026improving} the fundamental tradeoff between speculative sampling and watermarking has been discussed. 

\paragraph{Poisson Functional Representation (PFR).}

PFR is an exact sampling scheme introduced in~\citep{li2018strong}, with related Poisson process constructions appearing in Monte Carlo and $\mathrm{A}^*$ sampling~\citep{maddison2014sampling, maddison2016poisson}.
It selects a sample from a target distribution through a shared Poisson process and a proposal distribution, and has been applied to various information-theoretic settings~\citep{li2021unified, liu2025one, hentila2024communication, khisti2024unequal}, hypothesis testing~\citep{guo2024hypothesis}, and machine learning~\citep{li2020minimax, he2024accelerating}.
It can be naturally generalized to a list-coupling scheme~\citep{li2021unified, li2019pairwise} by the mapping theorem~\citep{last2017lectures} for Poisson processes, and achieves the best known bound on coupling without communication~\citep{daliri2025coupling, rowan2025list}. 
Related sampling constructions can be found in~\citep{flamich2024greedy, liu2024universal, li2024channel}. 




\section{Speculative Sampling}
\label{sec::spec}


We first briefly review the technical background  of speculative sampling~\citep{chen2023accelerating,leviathan2023fast,sun2023spectr}. 


Given volcabulary $\mathcal V$ and a context $x_{:t}:=(x_1,x_2, \ldots,x_t)$, an autoregressive language model $\mathcal{M}_b$ generates the next token by sampling from the conditional distribution $\mathcal{M}_b(\cdot| x_{:t})$ under temperature sampling~\citep{ackley1985learning,ficler2017controlling}. 
\emph{Speculative sampling} accelerates this process by introducing a smaller and cheaper \emph{draft model} $\mathcal{M}_s$, which proposes several candidate tokens before the target model verifies them. 
Given a prefix $x_{:t}$, one iteration of the speculative sampling consists of the following three stages:
\begin{enumerate}
    \item  \hspace{-3.9pt} \textbf{Draft Stage.} The draft model $\mathcal{M}_s$ sequentially samples $L$ candidate tokens $\tilde{x}_{t+1},\ldots,\tilde{x}_{t+L}$. For each position $i=1,\ldots,L$, we record the draft distribution $\mathcal{M}_s(\cdot| x_{:t}, \tilde{x}_{t+1:t+i-1})$. 

    \item \hspace{-3.9pt} \textbf{Computation Stage.} Conditioned on the drafted context, the target model $\mathcal{M}_b$ evaluates the corresponding conditional distributions, %as follows, 
    potentially in parallel:
    \[
    \mathcal{M}_b(\cdot| x_{:t}),
    \mathcal{M}_b(\cdot| x_{:t}, \tilde{x}_{t+1}), \ldots,
    \mathcal{M}_b(\cdot| x_{:t}, \tilde{x}_{t+1:t+L}), 
    \]
    
    \item \hspace{-3.9pt} \textbf{Selection Stage.} Using the distributions from $\mathcal{M}_s$ and $\mathcal{M}_b$, we accept the %a 
    longest prefix of the drafted tokens until the first rejection, and a \emph{correction token} is sampled from a pre-designed residual distribution. If all $L$ drafts are accepted, one additional token from the target is sampled. Then, decoding proceeds to the next iteration. 
\end{enumerate}

 
Crucially, the above acceptance-correction procedure ensures that the generated tokens follow \emph{exactly} the same distribution $\mathcal{M}_b(\cdot | x^t)$ as direct autoregressive sampling from the target model $\mathcal{M}_b$, i.e., decoding efficiency is improved without changing the target model's output law. 
The exactness is achieved via a recursive token-level maximal coupling~\citep{sun2023spectr,chen2023accelerating,leviathan2023fast}, as shown in Algorithm~\ref{alg::max_coupling}, where we use $P$ and $Q$ to denote the target model and draft model distributions, respectively, and achieve
\[
\mathrm{Pr}(X = Y) = \sum_{x\in\mathcal{V}} \min (P(x), Q(x)) = 1 - \Vert P - Q\Vert_\mathrm{TV}, 
\]
i.e., the closer $P$ and $Q$ are, the higher the probability that the draft token will be accepted. 
Together with the extra sampled token, there could be at most $L + 1$ tokens %been 
generated in each iteration, and therefore, the speedup is up to $(L+1)$ times if the decoding time of the draft model is negligible. 


 
\begin{figure}[H]
% \begin{figure}[t]
\centering
\begin{minipage}[t]{.51\textwidth}

\begin{algorithm}[H]
\SetAlgoLined
\DontPrintSemicolon
\SetKwIF{If}{ElseIf}{Else}{if}{:}{else if}{else}{end}
\SetKwInOut{Input}{Input}\SetKwInOut{Output}{Output}
\textbf{Input:} Target distribution $P$, draft distribution $Q$, and draft sample $X \sim Q$. \;

Compute the residual $P_{\mathrm{res}}$: for any $x\in \mathcal{V}$,
\[
P_{\mathrm{res}}(x)=\frac{P(x)-\min\{P(x),Q(x)\}}{1-\sum_y \min\{P(y),Q(y)\}}.
\]
Sample $U \sim \mathrm{Unif}(0,1)$
\;\vspace{2pt}
\If{$U \leq \min\left(1, {P(X)}/{Q(X)}\right)$}{
\Return{$Y=X$} \tcp*[r]{Accept}
}
\Return{$Y \sim P_{\mathrm{res}}$} \tcp*[r]{Correction}
\caption{Token-level sampling~\citep{sun2023spectr, chen2023accelerating, leviathan2023fast}}
\label{alg::max_coupling}
\end{algorithm}

\end{minipage}%
\hfill
\begin{minipage}[t]{.47\textwidth}

\begin{algorithm}[H]
\SetAlgoLined
\DontPrintSemicolon
\SetKwIF{If}{ElseIf}{Else}{if}{:}{else if}{else}{end}
\SetKwInOut{Input}{Input}\SetKwInOut{Output}{Output}
\textbf{Input:} Target distribution $P$ and draft distribution $Q$\;
\For{$i\in \mathcal{V}$}{
    $E_i\sim\mathrm{Exp}(1)$\;
}
$X \gets \arg\min_{i\in \mathcal{V}} E_i/Q(i)$\;
$Y \gets \arg\min_{i\in \mathcal{V}} E_i/P(i)$\;
\If{$X = Y$}{
\Return{$Y = X$} \tcp*[r]{Accept}
}
\Return{$Y$} \tcp*[r]{Correction}
\caption{Token-level sampling via exponential races~\citep{kobus2025speculative}}
\label{alg:ESGD_token}
\end{algorithm}
\end{minipage}
\end{figure}


\subsection*{Speculative Sampling via Exponential Races}

In~\citep{kobus2025speculative}, a new speculative sampling scheme (Algorithm~\ref{alg:ESGD_token}) was proposed, whose token-level coupling is based on exponential random variables rather than rejection sampling. Each token in $\mathcal{V}$ is assigned an exponential random variable, which is divided %. Dividing %, and dividing it 
%this value 
by the %token's 
probability of that token to produce a \emph{score}; 
%yields a \emph{score}, and 
the token with the smallest score wins the race.
The winner follows the desired distribution exactly: $X \sim Q$ and $Y \sim P$.
See~\citep{rowan2025list} for a multi-draft variant %algorithm 
with a similar spirit. 


One crucial property of Algorithm~\ref{alg:ESGD_token} is that the target and the draft models do not communicate, and the \emph{coupling without communication} yields a desirable \emph{drafter invariance} property~\citep{daliri2025coupling, rowan2025list}, see Section~\ref{sec::drafter_invar} for more discussions. 
In the next section, 
we formally present a single-draft version of our algorithm that %shares a similar idea, but uses 
leverages Poisson processes instead of exponential random variables for speculative sampling, and extend it to a multi-draft variant in Section \ref{sec::multi_draft}. 



\section{Watermarkable Speculative Sampling with Draft-Invariance}\label{sec:SingleVersion}


We present a novel speculative sampling scheme based on Poisson processes, which naturally embeds watermark during the sampling process based on pseudorandomness. 
It can be easily extended to a multi-draft variant %version 
to further accelerate autoregressive generation without compromising quality, still by the properties of Poisson processes, which %as 
will be elaborated in Section~\ref{sec::multi_draft}. 
Interestingly, our scheme is also drafter-invariant since there is no communication between the target and draft models. This property is particularly helpful for watermarking because it ties the watermark to the target-side keyed randomness rather than drafter-specific behavior; hence changing the drafter (e.g., via a model update) does not affect the watermark carried by the generated text. 
Before proceeding, %presenting our main algorithm, 
we first introduce two key components and then present our main algorithm. %some background. 

\subsection{Gumbel Watermark}
\label{subsec::gumbel_water}

A good watermark should be detectable from the generated text while preserving output quality, and, ideally, the model's marginal token distribution. 
\textit{Unbiased} (or distortion-free) watermarking schemes use pseudorandomness to couple generated tokens with secret keys, so that the watermark can be detected through statistical tests without significantly altering the marginal text distribution: 
suppose $\tilde{P}$ is a base distribution over $\mathcal{V}$, an unbiased watermark is a decoding rule such that, for each seed $\zeta$, it induces a watermarked distribution $\tilde{P}_\zeta$, which preserves $\tilde{P}$ after averaging over the seed, 
\[
\mathbf{E}_{\zeta} [\tilde{P}_\zeta(v)] = \tilde{P}(v), \qquad \forall v\in\mathcal{V}.
\]
The Gumbel watermark~\citep{aaronson2023watermarking}, which uses the Gumbel-max trick~\citep{gumbel1954statistical} and lets the seed $\zeta$ determine the Gumbel noise, is unbiased, and we leverage a similar idea for %, and the idea is similar to 
our watermarking scheme. 
Specifically, as explained in~\citep{zhao2024permute}, \textcolor{teal}{let $u_t(\cdot | x, y_{t-1}): \mathcal{V}\to \mathbb{R}$ be a real-valued function (a.k.a. \emph{logits}) that encodes the model's preferences on words and $\mathsf{T}$ be the temperature,} the softmax sampling is equivalent to
%\begin{equation}
%    y_t \sim p(y) = {\exp\big(u(y|x,y_{1:t-1})/T\big)}
%    \Big/ \, \Big({\sum_{\tilde{y}} \exp\big(u(\tilde{y}|x,y_{1:t-1})/T\big)\Big)} \label{eq::softmax}
%\end{equation}
%where $x$ is the \emph{prompt} for the current task, $y_{1:t-1}$ is the \emph{prefix} that include all previously generated words up to time $t$, and 
%$u(\cdot | x, y_{1:t-1}): \mathcal{V}\to \mathbb{R}$ is a real-value function (a.k.a. \emph{logits}) that encodes the model's preferences on words from the vocabulary $\mathcal{V}$. 
%By the Gumbel-max trick, \eqref{eq::softmax} is equivalent to 
\[
y_t =  \mathrm{argmax}_{y\in\mathcal{V}}  \big({u_t(y)} \big/ \mathsf{T} \big) + G_t(y),\quad G_t(y)\sim\textrm{Gumbel}(0,1) \text{ i.i.d for each } t,y,
\]
where  the Gumbel noise~\citep{gumbel1954statistical} can be generated by $\textrm{Gumbel}(0,1) \sim -\log(\log(1/\textrm{Unif}([0,1])))$. 



A random $r_t\sim  (\textrm{Unif}([0,1]))^{|\mathcal{V}|}$ can be represented by $G_t(y) = - \log (-  log (r_t(y)))$. 
We can replace $\textrm{Unif}([0,1])$ by a pseudo-random function $r_t(y) = F_{y_{t-m:t-1},\mathsf{k}}(y)$ with prefix length $m$, and employ
\begin{equation}
    y_t =  \mathrm{argmax}_{y\in\mathcal{V}}  \big({u_t(y)} \big/ \mathsf{T} \big) -\log(-\log(r_t(y))). \label{eq::decoding_pseudo}
\end{equation}
Conditioned on a secret key $\mathsf{k}$, $r_t(y)$ is a deterministic function, but over the distribution of $\mathsf{k}$, $r_t(y)$ is computationally indistinguishable from %sampled 
samples drawn from a truly i.i.d. uniform distribution. Therefore, %, and therefore we say 
the distribution of $y_t$ from~\eqref{eq::decoding_pseudo} is considered the same as the unwatermarked one. 

The detector, which has the secret key $\mathsf{k}$, can calculate the watermark score $ \sum^n_{t = m + 1}-\log\big(1 - r_t(y_t) \big)$. 
If $y_{1:n}$ is not watermarked, the score has an expectation $n - m$; otherwise, the expectation becomes larger and the watermark is detected. 
See Appendix~\ref{subapp::gumbel_water} for details. %a detailed introduction. 

Our algorithm relies on Poisson processes and also incorporates the use of pseudo-random functions. 


\subsection{Poisson Functional Representation (PFR)}
\label{subsec::PFR}


We leverage a sampling scheme called PFR~\citep{li2018strong, maddison2014sampling}, which 
has found a wide range of applications~\citep{li2021unified, theis2022algorithms, liu2024universal, hentila2024communication, yang2023introduction, song2025downlink, zamani2023privacy, theis2022lossy, zhou2026dual} in recent years.
\begin{definition}[Poisson functional representation (PFR)~\citep{li2018strong, maddison2014sampling}] \label{def::PFR}
Let $(T_{i})_{i}\sim \mathrm{PP}(1)$ be a Poisson process of rate $1$ (i.e., $T_{1},T_{2}-T_{1}\ldots \stackrel{\mathrm{iid}}{\sim} \mathrm{Exp}(1)$),
independent of $Z_{i}\stackrel{\mathrm{iid}}{\sim} \nu$ for $i=1,2,\ldots$. 
Then $(Z_{i},T_{i})_{i}$
is a Poisson process with intensity measure $Q\times\lambda_{[0,\infty)}$ \citep{last2017lectures}, where $\lambda_{[0,\infty)}$ is the Lebesgue measure over $[0,\infty)$.
Fix distribution $P$ over $\mathcal{Z}$ that is absolutely continuous with respect to $Q$. Let 
\begin{equation}
    \tilde{T}_{i}:=T_{i} \cdot \big({\mathrm{d}P}/{\mathrm{d}Q}(Z_{i})\big)^{-1}.
    \label{eq:PFR}
\end{equation}
By the mapping theorem \citep{last2017lectures}, $(Z_{i},\tilde{T}_{i})$ is a Poisson process with intensity measure $P\times\lambda_{[0,\infty)}$. 
PFR selects the point $Z=Z_{K}$ with the smallest associated $\tilde{T}_{K}$, 
i.e., let $K:=\mathrm{argmin}_{i}\tilde{T}_{i}$ and $Z := Z_K$.
\end{definition}
\begin{remark}
    By drawing a sequence $(Z_i)_i$ from the \emph{reference} distribution $Q$ and a sequence of \emph{times} $(T_i)_i$, the PFR selects a sample following the \emph{target} distribution $P$. 
    The sample $Z_i$ with the smallest $T_i$ only follows $Q$; to obtain a sample from $P$, we inflate the time by factor $({\mathrm{d}P}/{\mathrm{d}Q}(Z_{i}))^{-1}$. 
    Although PFR requires infinite shared randomness, it can be implemented easily in practice~\citep{theis2022algorithms}. 
\end{remark}



\subsection{Watermarkable Speculative Sampling via Poisson Processes}
\vspace{-5pt}

{\color{teal}

\paragraph{Core idea.} 
We propose a watermarkable speculative sampling scheme in which watermarking and acceleration are unified through a single keyed Poisson coupling. 
Every emitted token is target-distributed and watermarked, while the draft model only controls how many such samples can be emitted. 
The key distinction is that watermarking is performed at the \emph{coupling level}, rather than through a post-hoc reweighting of marginals. This architectural shift is what allows our scheme to escapte the no-go regime of~\citep{hu2024inevitable}, whose impossibility result applies to schemes that compose marginal reweighting with subsequent speculative coupling. 
}

% We integrate the ideas in Sections~\ref{subsec::gumbel_water} and \ref{subsec::PFR} to design a single-draft version of our algorithm. %first present a single-draft version of our algorithm that connects the ideas in Section~\ref{subsec::gumbel_water} and \ref{subsec::PFR}. 
% , and \textcolor{red}{we show our algorithm achieves the fundamental trade-off curve of watermark strength and sampling efficiency???}


Algorithm~\ref{alg:single_PFR_Spec_Water_simplified} is a sketch of our single-draft speculative sampling with watermark -- see a detailed implementation %is 
in Algorithm~\ref{alg:fwss_pfr}. 
The key step is fixing a Poison process $\Pi_\mathsf{k}$ with key $\mathsf{k}$, where %; 
PFR$(\bar{P}(\cdot | c),\Pi )$ denotes the sample selected by PFR (Definition~\ref{def::PFR}) with target distribution $\bar{P}(\cdot | c)$ and Poison process $\Pi$. %, and we 
We also label the Poisson processes by context $c$, denoted by $G(\mathsf{k}, c)$. 
% \begin{algorithm}[H]
\begin{algorithm}[tbp]
\SetAlgoLined
\DontPrintSemicolon
\SetAlgoNoEnd
\SetKwIF{If}{ElseIf}{Else}{if}{:}{else if}{else}{end}
\SetKwInOut{Input}{Given}

\Input{lookahead $K$, output length $N$, 
% target model $P$, draft model $Q$,
initial prompt $w_{1:n}$, Poisson generator $G$, secret key $\mathsf{k}$}

\While{$n < N$}{
    $h \gets w_{1:n}$; $L \gets \min\{K,\,N-n\}$; $accepted \gets \textbf{true}$; \;
    % \tcp{Draft a speculative block from keyed shared Poisson sources}
    \For{$s=1$ \KwTo $L$}{
        $c_s \gets h \,\|\, \tilde w_{1:s-1} $; 
        \textcolor{blue}{$\Pi_s \gets G(\mathsf{k}, c_s)$} \tcp*[r]{\textcolor{blue}{keyed Poisson source for watermark}}
        $\tilde w_s \gets \textsc{PFR}(Q(\cdot | c_s), \Pi_s)$\;
    }

    \For(\tcp*[f]{evaluation in parallel}){$s=1$ \KwTo $L$}{
        $y_s \gets \textsc{PFR}(P(\cdot | c_s), \Pi_s)$\;
        $a_s \gets \mathds{1} \left\{y_s = \tilde w_s\right\}$; 
        % \tcp*[r]{watermark-bearing acceptance bit}
        $w_{n+1} \gets y_s$; $n \gets n+1$; \;
        \If{$a_s = 0$}{
            $accepted \gets \textbf{false}$ and \textbf{break} \tcp*[r]{later draft contexts are invalid}
        }
    }
    \If(\tcp*[f]{all draft tokens accepted}){$accepted$ \textbf{and} $n < N$}{
    $\Pi^\star \gets G(\mathsf{k}, w_{1:n})$ \tcp*[r]{keyed Poisson source}
    $w_{n+1} \gets \textsc{PFR}(P(\cdot | w_{1:n}), \Pi^\star)$; 
    $n \gets n+1$ \tcp*[r]{bonus}
    }
}
\caption{Watermarkable Speculative Sampling via Poisson Processes (Simplified)}
\label{alg:single_PFR_Spec_Water_simplified} 
\vspace{-2pt}
\end{algorithm}


In LLM decoding, the target distributions are usually discrete; we prove in Appendix~\ref{subapp:PFR_discrete} that PFR reduces to using exponential random variables as in~\citep{kobus2025speculative}.
However, using Poisson processes becomes useful, and even necessary, when we extend to the multi-draft case in Section~\ref{sec::multi_draft}. In this case, %where 
we generate multiple samples, all following the desired distribution by the mapping theorem~\citep{last2017lectures} of Poisson processes, while achieving a higher token-level acceptance rate. 


{\color{teal}
We consider the watermark scores used in~\citep{aaronson2023watermarking, hu2024inevitable}, since our goal is to show that our scheme itself circumvents the tradeoff in~\citep{hu2024inevitable}, rather than avoiding it by choosing a different watermark metric, as in~\citep{he2026improving}. 
Moreover, as proved in Appendix~\ref{subsapp::PFR_gumbel}, our scheme is conceptually equivalent to the Gumbel watermark~\citep{aaronson2023watermarking} in the single-draft case. 
Therefore, using these scores not only leads to strong empirical performance, but also makes our scheme directly comparable to existing watermarking schemes.
}


If the output of key $\mathsf{k}$ is $w_{1:N}$ and $n_0$ is the initial prompt length, the detector evaluates the scored positions $\mathcal{T} \hspace{-1pt} := \hspace{-1pt}  \{n_0+1,\hspace{-1pt} \ldots\hspace{-1pt} ,N\}$. 
For $t\in\mathcal{T}$, let $c_t := w_{:t-1}$ be the prefix before position $t$, and the detector regenerates $\Pi_t\hspace{-1pt}  := \hspace{-1pt} G(\mathsf{k}, c_t)$. 
For each $v\hspace{-1pt} \in\hspace{-1pt} \mathcal{V}$, let $r_t(v)\hspace{-1pt} :=\hspace{-1pt}  \exp(-\mu(v)\tau_t(v))$ be the induced  uniform value where $\tau_t(v)\hspace{-1pt} \sim \hspace{-1pt} \mathrm{Exp}(\mu(v))$ is the first arrival time of $v$ in $\Pi_t$, and the detector computes
\vspace{-5pt}
\begin{equation}
    S_{\mathrm{A}}
    :=
    \sum\nolimits_{t\in\mathcal{T}} -\log(1-r_t(w_t)),\label{eq::S_A}
\end{equation}
which has expectation $\mathbf{E}[S_{\mathrm{A}}] \geq |\mathcal{T}| + (\pi^2/6-1)\sum_{t\in \mathcal{T}} \mathbf{E}[H(P(\cdot|w_{:t-1}))]$~\citep{aaronson2023watermarking}. 

In contrast, %otherwise, 
if the sequence is not generated with $\mathsf{k}$, the observed $w_t$ is independent of $\Pi_t$, %and 
$r_t(w_t)\sim \mathrm{Unif}[0,1]$ for $t\in\mathcal{T}$, and $\mathbf{E}[S_{\mathrm{A}}]  = |\mathcal{T}|$. 
Under the null, $S_{\mathrm{A}} \sim \mathrm{Gamma}(|\mathcal{T}|,1)$, and the $p$-value is \vspace{-2pt}
\[
p_{\mathrm{A}}
=
\mathbf{P}\big(\mathrm{Gamma}(|\mathcal{T}|, 1)\ge S_{\mathrm{A}}\big)
= \Gamma(M,S_{\mathrm{A}}) \,/\, \Gamma(|\mathcal{T}|), 
\]
where $\Gamma(|\mathcal{T}|)$ is the gamma function and $\Gamma(|\mathcal{T}|, S_{\mathrm{A}})$ is the upper incomplete gamma function, and the Average Negative Log P-value Per Token (ANLPPT) is $\mathrm{ANLPPT}_{\mathrm{A}} := -(1/|\mathcal{T}|) \log p_{\mathrm{A}}$. 
\vspace{-5pt}

\paragraph{At-a-Glance Experimental Validation.} Though most experiment results are in Section~\ref{sec::experiments}, we use Figure~\ref{fig::single_draft_qwen_cnn_new} to %\citep{hu2024inevitable}, and 
show that we avoid the tradeoff using the same score~\citep{hu2024inevitable}, $\sum_{t\in\mathcal{T}} r_t(w_t)$, that is easily connected to \eqref{eq::S_A}. 
The reason behind our improvement is, our keyed Poisson source simultaneously realizes the target token and couples it to the draft, rather than using a sequential route. 


% \vspace{-10pt}
\begin{figure}[htpb]
    \centering
    \includegraphics[scale=0.35]{images/qwen_cnn_figure1.png}
    \caption{
    The tradeoff between watermark strength and sampling efficiency applies to the algorithms in~\citep{hu2024inevitable}: Vanilla Speculative Sampling (VSpS), Maintain Watermark Strength (MWS), and Maintain Sampling Efficiency (MSE), but \textbf{not} to our algorithm based on the PFR, on which for both single-draft (lower-left) and multi-draft (upper) embedding watermark does not degrade the Average Accepted Tokens Per Step (AATPS). 
    The figure is from using Qwen2.5-7B-Instruct as the target model and Qwen2.5-0.5B-Instruct as the draft model on dataset \textsc{cnn\_dailymail} dataset; see Appendix~\ref{sec:single_draft} for results on other models and datasets
    }
    \label{fig::single_draft_qwen_cnn_new}
\end{figure}








\subsection{Drafter Invariance}
\label{sec::drafter_invar}


Conventional speculative sampling suffers from a major disadvantage~\citep{daliri2025coupling}: when the small drafter model $\mathcal{M}_d$ changes, e.g., because of a model update, the tokens generated by the large model may also change due to the coupling (Algorithm~\ref{alg::max_coupling}). 
In autoregressive language models, this can be even more problematic, since users usually do not want the model output to change under a fixed random seed: \emph{stable} models allow the output to be reliably reproduced for use, testing, and debugging. 


Based on the coupling without communication~\citep{bavarian2016optimality, li2018strong}, \emph{drafter invariant} sampling has recently been studied~\citep{daliri2025coupling}. It %, which 
refers to the case where, conditioned on the shared randomness $\mathcal R$ and the context $x_{:t}$, and potentially other things, the output-prefix law does not depend on the draft model. 
Let $Y_{1:\tau}, \widetilde{Y}_{1:\tau}$ denote the output sequences from different drafter models in block $\tau$, we define the following notion. 


\begin{definition}
\label{def:st_tm_drafter_invariance_single}
The algorithm $\mathcal{A}$ is said to be \emph{stopping-time drafter invariant} if, for every stopping time value $\tau_0$ in the support of $\tau$, every $1\le j\le \tau_0$, and every output prefix $y_{1:j}$, we have
\begin{equation*}
\mathbf{P}\big\{
Y_{1:j}=y_{1:j}\,\big|\,\mathcal{R},\,x_{:t},\,\tau=\tau_0
\big\}
=
\mathbf{P}\big\{
\widetilde{Y}_{1:j}=y_{1:j}\,\big|\,\mathcal{R},\,x_{:t},\,\widetilde{\tau}=\tau_0
\big\}.
\end{equation*}
\end{definition}

\begin{proposition} 
\label{prop::single_dft_st_tm_inv} 
Our Algorithm~\ref{alg:single_PFR_Spec_Water_simplified} is stopping-time drafter invariant. 
\end{proposition}

The proof can be found in Appendix~\ref{subsubapp::proof_single_dft_st_tm_inv}. 
The stopping-time drafter invariance is not satisfied by most existing schemes, e.g., SpecTr~\citep{sun2023spectr}, SpecInfer~\citep{miao2024specinfer}, and also \citep[Algorithm 2]{rowan2025list}. 




{\color{teal}
Intuitively, a stronger notion of drafter invariance helps watermarking but degrades sampling efficiency. 
The original definition, referred to as the ``strong drafter invariance''~\citep{daliri2025coupling}, was shown to have degraded sampling efficiency~\citep{rowan2025list}, in where the authors hence proposed a ``conditional drafter invariance'' that also conditions on the realized draft sequence. 
We hence propose the stopping-time drafter invariance that is weaker than the former to admit a better sampling efficiency, but stronger than the latter and more suitable for watermarking. 
The conditioning on $\tau$ is intentional: it separates the role of the drafter in determining the \emph{length} of the current speculative block from its role in determining the \emph{content} of the emitted tokens. 
Thus, the drafter is allowed to affect the stopping time, and hence the sampling efficiency, but once this stopping time is fixed, the law of every emitted prefix is governed by the target-side randomness and does not depend on the full realized draft sequence. 
}


In watermarking, such invariance is important. 
A watermark is usually tied to pseudorandomness and is detected from the final output sequence, while speculative sampling introduces an additional ambiguity: an emitted token may come from the draft, the target correction, or the bonus, which affects the watermark statistic, and is one reason behind the no-go tradeoff in~\citep{hu2024inevitable} and the use of acceptance-side information in~\citep{he2026improving}. 
However, a watermark detector typically observes only the final tokens and the secret key, not the realized draft sequence or acceptance-side information. 
Stopping-time drafter invariance is hence better aligned with watermarking: changing the drafter may affect the stopping time and the efficiency, but conditional on the block length, it does not introduce draft-specific dependence into the target-keyed token prefix. 
This makes the watermark directly tied to the target-side keyed randomness and leads to a more reproducible provenance signal. 


{\color{teal}
Moreover, \citep{he2026improving} proposed a new watermark metric that improves the tradeoff between watermark strength and sampling efficiency~\citep{hu2024inevitable}; see, e.g., \citep[Figure~1]{he2026improving}. 
However, their analysis relies on the maximal coupling probability $\mathbf{P}(X=Y)=1-\Vert P-Q\Vert_{\mathrm{TV}}$ for coupling $X\sim P, Y\sim Q$~\citep{thorisson2000coupling, den2012probability}. 
In contrast, the best matching probability under no-communication coupling is fundamentally different:
\begin{equation}
    \mathbf{P}(X = Y) \leq (1 - \Vert P - Q\Vert_{\mathrm{TV}}) / (1 + \Vert P - Q\Vert_{\mathrm{TV}}).
    \label{eq::cpl_wo_commu_bd}
\end{equation}
Thus, enforcing drafter invariance may appear to fundamentally reduce sampling efficiency. 
We argue, however, that this loss is not inevitable: by using Poisson processes, we derive a natural multi-draft extension that greatly improves the token acceptance rate, as elaborated in the next section.
}



\section{Multi-Draft Speculative Sampling}
\label{sec::multi_draft}


We now extend this to a watermarkable multi-draft algorithm by using a multi-sample generalization of PFR built on the mapping theorem for Poisson processes~\citep{li2021unified, last2017lectures}, which is likewise drafter invariant. 


We use $\mathcal{M}_t$ and $\mathcal{M}_d^{(b)}$ to denote the \emph{target} and \emph{draft} language models with %, respectively, where 
$1\leq b\leq B$. They induce conditional distributions $\mathcal{M}_t(\cdot| x_{:t})$ and $\mathcal{M}_d^{(b)}(\cdot| x_{:t})$, which give the probability of the next token given the current context $x_{:t}$, usually denoted by $c$.
In multi-draft speculative sampling~\citep{miao2024specinfer, sun2023spectr, khisti2024multi, rowan2025list}, at each context $c$, we generate a list of $B$ proposal drafts, often i.i.d. in practice~\citep{rowan2025list, sun2023spectr}  from $\mathcal{M}_d^{(1)}, \ldots, \mathcal{M}_d^{(B)}$. These proposal drafts %, which 
induce a depth-$L$ speculative tree rooted at the current prefix, which is then verified %and verify them 
in parallel with $\mathcal{M}_t$.
At each step, %Among the $B$ candidate tokens, 
if at least one of the $B$ candidate tokens is accepted, %at each step, 
the first accepted token will be appended to the final output; otherwise, the verification procedure stops. 
The output is $Y_{1:\tau}$, where $\tau$ is the number of accepted tokens plus one.

\subsection{Algorithm Design}

The PFR (Definition~\ref{def::PFR}) selects a sample $\tilde{Z}_P$ from $(Z_i)_i$ that has the smallest \emph{score} $T_i\cdot \big(\frac{\mathrm{d}P}{\mathrm{d}Q}(Z_i)\big)^{-1}$, where $(T_i)_i\sim \mathrm{PP}(1)$, and it follows the target distribution $P$. 
To get $B$ samples, multi-sample PFR (MPFR) selects samples with the smallest, second smallest, and so on up to the $B$-th smallest scores:
\begin{equation}
    \tilde{Z}_P(1), \tilde{Z}_P(2),\ldots, \tilde{Z}_P(B) \stackrel{\mathrm{iid}}{\sim} P, \label{eq::MPFR}
\end{equation}
where the exactness follows~\citep{li2021unified}; see Appendix~\ref{subapp::mapped_PP} for the formal definition and implementation. 


For multi-draft speculative sampling, we use $B$ i.i.d. drafts from a model $\mathcal{M}_d$, and at each sampling step we consider $P := \mathcal{M}_t(\cdot| x_{:t})$ as the target distribution and $Q := \mathcal{M}_d(\cdot| x_{:t})$ as the proposal distribution. 
Note that we consider \emph{identical} proposals~\citep{sun2023spectr} in this section, but our framework can also be generalized to non-identically distributed proposals~\citep{rowan2025list}, which is deferred to Appendix~\ref{app::non_id_prop}. 


Our scheme is %briefly 
sketched in Algorithm~\ref{alg:multi_PFR_spec_water_simplified}; see Algorithm~\ref{alg:multi_fwss_pfr_formal} for a detailed implementation. 
The core idea is as follows. 
To implement a multi-draft algorithm, Algorithm~\ref{alg:multi_PFR_spec_water_simplified} uses \emph{context-indexed} mapped Poisson processes.
Starting from the root prefix $h=w_{1:n}$, we build a depth-$\ell$ speculative tree, where each occupied context $c$ is assigned a keyed source $\Pi(c)$ and generates $\nu(c)$ drafts from $Q(\cdot |  c)$ via MPFR.
We then compute the target-side sample $\textsc{MPFR}(P,\Pi(c),1)$ for every $c$ and emit tokens by following the unique realized target path, continuing as long as the current target winner belongs to the draft set. 


Importantly, the Poisson process is indexed by contexts rather than by drafter identities, and thus MPFR works in a way that makes the emitted tokens tightly tied to the prefix and the secret key, which helps watermark detection.
In comparison, \citep{rowan2025list} uses draft-indexed algorithms, keeps a set of currently viable draft indices, and chooses the next emitted token using only the active drafts.
This is a natural implementation, but the active set makes the next emitted token depend on the realized draft sequences, causes the output to retain hidden dependence on draft-side trajectories, and hence only satisfies their conditional drafter invariance, but not stopping-time drafter invariance (Definition~\ref{def:st_tm_drafter_invariance_single}).


\begin{algorithm}[tbp]
% \begin{algorithm}[H]
\SetAlgoLined
\DontPrintSemicolon
\SetAlgoNoEnd
\SetKwIF{If}{ElseIf}{Else}{if}{:}{else if}{else}{end}
\SetKwInOut{Input}{Given}

\Input{lookahead $L$, number of homogeneous drafts $B$, output length $N$, target model $P$,
draft model $Q$, initial prompt $w_{1:n}$, keyed source generator $G$,
and secret key $\mathsf{k}$}

\While{$n < N$}{
    $h \gets w_{1:n}$; $\ell \gets \min\{L,\,N-n\}$\tcp*[r]{root context of current block} 
    initialize root context $h$ with multiplicity $B$\;
    
    \For(\tcp*[f]{build context tree up to depth $\ell$}){$s=1$ \KwTo $\ell$}{
        \For{each currently occupied context $c$}{
            $\Pi(c) \gets G(\mathsf{k},c)$ \tcp*[r]{\textcolor{blue}{keyed Poisson source for watermark}}
            
            generate $\nu(c)$ drafts by $Q(\cdot | c)$ using $\textsc{MPFR} (Q(\cdot | c),\Pi(c), \cdot)$ \tcp*[r]{Eq.\eqref{eq::MPFR}} 
            % \[
            % \textsc{MPFR}(Q(\cdot | c),\Pi(c),1),\ldots,
            % \textsc{MPFR}(Q(\cdot | c),\Pi(c),\nu(c))
            % \]
            merge equal draft tokens into child contexts and update their multiplicities\;
        }
    }

    compute target-side sample $y(c)=\textsc{MPFR}(P(\cdot | c),\Pi(c),1)$ for every $c$ \tcp*[r]{in parallel}

    $c^\star \gets h$; $accepted \gets \textbf{true}$\;
    \For{$s=1$ \KwTo $\ell$}{
        emit the target token $y(c^\star)$ and append it to the output\;
        \If(\tcp*[f]{$\mathcal{D}(c^\star)$ contains distinct draft tokens at $c$}){$y(c^\star)\notin \mathcal{D}(c^\star)$}{
            $accepted \gets \textbf{false}$ and \textbf{break}
            \tcp*[r]{first rejection ends the block}
        }
        \If{$n = N$}{
            \textbf{break}
        }
        $c^\star \gets c^\star \,\|\, y(c^\star)$
        \tcp*[r]{move to the next realized context}
    }

    \If{$accepted$ \textbf{and} $n < N$}{
        emit one bonus token by $\textsc{MPFR}(P(\cdot | c^\star),G(\mathsf{k},c^\star),1)$ 
    }
}
\caption{Watermarkable Multi-Draft Speculative Sampling (Simplified)}
\label{alg:multi_PFR_spec_water_simplified}
\end{algorithm}



\subsection{Theoretic Guarantee}
\label{subsec::MPFR_theory}
\begin{theorem}
    \label{thm::PML_accept}
    With $B$ draft models, the acceptance probability of at least one token at a step satisfies 
    \begin{equation}
        \label{eq::PML_accept}
        \mathbf{P}(\mathrm{accept}) \geq 
        1-\sum_{i\in\mathcal{V}} P(i)\Big(\frac{P(i)}{P(i) + Q(i)}\Big)^B. 
    \end{equation}
\end{theorem}
The proof is in Appendix~\ref{app::theory}, where we also give 
a generalized version that accounts for the number of active drafts, which %whose number 
can be smaller than $B$, the number of drafts initialized at the start of the block. 
We demonstrate the advantage of our acceptance probability over~\citep{rowan2025list} in Appendix~\ref{app::theory}. 





\section{Experiments}
\label{sec::experiments}

We evaluate whether the proposed Poisson-process sampler improves the empirical
frontier between speculative decoding efficiency and watermark detectability, while
preserving generation quality. Unless otherwise stated, we use Qwen2.5-7B-Instruct as
the target model and Qwen2.5-0.5B-Instruct as the draft model on
\textsc{CNN/DailyMail}. Additional model-dataset pairs are reported in the appendix.

\paragraph{Baselines and evaluation protocol.}
We compare against four groups of baselines. First, \textsc{basic\_uwm} denotes
autoregressive Gumbel-style unbiased watermarking without speculation. It serves as
a watermark-only reference for the detection signal achievable when no speculative
verification is used. Second, we compare with the speculative watermarking baselines
of \citet{hu2024inevitable}: \textsc{VSpS}, \textsc{MWS}, and \textsc{MSE}.
Here \textsc{VSpS} applies vanilla speculative sampling together with the unbiased
watermarking rule; \textsc{MWS} is designed to maintain watermark strength, but
typically pays a speculative-acceptance cost; and \textsc{MSE} is designed to
maintain sampling efficiency, but typically weakens the watermark signal. These
methods represent different points on the previously reported
efficiency--watermark trade-off. Third, for fixed-FPR detection comparisons, we also
include \textsc{MSE\_PSEUDO}, the pseudo-randomized variant of
\citet{he2026improving}, which improves the trade-off by injecting additional
pseudorandomness into the speculative acceptance or residual-sampling mechanism.
Fourth, in the multi-draft setting, we compare \textsc{MPFR} against the
unwatermarked List sampling~\citep{rowan2025list} decoder of \citet{rowan2025list}, which is the
closest list-coupling baseline for speculative efficiency. We also include
no-watermark variants of our methods to separate the cost of watermarking from the
cost of Poisson coupling itself.

We report three classes of metrics. For speculative efficiency, we use average
accepted tokens per speculative step (AATPS), and report token rate in the appendix
as a system-level diagnostic. For watermark detectability, we report ANLPPT-U and
TPR@1\%FPR. The TPR threshold is calibrated on non-watermarked generations to obtain
a 1\% false-positive rate and then evaluated on watermarked generations. For
generation quality, we report target-model log perplexity (LPPL), together with
task metrics such as ROUGE-L when references are available. This separation is
important: AATPS measures the algorithmic acceptance behavior, ANLPPT-U and
TPR@1\%FPR measure detector-visible watermark evidence, and LPPL audits whether the
target-model marginal is preserved.

\paragraph{Single-draft efficiency-watermark frontier.}
Figure~\ref{fig::single_draft_qwen_cnn_new} shows the single-draft
efficiency--watermark frontier. The prior speculative watermarking baselines exhibit
the expected trade-off: methods that preserve stronger watermark evidence tend to
sacrifice speculative efficiency, while methods optimized for efficiency produce
weaker watermark evidence. In contrast, \textsc{PFR} lies in the upper-right region
of the plot, achieving both high AATPS and strong ANLPPT-U. Moreover,
\textsc{PFR} and its no-watermark variant have nearly identical AATPS, indicating
that the watermark signal is introduced through the target-side keyed race rather
than through an acceptance-penalizing reweighting step. Additional single-draft
results on other model--dataset pairs are provided in Appendix~\ref{sec:single_draft}.

\paragraph{Multi-draft scaling.}
We next evaluate whether the same construction scales to multiple drafts. We compare
\textsc{MPFR} with the unwatermarked List sampling~\citep{rowan2025list} decoder of
\citet{rowan2025list} using lookahead $L=4$ and draft counts $B\in\{2,4,6,8\}$.
 Full results for
other model--dataset pairs, the additional speculative-watermarking baselines
\textsc{MWS}/\textsc{MSE} of \citet{hu2024inevitable}, the remaining ANLPPT
variants (Li, PL), and LPPL are deferred to Appendix~\ref{sec:multi_draft}.



\paragraph{Detection at fixed false-positive rate.}
ANLPPT-U measures average evidence per token, while TPR@1\%FPR measures operational
detectability under a calibrated false-positive budget. Figure~\ref{fig:drafter_invariance}
(left) compares detection power across methods, including the pseudo-randomized
\textsc{MSE\_PSEUDO} variant of \citet{he2026improving}. \textsc{PFR} and
\textsc{MPFR} closely track the autoregressive watermark reference and the
strong-watermark baseline \textsc{MWS}, while efficiency-oriented baselines such as
\textsc{MSE} produce weaker detection. \textsc{MSE\_PSEUDO} improves over
\textsc{MSE}, but still remains below the Poisson-race methods in our evaluated
settings. This supports the claim that the Poisson construction preserves
detector-visible watermark evidence while still benefiting from speculative
verification.

\begin{figure}[t]
\centering
\includegraphics[width=\linewidth]{images/drafter_invariance_composite.png}
\caption{
\textbf{Detection and robustness to drafter substitution.}
\textbf{(a)} Detection at a fixed false-positive rate on the default drafter.
\textsc{PFR} and \textsc{MPFR} achieve TPR@1\%FPR close to the strong-watermark
references \textsc{Basic-UWM} and \textsc{MWS}, and outperform the
efficiency-oriented baselines \textsc{MSE} and \textsc{MSE-Pseudo}.
\textbf{(b)} Robustness under drafter substitution. We vary the drafter across
four conditions while keeping the target model, watermark key, and prompts fixed.
\textsc{PFR} remains nearly unchanged, whereas \textsc{MSE} and
\textsc{MSE-Pseudo} exhibit substantially larger drift.
}
\label{fig:drafter_invariance}
\end{figure}

\paragraph{Drafter-substitution ablation.}
Finally, we test whether the watermark signal remains stable when the proposal
mechanism changes. We keep the target model, watermark key, and evaluation prompts
fixed, and vary only the drafter by either changing its model scale or its sampling
temperature. Figure~\ref{fig:drafter_invariance} (right) shows that \textsc{PFR}
has at most one percentage point of TPR drift at
$T_{\mathrm{eval}}\in\{64,128\}$ across the four drafter conditions. In contrast,
\textsc{MSE} drifts by $12$--$22$ percentage points, and \textsc{MSE\_PSEUDO}
drifts by $6$--$11$ percentage points across drafter conditions. This behavior is
consistent with the intended role of target-side keyed Poisson races: the drafter can
affect how many tokens are verified in a speculative block, but the watermark
statistic of the emitted tokens is tied to the realized prefix and the target-side
keyed randomness. We further verify that generation quality is stable across
decoders; full LPPL and ROUGE tables are reported in Appendix~\ref{app:ablation}.




\section{Limitations}


We have proposed a draft-invariant algorithm that sidesteps the no-go theorem of~\citep{hu2024inevitable}, but a fundamental characterization of the tradeoff among target-model communication, watermark strength, and speculative sampling efficiency remains unclear. 
Moreover, compared with~\citep{rowan2025list}, although our method achieves several advantages, such as slightly higher AATPS and watermarking capability, its token rate is slightly lower, mainly because MPFR performs more coupling work per context and requires additional watermark bookkeeping. 
A more practical implementation of our multi-draft algorithm with a higher token rate remains to be studied.

\begin{table}[h]
\centering
\caption{Token rate (tok/s) comparison across three measurement runs at matched drafter ratio 1:103, CNN/DailyMail, $L=K=4$.
Column~1 reproduces Hu \& Huang~\cite{hu2024inevitable} Table~1 (LLaMa-7B + LLaMa-68m, $K=4$, DeltaGumbel; original units PTT in ms/tok converted to tok/s).
Columns~2 and~3 report two independent measurements on Vicuna-7B + vicuna-68m, $L=4$, $n=1000$ prompts on 4$\times$A6000.
Column~2 is a single-pass run; Column~3 is a separate run with per-prompt timing aggregated as mean $\pm$ std.
.}
\label{tab:token_rate_three_way}
\small
\setlength{\tabcolsep}{6pt}
\begin{tabular}{l c c c}
\toprule
\multirow{2}{*}{Method}
 & \cite{hu2024inevitable}
 & Ours (single-pass)
 & Ours (mean $\pm$ std) \\
 & LLaMa-7B/68m & Vicuna-7B/68m & Vicuna-7B/68m \\
\midrule
Basic / \texttt{basic\_uwm}             & $33.81$ & $35.29$ & $37.06 \pm 0.85$ \\
\texttt{VSpS} / \texttt{mc} (no wm)     & $51.89$ & $50.57$ & $53.94 \pm 8.24$ \\
\texttt{pfr\_no\_watermark}              & ---     & $47.89$ & $40.08 \pm 7.88$ \\
\textbf{\texttt{pfr} (ours)}             & ---     & $\mathbf{52.86}$ & $\mathbf{41.80 \pm 9.36}$ \\
\texttt{MSE} / \texttt{mc\_uwm\_speed}   & $48.10$ & $50.23$ & $32.73 \pm 6.78$ \\
\texttt{MWS} / \texttt{mc\_uwm\_strength} & $45.05$ & $48.82$ & $27.25 \pm 4.41$ \\
\texttt{mse\_pseudo}~\cite{he2026improving} & ---  & $48.31$ & $33.01 \pm 6.81$ \\
\bottomrule
\end{tabular}
\end{table}


\newpage













\newpage
\bibliography{ref.bib}
\bibliographystyle{plainnat}

\appendix



\section{More on PFR and Related Works}
\label{app::Poisson}

\subsection{PFR for Discrete Distributions}
\label{subapp:PFR_discrete}

It is known that the general PFR~\citep{li2018strong, li2021unified} reduces to the use of exponential random variables~\citep{li2018strong, liu2025one} in the case of simulating discrete distributions, which is similar to the Gumbel-max trick~\citep{gumbel1954statistical} used for speculative sampling~\citep{kobus2025speculative}.

In~\citep{li2024channel}, it is discussed how to extend the discrete case to the general case. However, the other direction, namely reducing the general case to the discrete case, has not been discussed in the literature, to the best of the authors' knowledge.
To make the paper self-contained, we provide formal statements and proofs for this direction.


The first define the exponential functional representation (EFR) as follows (also see \citep{li2018strong, liu2025one}). 
Let $\mathcal{Z}$ be a finite or countably infinite set, and let $P$ be a distribution on $\mathcal{Z}$. 
Let $(E_z)_{z\in\mathcal{Z}}$ be independent $\mathrm{Exp}(1)$ random variables. 
With the convention that $E_z/P(z)=+\infty$ whenever $P(z)=0$, define
\begin{equation}
    Z_{\mathrm{EFR}}
    :=
    \arg\min_{z\in\mathcal{Z}} \frac{E_z}{P(z)} .
    \label{eq:EFR}
\end{equation}
As proved in Appendix~\ref{subsapp::PFR_gumbel}, EFR is equivalent to the Gumbel-max trick.
In the following, we prove that PFR reduces to EFR in the discrete case, i.e., when the target and proposal distributions are discrete. 



\begin{theorem}
\label{thm::PFR_reduces_to_EFR}
Let $\mathcal{Z}$ be a finite or countably infinite set. Let $Q$ and $P$ be distributions on $\mathcal{Z}$ such that $P\ll Q$. Write $q_z:=Q(z), p_z:=P(z), z\in\mathcal{Z}$. 
The PFR considers a Poisson process $(Z_i,T_i)_{i\geq 1}$ on $\mathcal{Z}\times[0,\infty)$ with intensity measure
$Q\times\lambda_{[0,\infty)}$, and $\tilde T_i = T_i\left(\frac{\mathrm{d}P}{\mathrm{d}Q}(Z_i)\right)^{-1}$. 
For each $z$ with $q_z>0$, define the first arrival time of symbol $z$ by
\begin{equation}
    S_z:=\inf\{T_i:Z_i=z\},
    \label{eq:first_arrival_symbol_z}
\end{equation}
and define
\begin{equation}
    E_z:=q_z S_z .
    \label{eq:Ez_from_PFR}
\end{equation}
For $z$ with $q_z=0$, define $E_z$ to be an arbitrary independent $\mathrm{Exp}(1)$ random variable, independent of everything else. Then the random variables $(E_z)_{z\in\mathcal{Z}}$ are independent $\mathrm{Exp}(1)$ random variables, and the PFR output satisfies
\begin{equation}
    Z_{\mathrm{PFR}}
    =
    \arg\min_i \tilde T_i
    =
    \arg\min_{z\in\mathcal{Z}}\frac{E_z}{p_z}
    =
    Z_{\mathrm{EFR}}
    \qquad \text{a.s.},
    \label{eq:PFR_equals_EFR}
\end{equation}
where we use the convention that $E_z/p_z=+\infty$ if $p_z=0$.
\end{theorem}

\begin{proof}
For $z\in\mathcal{Z}$ and $t\geq 0$, define
\[
    N_z(t):=\#\{i:T_i\leq t,\ Z_i=z\}.
\]
Since $(Z_i,T_i)_{i\geq 1}$ is a Poisson process with intensity measure
$Q\times\lambda_{[0,\infty)}$, for each $z$ with $q_z>0$, the process
$(N_z(t))_{t\geq 0}$ is a Poisson process with rate $q_z$. Moreover, for distinct symbols
$z_1,\ldots,z_m$, the processes
\[
    (N_{z_1}(t))_{t\geq 0},\ldots,(N_{z_m}(t))_{t\geq 0}
\]
are independent, because they count disjoint subsets of the underlying Poisson process.

Therefore, for $z$ with $q_z>0$, the first arrival time
\[
    S_z=\inf\{t:N_z(t)\geq 1\}
\]
satisfies
\[
    \mathbf{P}(S_z>s)
    =
    \mathbf{P}(N_z(s)=0)
    =
    e^{-q_z s},
    \qquad s\geq 0.
\]
Thus $S_z\sim\mathrm{Exp}(q_z)$. Hence
\[
    E_z=q_z S_z\sim \mathrm{Exp}(1).
\]
The independence of the Poisson counting processes for different symbols implies that the random variables
\[
    (E_z)_{z:q_z>0}
\]
are independent. For symbols with $q_z=0$, we have $p_z=0$ because $P\ll Q$, and these symbols can never be selected by PFR. Defining independent $\mathrm{Exp}(1)$ variables $E_z$ for those symbols therefore does not change the output. Hence the full family $(E_z)_{z\in\mathcal{Z}}$ is independent and identically distributed as $\mathrm{Exp}(1)$.

It remains to show that the PFR selection rule reduces to the EFR selection rule. Since $P\ll Q$, for every $z$ with $p_z>0$, we also have $q_z>0$, and
\[
    \frac{\mathrm{d}P}{\mathrm{d}Q}(z)
    =
    \frac{p_z}{q_z}.
\]
Therefore, for every index $i$ with $Z_i=z$ and $p_z>0$,
\[
    \tilde T_i
    =
    T_i\left(\frac{p_z}{q_z}\right)^{-1}
    =
    T_i\frac{q_z}{p_z}.
\]
Taking the minimum over all Poisson points with mark $z$, we get
\[
    \inf_{i:Z_i=z}\tilde T_i
    =
    \left(\inf_{i:Z_i=z} T_i\right)\frac{q_z}{p_z}
    =
    S_z\frac{q_z}{p_z}
    =
    \frac{E_z}{p_z}.
\]
If $p_z=0$, then $\mathrm{d}P/\mathrm{d}Q(z)=0$ whenever $q_z>0$, and hence
\[
    \tilde T_i
    =
    T_i\left(\frac{\mathrm{d}P}{\mathrm{d}Q}(z)\right)^{-1}
    =
    +\infty
\]
for every Poisson point with mark $z$. This agrees with the EFR convention
\[
    \frac{E_z}{p_z}=+\infty .
\]
If $q_z=0$, then also $p_z=0$, and no Poisson point has mark $z$ almost surely, so such a symbol is irrelevant in both constructions.

Thus,
\[
    \inf_i \tilde T_i
    =
    \inf_{z\in\mathcal{Z}}\inf_{i:Z_i=z}\tilde T_i
    =
    \inf_{z\in\mathcal{Z}}\frac{E_z}{p_z}.
\]
Moreover, the minimizer is unique almost surely. Indeed, for every $z$ with $p_z>0$,
\[
    \frac{E_z}{p_z}\sim \mathrm{Exp}(p_z),
\]
and these variables are independent. Since exponential random variables are continuous, ties have probability zero. Therefore, almost surely, 
\[
    Z_{\mathrm{PFR}}
    =
    Z_K
    =
    \arg\min_{z\in\mathcal{Z}}\frac{E_z}{p_z}
    =
    Z_{\mathrm{EFR}}.  
\]
\end{proof}


\subsection{Relation between PFR and Gumbel-Max Trick}
\label{subsapp::PFR_gumbel}


As discussed in~\citep{zhao2024permute}, the Gumbel Watermark~\citep{aaronson2023watermarking} is based on  
\begin{equation*}
    y_t = \mathrm{argmin}_{y\in\mathcal{V}} \frac{u_t(y)}{T} + G_t(y)
\end{equation*}
where $G(y)\sim \mathrm{Gumbel}(0,1)$ i.i.d. for each $t$ and $y$. 
Here, we show that it is equivalent to the PFR in the discrete case. 
Recall $\mathrm{Gumbel}(0,1) \sim -\log\big(\log(1/\mathrm{Unif}[0,1]) \big)$. 
We calculate
\begin{align*}
y_t
&= \arg\max_{y\in\mathcal{V}}\left(\frac{u_t(y)}{T} + G_t(y)\right)\\
&= \arg\max_{y\in\mathcal{V}}\left(\frac{u_t(y)}{T} - \log Z_t(y)\right)
\qquad\Bigl(Z_t(y) \stackrel{\mathrm{iid}}{\sim} \mathrm{Exp}(1),\; G_t(y)\stackrel{d}{=}-\log Z_t(y)\Bigr)\\
&= \arg\max_{y\in\mathcal{V}}\left(-\log \bigl(Z_t(y)e^{-u_t(y)/T}\bigr)\right)\\
&= \arg\min_{y\in\mathcal{V}} Z_t(y)e^{-u_t(y)/T}\\
&= \arg\min_{y\in\mathcal{V}} \frac{Z_t(y)}{\exp(u_t(y)/T)}\\
&= \arg\min_{y\in\mathcal{V}} \frac{Z_t(y)}{w_t(y)}
\qquad\Bigl(w_t(y):=\exp(u_t(y)/T)\Bigr)\\
&= \arg\min_{y\in\mathcal{V}} \frac{Z_t(y)}{P_t(y)},
\end{align*} 
where the last equality is because $P_t(y):=\tfrac{w_t(y)}{\sum_{y'\in\mathcal{V}} w_t(y')}$ and scaling by a positive constant doesn't change the argmin. 



\subsection{Mapped Poisson Process}
\label{subapp::mapped_PP}

\subsubsection{Definitions}

\begin{definition}[Mapped Poisson Process~\citep{li2021unified}] \label{def::mapped_PP}
Let $(T_{i})_{i}$ be a Poisson process of rate $1$,
independent of $Z_{i}\stackrel{\mathrm{iid}}{\sim}Q$ for $i=1,2,\ldots$. 
Let $i_{P, 1},i_{P, 2}, \ldots\in\mathbb{N}$ be a sequence of distinct integers such that $\bigcup_{j=1}^\infty \big\{i_{P, j} \big\} = \big\{i:\frac{\mathrm{d}P}{\mathrm{d}Q}(Z_{i})>0 \big\}$ and $\big\{T_{i_{P, j}} \big(\frac{\mathrm{d}P}{\mathrm{d}Q}(Z_{i_{P, j}})\big)^{-1} \big\}_{j \in \mathbb{N}}$ is sorted in ascending order with arbitrary tie-breaking.\footnote{Note a tie occurs with probability $0$.}
For $j \in\mathbb{N}$ and suppose $P$ is the target distribution, we denote
\begin{equation}
    \left\{
    \tilde{Z}_P(j), \tilde{T}(j)
    \right\}_{j\in\mathbb{N}} 
    := \left\{ Z_{i_{P, j}}, T_{i_{P, j}}
    \big(\frac{\mathrm{d}P}{\mathrm{d}Q}(Z_{i_{P, j}})\big)^{-1} \right\}_{j\in\mathbb{N}}. 
\end{equation}
Give a positive integer $B$, the mapped Poisson Process selects the first $B$ samples in $\left\{ \tilde{Z}_P(j), \tilde{T}(j) \right\}_{j\in\mathbb{N}}$. 
\end{definition}


\begin{proposition}
\label{prop:ppr_output_dist}
By the mapping theorem~\citep{last2017lectures}, $\big\{ \tilde{Z}_P(j), \tilde{T}(j) \big\}_{j\in\mathbb{N}}$ is a Poisson process with intensity measure $P\times \lambda_{\mathbb{R}_{\geq 0}}$ where $\mathbb{R}_{\geq 0}$ to denote nonnegative real numbers, and therefore we have 
\begin{equation}
    \tilde{Z}_P (1), \tilde{Z}_P (2), \ldots \stackrel{\mathrm{iid}}{\sim} P.
\end{equation}
\end{proposition}







\subsubsection{Direct Implementations}
\label{subapp::direct_implementation_MPFR}
We can implement the MPFR just like~\citep{theis2022algorithms} for PFR, as follows: 

\begin{algorithm}[H]
\SetAlgoLined
\DontPrintSemicolon
\SetAlgoNoEnd
\SetKwIF{If}{ElseIf}{Else}{if}{:}{else if}{else}{end}
\SetKwInOut{Input}{Given}
\SetKwInOut{Output}{Return}

\Input{target distribution $P$, proposal distribution $Q$, likelihood ratio
$r=dP/dQ$, number of samples $B$, and a constant $\underline w>0$ such that
$1/r(u)\ge \underline w$ $Q$-a.s. on $\{r(u)>0\}$}
\Output{$B$ samples from $P$}

$H \gets \emptyset$ \tcp*[r]{$H$ stores the current $B$ smallest-score candidates}
$t \gets 0$; $i \gets 1$\;

\While{$|H|<B$ \textbf{or} $\max_{(s,j,z)\in H} s > t\underline w$}{
    Sample $E_i\sim \mathrm{Exp}(1)$ and set $t\gets t+E_i$ \tcp*[r]{Poisson process}
    
    Sample $Z_i\sim Q$\;
    
    \eIf{$r(Z_i)>0$}{
        $S_i \gets t/r(Z_i)$\;
    }{
        $S_i \gets +\infty$\;
    }

    \If{$|H|<B$}{
        Insert $(S_i,i,Z_i)$ into $H$\;
    }
    \ElseIf{$S_i < \max_{(s,j,z)\in H} s$}{
        Remove an element $(s^\star,j^\star,z^\star)\in H$ satisfying
        $s^\star=\max_{(s,j,z)\in H} s$\;
        
        Insert $(S_i,i,Z_i)$ into $H$\;
    }

    $i\gets i+1$\;
}
Sort the elements of $H$ as
\[
(S_{i_1},i_1,Z_{i_1}),\ldots,(S_{i_B},i_B,Z_{i_B})
\]
so that
\[
S_{i_1}\le S_{i_2}\le \cdots \le S_{i_B}.
\]

\Return{$(Z_{i_1},\ldots,Z_{i_B})$}

\caption{Finite implementation of the $B$-sample mapped PFR}
\label{alg:B_sample_finite_PFR}
\end{algorithm}

 
\subsubsection{Proof of Correctness}

We then prove that the implementation of Algorithm~\ref{alg:B_sample_finite_PFR} is correct. 


\begin{theorem}
Let $P\ll Q$ be probability measures on a Polish space $\mathcal U$, and let
$r:=dP/dQ$. Let $Z_1,Z_2,\ldots\stackrel{iid}{\sim} Q$ and let
$0<T_1<T_2<\cdots$ be the arrival times of a unit-rate Poisson process,
independent of $(Z_i)_{i\ge1}$. Define
\[
S_i :=
\begin{cases}
T_i/r(Z_i), & r(Z_i)>0,\\
+\infty, & r(Z_i)=0.
\end{cases}
\]
Assume that there exists $\underline w>0$ such that
\[
1/r(u)\ge \underline w
\qquad Q\text{-a.s. on }\{r(u)>0\}.
\] 
Fix $B\in\mathbb N$. Consider the algorithm that scans the candidates in
the order $i=1,2,\ldots$, keeps the $B$ smallest observed scores, and stops
at the first time $n$ such that the current $B$-th smallest observed score
is at most $T_n\underline w$. Then the algorithm terminates almost surely
after finitely many iterations. Moreover, at termination, the retained
candidates are exactly the $B$ smallest-score candidates among the infinite
sequence. Consequently, the returned marks are distributed as
\[
(\widetilde U_P(1),\ldots,\widetilde U_P(B)),
\]
the first $B$ points of the mapped Poisson process with respect to $P$,
and hence are i.i.d. according to $P$.
\end{theorem}



\begin{proof}
Let
\[
\Pi_Q:=\{(Z_i,T_i):i\ge 1\}.
\]
Since $0<T_1<T_2<\cdots$ are the arrival times of a unit-rate Poisson
process and $Z_1,Z_2,\ldots\stackrel{iid}{\sim}Q$ are independent marks,
$\Pi_Q$ is a Poisson point process on $\mathcal U\times \mathbb R_{\ge 0}$
with intensity measure $Q\times \lambda$, where $\lambda$ denotes the
Lebesgue measure on $\mathbb R_{\ge 0}$.

Consider the measurable mapping
\[
\Phi:\mathcal U\times \mathbb R_{\ge 0}\to
\mathcal U\times \mathbb R_{\ge 0},
\qquad
\Phi(u,t)=\left(u,\frac{t}{r(u)}\right)
\]
on the set $\{u:r(u)>0\}$. Points with $r(u)=0$ are assigned score
$+\infty$ and hence never appear among the finite mapped arrival times.
For any nonnegative measurable function $g$ on
$\mathcal U\times \mathbb R_{\ge 0}$, we have
\[
\begin{aligned}
\int_{\mathcal U}\int_0^\infty
g\left(u,\frac{t}{r(u)}\right)\,dt\,Q(du)
&=
\int_{\{r>0\}}\int_0^\infty
g\left(u,\frac{t}{r(u)}\right)\,dt\,Q(du) \\
&=
\int_{\{r>0\}}\int_0^\infty
g(u,s)r(u)\,ds\,Q(du) \\
&=
\int_{\mathcal U}\int_0^\infty
g(u,s)\,ds\,P(du).
\end{aligned}
\]
Therefore, by the mapping theorem for Poisson point processes, the mapped
point process
\[
\Pi_P:=\left\{\left(Z_i,S_i\right):i\ge 1,\ r(Z_i)>0\right\},
\qquad
S_i:=\frac{T_i}{r(Z_i)},
\]
is a Poisson point process on $\mathcal U\times \mathbb R_{\ge 0}$ with
intensity measure $P\times \lambda$.

Let
\[
S_{i_1}<S_{i_2}<\cdots
\]
be the increasing ordering of the finite scores in $\Pi_P$. Since the
time-marginal intensity of $\Pi_P$ is $\lambda$, the ordered times
$S_{i_1},S_{i_2},\ldots$ are the arrival times of a unit-rate Poisson
process. Moreover, conditional on these arrival times, the associated marks
are independent with common distribution $P$. Hence
\[
(Z_{i_1},Z_{i_2},\ldots)\stackrel{d}{=}(U_1,U_2,\ldots),
\]
where $U_1,U_2,\ldots\stackrel{iid}{\sim}P$. In particular,
\[
(Z_{i_1},\ldots,Z_{i_B})
\]
are exactly the first $B$ points of the mapped Poisson process with respect
to $P$, and they are i.i.d. according to $P$.

It remains to show that the finite algorithm stops safely and terminates
almost surely. After observing candidates $1,\ldots,n$, let
\[
S_{n,(B)}
\]
denote the $B$-th smallest value among $S_1,\ldots,S_n$, whenever at least
$B$ finite scores have been observed. By construction, the algorithm keeps
the $B$ smallest observed scores, so the largest score retained by the
algorithm is exactly $S_{n,(B)}$.

Suppose the algorithm stops at time $n$. Then
\[
S_{n,(B)}\le T_n\underline w.
\]
For every future candidate $m>n$, we have $T_m>T_n$. If $r(Z_m)=0$, then
$S_m=+\infty$, so such a candidate cannot improve the current list. If
$r(Z_m)>0$, then by the assumption $1/r(u)\ge \underline w$ $Q$-a.s.,
\[
S_m
=
\frac{T_m}{r(Z_m)}
=
T_m\frac{1}{r(Z_m)}
\ge
T_m\underline w
>
T_n\underline w
\ge
S_{n,(B)}.
\]
Thus every future score is strictly larger than the current $B$-th smallest
observed score. Therefore no future candidate can enter the set of the
global $B$ smallest scores. Hence, when the algorithm stops, the retained
candidates are exactly the global $B$ smallest-score candidates among the
infinite sequence.

Finally, we prove almost sure finite termination. Since $\Pi_P$ is a
Poisson point process with intensity $P\times \lambda$, its first $B$ mapped
arrival times are finite almost surely. Let
\[
S_{i_B}
\]
denote the $B$-th smallest mapped score, and let
\[
N_0:=\max\{i_1,\ldots,i_B\}.
\]
Then $N_0<\infty$ almost surely. After the algorithm has processed
candidates $1,\ldots,N_0$, it has already observed the true global first
$B$ mapped candidates, so the largest retained score is $S_{i_B}$. Since
\[
T_n\to \infty
\qquad\text{almost surely},
\]
and since $\underline w>0$, there exists an almost surely finite random
index $N\ge N_0$ such that
\[
T_N\underline w\ge S_{i_B}.
\]
At time $N$, the algorithm therefore satisfies the stopping condition
\[
S_{N,(B)}=S_{i_B}\le T_N\underline w.
\]
Thus the algorithm terminates almost surely after finitely many iterations.

Combining the safe stopping argument with the mapped Poisson process
argument, the algorithm returns exactly
\[
(Z_{i_1},\ldots,Z_{i_B}),
\]
which are the first $B$ points of the mapped Poisson process with respect to
$P$. Consequently, the returned samples are i.i.d. according to $P$.
\end{proof}





\subsection{A better finite-support MPFR implementation}
\label{subsec::direct_topk_MPFR}

The proposal-scanning implementation in Appendix~\ref{subapp::direct_implementation_MPFR} can be slow in finite-vocabulary LLM decoding. 
Its stopping rule depends on a lower bound on the likelihood ratio, which can be very small when a rare token has small proposal probability but non-negligible target probability. 
In our experiments, the logits are processed by temperature and top-$k$ truncation, so the resulting distribution has finite support of size at most $K$. 
In this setting, we can simulate the relevant Poisson clocks directly and avoid proposal scanning altogether.

Let $P$ be a distribution on a finite vocabulary $\mathcal U$, and define
\[
\mathcal A_P:=\{u\in\mathcal U:P(u)>0\}.
\]
For each $u\in\mathcal A_P$, let
\[
G_{u,j}:=\sum_{m=1}^{j}E_{u,m},\qquad j\geq 1,
\]
where the $E_{u,m}$'s are independent $\mathrm{Exp}(1)$ random variables. 
We associate to $(u,j)$ the mapped-Poisson score
\[
S_{u,j}:=\frac{G_{u,j}}{P(u)}.
\]
The first $B$ MPFR samples are the labels of the $B$ smallest scores among
\[
\{S_{u,j}:u\in\mathcal A_P,\ j\geq 1\}.
\]
Since no token can contribute more than $B$ arrivals among the first $B$ total arrivals, it is enough to generate $j=1,\ldots,B$ for every $u\in\mathcal A_P$.


\begin{algorithm}[tbp]
\SetAlgoLined
\DontPrintSemicolon
\SetAlgoNoEnd
\SetKwInOut{Input}{Given}
\SetKwInOut{Output}{Return}

\Input{finite distribution $P$ on $\mathcal U$, number of samples $B$, and keyed source $\Pi$}
\Output{$B$ samples from $P$}

$\mathcal A_P\gets\{u\in\mathcal U:P(u)>0\}$\;
$H\gets\emptyset$\tcp*[r]{candidate triples $(S,u,j)$}

\For{each $u\in\mathcal A_P$}{
    $G\gets 0$\;
    \For{$j=1$ \KwTo $B$}{
        Generate $E_{u,j}\sim\mathrm{Exp}(1)$ from $\Pi$\;
        $G\gets G+E_{u,j}$\;
        $S\gets G/P(u)$\;
        Insert $(S,u,j)$ into $H$\;
    }
}

Let
\[
(S_{(1)},U_{(1)},J_{(1)}),\ldots,(S_{(B)},U_{(B)},J_{(B)})
\]
be the $B$ elements of $H$ with smallest scores, ordered by increasing score\;

\Return{$(U_{(1)},\ldots,U_{(B)})$}

\caption{Direct finite-support implementation of $B$-sample MPFR}
\label{alg:direct_finite_support_MPFR}
\end{algorithm}

\begin{theorem}[Exactness of direct finite-support MPFR]
\label{thm::direct_MPFR_exactness}
Let $P$ be a distribution on a finite set $\mathcal U$. 
The output $(U_{(1)},\ldots,U_{(B)})$ of Algorithm~\ref{alg:direct_finite_support_MPFR} consists of $B$ independent samples from $P$. 
That is, for every $(u_1,\ldots,u_B)\in\mathcal U^B$,
\[
\mathbf P\{U_{(1)}=u_1,\ldots,U_{(B)}=u_B\}
=
\prod_{b=1}^B P(u_b).
\]
\end{theorem}

\begin{proof}
For each $u\in\mathcal A_P$, the sequence $(G_{u,j})_{j\geq 1}$ is the arrival process of a unit-rate Poisson process. 
Therefore,
\[
\left(\frac{G_{u,j}}{P(u)}\right)_{j\geq 1}
\]
is a Poisson process on $\mathbb R_+$ with rate $P(u)$, since
\[
\#\left\{j:\frac{G_{u,j}}{P(u)}\leq t\right\}
=
\#\{j:G_{u,j}\leq P(u)t\}
\]
is Poisson with mean $P(u)t$. 
The processes corresponding to different tokens are independent. 
Their superposition is therefore a marked Poisson process on $\mathbb R_+\times\mathcal U$ with intensity measure
\[
dt\otimes P.
\]
Equivalently, the total arrival rate is $\sum_{u\in\mathcal A_P}P(u)=1$, and the marks of the ordered arrivals are independent with common distribution $P$.

Algorithm~\ref{alg:direct_finite_support_MPFR} constructs the first $B$ candidate arrivals of each token clock. 
This is sufficient because among the first $B$ arrivals of the superposed process, no individual token clock can contribute more than $B$ arrivals. 
Hence the $B$ smallest elements in $H$ are exactly the first $B$ arrivals of the superposed marked Poisson process, and their labels are $B$ independent samples from $P$.
\end{proof}


We now describe the batched speculative decoding implementation. 
For a context $c$, let $P_k(\cdot|c)$ and $Q_k(\cdot|c)$ denote the processed target and draft distributions after applying the same logits processing, such as temperature and top-$k$ truncation. 
The draft side builds a tree of candidate continuations using Algorithm~\ref{alg:direct_finite_support_MPFR} with $Q_k$. 
The target side then evaluates all contexts needed for verification in a batched manner and uses the same MPFR clocks with $P_k$ to follow the realized target path. 
This batching changes only the order and grouping of model evaluations, not the sampling law.

\begin{algorithm}[tbp]
\SetAlgoLined
\DontPrintSemicolon
\SetAlgoNoEnd
\SetKwIF{If}{ElseIf}{Else}{if}{:}{else if}{else}{end}
\SetKwInOut{Input}{Given}

\Input{lookahead $L$, number of drafts $B$, output length $N$, target model $P$, draft model $Q$, prompt $w_{1:n}$, key $\mathsf{k}$}
% \Output{generated continuation from the processed target model $P_k$}

\While{$n<N$}{
    $h\gets w_{1:n}$, $\ell\gets \min\{L,N-n\}$\;
    $\mathcal C_0\gets\{h\}$ and $\nu(h)\gets B$\tcp*[r]{context multiplicities}

    \tcp{Draft tree construction}
    \For{$s=1$ \KwTo $\ell$}{
        $\mathcal C_s\gets\emptyset$\;
        \For{each $c\in\mathcal C_{s-1}$}{
            Generate
            \[
            (X_1(c),\ldots,X_{\nu(c)}(c))
            \gets
            \textsc{MPFR}(Q_k(\cdot|c),\Pi_{\mathsf{k}}(c),\nu(c)).
            \]
            Let $\mathcal D(c)$ be the set of distinct tokens among these samples\;
            \For{each $u\in\mathcal D(c)$}{
                $m(c,u)\gets |\{i:X_i(c)=u\}|$\;
                Add $c\|u$ to $\mathcal C_s$\;
                $\nu(c\|u)\gets \nu(c\|u)+m(c,u)$\;
            }
        }
    }

    \tcp{Batched target evaluation}
    Compute the processed target distributions $P_k(\cdot|c)$ for all contexts
    \[
    c\in \bigcup_{s=0}^{\ell}\mathcal C_s
    \]
    using batched target-model forward passes\;
    \For{each such context $c$}{
        Generate
        \[
        Y(c)\gets
        \textsc{MPFR}(P_k(\cdot|c),\Pi_{\mathsf{k}}(c),1).
        \]
    }

    \tcp{Follow the realized target path}
    $c^\star\gets h$, $\mathrm{accepted}\gets\textbf{true}$\;
    \For{$s=1$ \KwTo $\ell$}{
        Emit $Y(c^\star)$ and append it to the output\;
        \If{$Y(c^\star)\notin\mathcal D(c^\star)$}{
            $\mathrm{accepted}\gets\textbf{false}$\;
            \textbf{break}\tcp*[r]{first rejection ends the block}
        }
        \If{$n=N$}{
            \textbf{break}\;
        }
        $c^\star\gets c^\star\|Y(c^\star)$\;
    }

    \If{$\mathrm{accepted}$ \textbf{and} $n<N$}{
        Emit the precomputed bonus token $Y(c^\star)$\;
    }
}

\caption{Batched-target direct-top-$k$ MPFR speculative sampling}
\label{alg:batched_direct_topk_MPFR_spec}
\end{algorithm}

\begin{theorem}[Exactness of batched-target direct-top-$k$ MPFR decoding]
\label{thm::batched_direct_topk_MPFR_spec_exactness}
Assume that for every context $c$, the keyed source $\Pi_{\mathsf{k}}(c)$ provides independent exponential clocks across tokens and local arrival indices, and that the sources $\{\Pi_{\mathsf{k}}(c)\}_c$ are independent across distinct contexts. 
Then Algorithm~\ref{alg:batched_direct_topk_MPFR_spec} generates exactly from the autoregressive model with transition kernel $P_k(\cdot|c)$. 
That is, for every finite sequence $u_1,\ldots,u_m$,
\[
\mathbf P\{W_{n+1}=u_1,\ldots,W_{n+m}=u_m\}
=
\prod_{t=1}^{m}
P_k(u_t|w_{1:n},u_1,\ldots,u_{t-1}).
\]
\end{theorem}

\begin{proof}
Fix a context $c$. 
By Theorem~\ref{thm::direct_MPFR_exactness},
\[
Y(c)=\textsc{MPFR}(P_k(\cdot|c),\Pi_{\mathsf{k}}(c),1)
\]
has distribution $P_k(\cdot|c)$. 
The draft samples generated from $Q_k(\cdot|c)$ may use the same keyed clocks, and hence are coupled with $Y(c)$, but this does not change the marginal distribution of $Y(c)$.

Let
\[
C_t:=(w_{1:n},W_{n+1},\ldots,W_{n+t-1})
\]
be the random context before generating $W_{n+t}$. 
The event $\{C_t=c\}$ is determined by the keyed sources associated with strict prefix contexts of $c$. 
By the assumed independence across contexts, $\Pi_{\mathsf{k}}(c)$ is independent of this event. 
Therefore, conditional on $C_t=c$,
\[
W_{n+t}=Y(c)\sim P_k(\cdot|c),
\]
and hence
\[
\mathbf P\{W_{n+t}=u|C_t=c\}=P_k(u|c).
\]
Applying the chain rule gives
\[
\begin{aligned}
&\mathbf P\{W_{n+1}=u_1,\ldots,W_{n+m}=u_m\} \\
&\qquad =
\prod_{t=1}^{m}
\mathbf P\{W_{n+t}=u_t
| W_{n+1}=u_1,\ldots,W_{n+t-1}=u_{t-1}\} \\
&\qquad =
\prod_{t=1}^{m}
P_k(u_t|w_{1:n},u_1,\ldots,u_{t-1}).
\end{aligned}
\]
Thus the generated sequence has exactly the desired autoregressive law. 
Finally, batched target evaluation only changes the computational order of deterministic model evaluations. 
It does not change the context-indexed MPFR clocks or the resulting samples $Y(c)$, so it does not change the output distribution.
\end{proof}


This implementation removes the proposal-scanning overhead by simulating the token-wise Poisson clocks directly. 
For support size at most $K$, producing $B$ MPFR samples requires $KB$ exponential variables followed by a top-$B$ selection. 
The batched target implementation further reduces the number of target-model calls: instead of verifying the realized path sequentially with roughly one target forward pass per generated token, it evaluates the contexts in a speculative block together and extracts all logits needed for verification. 
If a block emits $\mathrm{BE}$ tokens on average, the target-forward cost per emitted token is reduced from approximately $1$ to approximately $1/\mathrm{BE}$, while preserving the exact target marginal distribution.



\section{Theoretic Guarantees}
\label{app::theory}


\subsection{Proof of Theorem~\ref{thm::PML_accept}}
\label{subapp::proof_PML_accept}

\begin{proof}
The proof is based on the Poisson matching lemma~\citep{li2021unified}. 
Let $p_i := P(i)$ and $q_i := Q(i)$ for $i\in\mathcal{V}$. 
Let $Y := \tilde Z_P(1)$ and $X^{(b)} := \tilde Z_Q(b)$ for $b=1,\ldots,B$ be generated from the same mapped Poisson process, and define the token-level acceptance event
\[
\mathcal{A}_t := \Big\{Y \in \{X^{(1)},\ldots,X^{(B)}\}\Big\}, 
\]
which is the acceptance event. Apply the ~\citep[Lemma 3]{li2021unified} with $j=1$ and $k=B$.
Conditioned on $Y=\tilde Z_P(1)=i$, the event $\mathcal{A}_t^c$ is exactly the event that
the first mapped $P$-sample does not appear among the first $B$ mapped $Q$-samples.
Hence
\[
\mathbf{P}(\mathcal{A}_t^c  | Y=i)
\le
\left(1-\left(1+\frac{p_i}{q_i}\right)^{-1}\right)^B
=
\left(\frac{p_i}{p_i+q_i}\right)^B.
\]
Taking complements gives
\[
\mathbf{P}(\mathcal{A}_t  | Y=i)
\ge
1-\left(\frac{p_i}{p_i+q_i}\right)^B.
\]
Finally, average over $Y\sim P$:
\[
\mathbf{P}(\mathcal{A}_t)
=
\sum_{i\in\mathcal{V}} p_i \mathbf{P}(\mathcal{A}_t  | Y=i)
\ge
\sum_{i\in\mathcal{V}}
p_i\left[1-\left(\frac{p_i}{p_i+q_i}\right)^B\right].
\]
For $B=1$, this simplifies to
\[
1-\sum_{i\in\mathcal{V}}\frac{p_i^2}{p_i+q_i}
=
\sum_{i\in\mathcal{V}}\frac{p_iq_i}{p_i+q_i}.
\qedhere
\]
\end{proof}





\subsection{Comparison with~\citep{rowan2025list} on Acceptance Probability}
\label{subapp::compare_Rowan_acpt_prob}


We compare our token-level acceptance probability with the token-level acceptance probability from \citep{rowan2025list} in two levels, and we prove that in each level our algorithm's acceptance probability is strictly better. 




\subsubsection{}

We first prove a stronger version of  Theorem~\ref{thm::PML_accept}, but not as clean as it. 


\begin{theorem}
\label{thm::PML_accept_exact}
Let $P$ be the target distribution and $Q$ be the draft distribution on a finite
vocabulary $\mathcal V$. Assume first that $P(i),Q(i)>0$ for all
$i\in\mathcal V$. For each $i\in\mathcal V$, define
\[
    \alpha_i
    :=
    P(i)\sum_{j\in\mathcal V}
    \left(
        \frac{Q(j)}{Q(i)}-\frac{P(j)}{P(i)}
    \right)_+ .
\]
If the target token is generated as $\tilde U_P(1)$ and the $B$ draft tokens are
generated as $\tilde U_Q(1),\ldots,\tilde U_Q(B)$ from the same underlying
Poisson process, then the acceptance probability of at least one token satisfies
\begin{equation}
    \label{eq::PML_accept_exact}
    \mathbf P(\mathrm{accept})
    \ge
    1-\sum_{i\in\mathcal V}
    P(i)
    \left(
        \frac{\alpha_i}{1+\alpha_i}
    \right)^B .
\end{equation}
\end{theorem}



\begin{proof}
Let the target token be
\[
    Y:=\tilde U_P(1),
\]
and let the $B$ draft tokens be
\[
    X_b:=\tilde U_Q(b),\qquad b=1,\ldots,B.
\]
The acceptance event is
\[
    \mathrm{accept}
    =
    \{Y\in\{X_1,\ldots,X_B\}\}.
\]
Conditioning on $Y=i$, the failure event is
\[
    \{Y\notin\{X_1,\ldots,X_B\}\}.
\]
By the exact $j=1$ tail formula in the generalized Poisson matching lemma,
\[
    \mathbf P\!\left(
        Y\notin\{X_1,\ldots,X_B\}
        \,\middle|\, Y=i
    \right)
    \le
    \left(
        \frac{\alpha_i}{1+\alpha_i}
    \right)^B,
\]
where
\[
    \alpha_i
    =
    P(i)\sum_{j\in\mathcal V}
    \left(
        \frac{Q(j)}{Q(i)}-\frac{P(j)}{P(i)}
    \right)_+ .
\]
Therefore,
\[
    \mathbf P(\mathrm{accept}| Y=i)
    \ge
    1-
    \left(
        \frac{\alpha_i}{1+\alpha_i}
    \right)^B .
\]
Averaging over $Y\sim P$ gives
\[
\begin{aligned}
    \mathbf P(\mathrm{accept})
    &=
    \sum_{i\in\mathcal V} P(i)
    \mathbf P(\mathrm{accept}| Y=i)\\
    &\ge
    \sum_{i\in\mathcal V} P(i)
    \left[
        1-
        \left(
            \frac{\alpha_i}{1+\alpha_i}
        \right)^B
    \right]\\
    &=
    1-\sum_{i\in\mathcal V}
    P(i)
    \left(
        \frac{\alpha_i}{1+\alpha_i}
    \right)^B .
\end{aligned}
\]
This proves~\eqref{eq::PML_accept_exact}.
\end{proof}


Moreover,
\begin{equation}
    \label{eq::PML_accept_exact_implies_weak}
    1-\sum_{i\in\mathcal V}
    P(i)
    \left(
        \frac{\alpha_i}{1+\alpha_i}
    \right)^B
    \ge
    1-\sum_{i\in\mathcal V}
    P(i)
    \left(
        \frac{P(i)}{P(i)+Q(i)}
    \right)^B .
\end{equation}
Thus~\eqref{eq::PML_accept_exact} strictly improves the weaker bound
\[
    \mathbf P(\mathrm{accept}) \ge
    1-\sum_{i\in\mathcal V} P(i)
    \left(
        \frac{P(i)}{P(i)+Q(i)}
    \right)^B
\]
whenever the inequality in~\eqref{eq::PML_accept_exact_implies_weak} is strict.


We then show above Theorem~\ref{thm::PML_accept_exact} is tighter than \citep[Proposition 2]{rowan2025list}. 

\begin{theorem}
\label{thm::exact_PML_dominates_LML}
Let $P$ be the target distribution and $Q$ be the draft distribution on a finite
vocabulary $\mathcal V$. Write
\[
    p_i := P(i),\qquad q_i := Q(i),\qquad i\in\mathcal V,
\]
and assume first that $p_i,q_i>0$ for all $i\in\mathcal V$. For each
$i\in\mathcal V$, define
\[
    \alpha_i
    :=
    p_i\sum_{j\in\mathcal V}
    \left(
        \frac{q_j}{q_i}-\frac{p_j}{p_i}
    \right)_+ .
\]
Let
\[
    L_{\mathrm{PML}}
    :=
    1-\sum_{i\in\mathcal V}
    p_i
    \left(
        \frac{\alpha_i}{1+\alpha_i}
    \right)^B
\]
be the acceptance lower bound obtained from the exact $j=1$ tail formula of the
generalized Poisson matching lemma. Let
\[
    L_{\mathrm{LML}}
    :=
    \sum_{i\in\mathcal V}
    \frac{B}{
    \sum_{j\in\mathcal V}
    \left[
        \max\left\{
            \frac{p_j}{p_i},
            \frac{q_j}{q_i}
        \right\}
        +(B-1)\frac{p_j}{p_i}
    \right]}
\]
be the lower bound in Eq.~(3) of the list matching lemma, written in the same
target--draft notation. Then
\[
    L_{\mathrm{PML}}\ge L_{\mathrm{LML}}.
\]
Moreover, if $B>1$ and $P\neq Q$, then
\[
    L_{\mathrm{PML}}>L_{\mathrm{LML}}.
\]
\end{theorem}

\begin{proof}
The proof is based on the Poisson matching lemma~\citep{li2021unified}.
Let $Y:=\tilde Z_P(1)$ and $X^{(b)}:=\tilde Z_Q(b)$ for
$b=1,\ldots,B$ be generated from the same mapped Poisson process. By the
mapping theorem, $\tilde Z_Q(1),\tilde Z_Q(2),\ldots$ are i.i.d. according
to $Q$, and $\tilde Z_P(1)\sim P$. Define the token-level acceptance event
\[
    \mathcal A_t
    :=
    \Big\{
        Y\in\{X^{(1)},\ldots,X^{(B)}\}
    \Big\}.
\]
Apply the exact $j=1$ tail formula of the generalized Poisson matching lemma
with the first distribution equal to $P$ and the second distribution equal to
$Q$. Conditioned on $Y=\tilde Z_P(1)=i$, the event $\mathcal A_t^c$ is the
event that the first mapped $P$-sample does not appear among the first $B$
mapped $Q$-samples. Hence
\[
    \mathbf{P}(\mathcal A_t^c| Y=i)
    \le
    \left(
        \frac{\alpha_i}{1+\alpha_i}
    \right)^B,
\]
where
\[
    \alpha_i
    =
    p_i\sum_{j\in\mathcal V}
    \left(
        \frac{q_j}{q_i}-\frac{p_j}{p_i}
    \right)_+ .
\]
Taking complements gives
\[
    \mathbf{P}(\mathcal A_t| Y=i)
    \ge
    1-
    \left(
        \frac{\alpha_i}{1+\alpha_i}
    \right)^B .
\]
Finally, averaging over $Y\sim P$ gives
\[
\begin{aligned}
    \mathbf{P}(\mathcal A_t)
    &=
    \sum_{i\in\mathcal V}
    p_i\mathbf{P}(\mathcal A_t| Y=i)\\
    &\ge
    \sum_{i\in\mathcal V}
    p_i
    \left[
        1-
        \left(
            \frac{\alpha_i}{1+\alpha_i}
        \right)^B
    \right]\\
    &=
    1-\sum_{i\in\mathcal V}
    p_i
    \left(
        \frac{\alpha_i}{1+\alpha_i}
    \right)^B
    =
    L_{\mathrm{PML}}.
\end{aligned}
\]

We now rewrite the LML bound in terms of the same quantities $\alpha_i$.
For each fixed $i\in\mathcal V$,
\begin{align*}
    p_i
    \sum_{j\in\mathcal V}
    \max\left\{
        \frac{p_j}{p_i},
        \frac{q_j}{q_i}
    \right\}
    &=
    \sum_{j\in\mathcal V}
    \max\left\{
        p_j,
        \frac{p_i}{q_i}q_j
    \right\} \\
    &=
    \sum_{j\in\mathcal V}
    \left[
        p_j+
        \left(
            \frac{p_i}{q_i}q_j-p_j
        \right)_+
    \right] \\
    &=
    1+
    p_i\sum_{j\in\mathcal V}
    \left(
        \frac{q_j}{q_i}-\frac{p_j}{p_i}
    \right)_+ \\
    &=
    1+\alpha_i.
\end{align*}
Therefore,
\begin{align*}
    \sum_{j\in\mathcal V}
    \left[
        \max\left\{
            \frac{p_j}{p_i},
            \frac{q_j}{q_i}
        \right\}
        +(B-1)\frac{p_j}{p_i}
    \right]
    &=
    \frac{1+\alpha_i}{p_i}
    +(B-1)\frac{1}{p_i} \\
    &=
    \frac{B+\alpha_i}{p_i}.
\end{align*}
Thus Eq.~(3) of the list matching lemma can be written as
\[
    L_{\mathrm{LML}}
    =
    \sum_{i\in\mathcal V}
    p_i\frac{B}{B+\alpha_i}.
\]

It remains to compare the two expressions term by term. For every
$\alpha\ge 0$, we claim that
\[
    1-\left(\frac{\alpha}{1+\alpha}\right)^B
    \ge
    \frac{B}{B+\alpha}.
\]
If $\alpha=0$, both sides are equal to $1$. If $\alpha>0$, the inequality is
equivalent to
\[
    \left(\frac{\alpha}{1+\alpha}\right)^B
    \le
    \frac{\alpha}{B+\alpha},
\]
or equivalently
\[
    \left(1+\frac{1}{\alpha}\right)^B
    \ge
    1+\frac{B}{\alpha}.
\]
This follows immediately from the binomial theorem:
\[
    \left(1+\frac{1}{\alpha}\right)^B
    =
    1+\frac{B}{\alpha}
    +\sum_{m=2}^B \binom{B}{m}\alpha^{-m}
    \ge
    1+\frac{B}{\alpha}.
\]
Hence, for each $i\in\mathcal V$,
\[
    p_i
    \left[
        1-\left(\frac{\alpha_i}{1+\alpha_i}\right)^B
    \right]
    \ge
    p_i\frac{B}{B+\alpha_i}.
\]
Summing over $i\in\mathcal V$ yields
\[
    L_{\mathrm{PML}}
    \ge
    L_{\mathrm{LML}}.
\]

Finally, suppose $B>1$ and $P\neq Q$. We show that the inequality is strict.
If $\alpha_i=0$ for every $i\in\mathcal V$, then for every $i,j\in\mathcal V$,
\[
    \frac{q_j}{q_i}
    \le
    \frac{p_j}{p_i}.
\]
Equivalently,
\[
    \frac{q_j}{p_j}
    \le
    \frac{q_i}{p_i}.
\]
Since this holds for every pair $(i,j)$, all ratios $q_i/p_i$ must be equal to
a common constant. Because both $P$ and $Q$ are probability distributions, this
constant must be $1$, and hence $P=Q$, a contradiction. Therefore, when
$P\neq Q$, there exists at least one $i$ such that $\alpha_i>0$.

For $B>1$ and $\alpha>0$, the binomial expansion above is strict:
\[
    \left(1+\frac{1}{\alpha}\right)^B
    >
    1+\frac{B}{\alpha}.
\]
Hence
\[
    1-\left(\frac{\alpha}{1+\alpha}\right)^B
    >
    \frac{B}{B+\alpha}.
\]
Applying this to an index $i$ with $\alpha_i>0$ gives a strict term in the
sum. Therefore
\[
    L_{\mathrm{PML}}
    >
    L_{\mathrm{LML}}.
\]
This completes the proof.
\end{proof}




\subsubsection{}

We then show the clean version,  Theorem~\ref{thm::PML_accept}, is already tighter than \citep[Eq. (4)]{rowan2025list}. 
	
	\begin{theorem}
\label{thm::PML_weak_dominates_relaxed_LML}
Let $P$ be the target distribution and $Q$ be the draft distribution on a finite
vocabulary $\mathcal V$. Write
\[
    p_i:=P(i),\qquad q_i:=Q(i),\qquad i\in\mathcal V.
\]
Assume first that $p_i,q_i>0$ for all $i\in\mathcal V$. Define
\[
    L_{\mathrm{PML,weak}}
    :=
    1-\sum_{i\in\mathcal V}
    p_i
    \left(
        \frac{p_i}{p_i+q_i}
    \right)^B
\]
and
\[
    L_{\mathrm{LML,rel}}
    :=
    \sum_{i\in\mathcal V}
    p_i
    \left(
        1+\frac{p_i}{Bq_i}
    \right)^{-1}
    =
    \sum_{i\in\mathcal V}
    p_i
    \frac{Bq_i}{p_i+Bq_i}.
\]
Then
\[
    L_{\mathrm{PML,weak}}
    \ge
    L_{\mathrm{LML,rel}}.
\]
Moreover, if $B>1$, then the inequality is strict.
\end{theorem}

\begin{proof}
The proof is based on the Poisson matching lemma~\citep{li2021unified}.
Let $Y:=\tilde Z_P(1)$ and let
\[
    X^{(b)}:=\tilde Z_Q(b),\qquad b=1,\ldots,B,
\]
be generated from the same mapped Poisson process. Define the token-level
acceptance event
\[
    \mathcal A_t
    :=
    \Big\{
        Y\in\{X^{(1)},\ldots,X^{(B)}\}
    \Big\}.
\]
Applying the relaxed Poisson matching bound with $j=1$ and $k=B$, conditioned
on $Y=\tilde Z_P(1)=i$, gives
\[
    \mathbf{P}(\mathcal A_t^c| Y=i)
    \le
    \left(
        1-\left(1+\frac{p_i}{q_i}\right)^{-1}
    \right)^B
    =
    \left(
        \frac{p_i}{p_i+q_i}
    \right)^B .
\]
Taking complements yields
\[
    \mathbf{P}(\mathcal A_t| Y=i)
    \ge
    1-
    \left(
        \frac{p_i}{p_i+q_i}
    \right)^B .
\]
Averaging over $Y\sim P$ gives
\[
\begin{aligned}
    \mathbf{P}(\mathcal A_t)
    &=
    \sum_{i\in\mathcal V}
    p_i\mathbf{P}(\mathcal A_t| Y=i)\\
    &\ge
    \sum_{i\in\mathcal V}
    p_i
    \left[
        1-
        \left(
            \frac{p_i}{p_i+q_i}
        \right)^B
    \right]\\
    &=
    1-\sum_{i\in\mathcal V}
    p_i
    \left(
        \frac{p_i}{p_i+q_i}
    \right)^B
    =
    L_{\mathrm{PML,weak}}.
\end{aligned}
\]

We now compare this lower bound with the relaxed version of the list matching
lemma. In the notation of the list matching lemma paper, the proposal
distribution is denoted by $p$ and the target distribution by $q$. Since here
we use $P$ for the target and $Q$ for the draft, their relaxed conditional
bound, Eq.~(4), becomes
\[
    \mathbf{P}(\mathcal A_t| Y=i)
    \ge
    \left(
        1+\frac{p_i}{Bq_i}
    \right)^{-1}
    =
    \frac{Bq_i}{p_i+Bq_i}.
\]
Thus the corresponding unconditional relaxed LML bound is
\[
    L_{\mathrm{LML,rel}}
    =
    \sum_{i\in\mathcal V}
    p_i
    \frac{Bq_i}{p_i+Bq_i}.
\]

It remains to compare the two conditional terms. Let
\[
    r_i:=\frac{p_i}{q_i}.
\]
Then
\[
    1-
    \left(
        \frac{p_i}{p_i+q_i}
    \right)^B
    =
    1-
    \left(
        \frac{r_i}{1+r_i}
    \right)^B,
\]
whereas
\[
    \frac{Bq_i}{p_i+Bq_i}
    =
    \frac{B}{B+r_i}.
\]
Hence it suffices to prove that, for every $r>0$,
\[
    1-
    \left(
        \frac{r}{1+r}
    \right)^B
    \ge
    \frac{B}{B+r}.
\]
Equivalently,
\[
    \left(
        \frac{r}{1+r}
    \right)^B
    \le
    \frac{r}{B+r}.
\]
Since $r>0$, this is equivalent to
\[
    \left(
        1+\frac{1}{r}
    \right)^B
    \ge
    1+\frac{B}{r}.
\]
This follows from the binomial theorem:
\[
    \left(
        1+\frac{1}{r}
    \right)^B
    =
    1+\frac{B}{r}
    +
    \sum_{m=2}^B
    \binom{B}{m}r^{-m}
    \ge
    1+\frac{B}{r}.
\]
Therefore, for each $i\in\mathcal V$,
\[
    1-
    \left(
        \frac{p_i}{p_i+q_i}
    \right)^B
    \ge
    \frac{Bq_i}{p_i+Bq_i}.
\]
Multiplying by $p_i$ and summing over $i\in\mathcal V$ gives
\[
    L_{\mathrm{PML,weak}}
    \ge
    L_{\mathrm{LML,rel}}.
\]

Finally, if $B>1$, then for every $r>0$,
\[
    \sum_{m=2}^B
    \binom{B}{m}r^{-m}
    >
    0.
\]
Thus the above inequality is strict for every $i$ with $p_i,q_i>0$. Under full
support, all such terms have positive weight $p_i$, and hence
\[
    L_{\mathrm{PML,weak}}
    >
    L_{\mathrm{LML,rel}}.
\]
This completes the proof.
\end{proof}






\subsection{Generalized Theoretic Guarantee and Proof}
\label{subapp::proof_PML_accept_generalized}

\begin{theorem}
\label{thm:accept_bound_pfr_active}
Consider the exact draft-indexed active-set version of multi-draft speculative sampling
with total draft count $B$.
Fix a current decoding step $t$, and let $S_t\subseteq\{1,\ldots,B\}$ be the active set
before sampling the next token. Write $J_t := |S_t|$. 
Let the current realized context be $c_t$, and similarly define $P_t := \mathcal{M}_t(\cdot | c_t), Q_t := \mathcal{M}_d(\cdot | c_t)$. 
For $i\in\mathcal{V}$, write $p_i^{(t)} := P_t(i)$ and $q_i^{(t)} := Q_t(i)$. 
Then the token-level acceptance probability at step $t$ satisfies
\[
\mathbf{P}(\mathcal{A}_t  | c_t,S_t)
\ge
\sum_{i\in\mathcal{V}}
p_i^{(t)}
\left[
1-\left(1-\left(1+\frac{p_i^{(t)}}{q_i^{(t)}}\right)^{-1}\right)^{J_t}
\right]
=
1-\sum_{i\in\mathcal{V}}
p_i^{(t)}
\left(\frac{p_i^{(t)}}{p_i^{(t)}+q_i^{(t)}}\right)^{J_t}.
\]
In particular, at the first step of each speculative block one has $J_t=B$, so
\[
\mathbf{P}(\mathcal{A}_t  | c_t)
\ge
1-\sum_{i\in\mathcal{V}}
p_i^{(t)}
\left(\frac{p_i^{(t)}}{p_i^{(t)}+q_i^{(t)}}\right)^B.
\]
\end{theorem}

\begin{proof}
Conditioned on the realized context $c_t$ and the active set $S_t$, the exact active-set
algorithm uses only the $J_t$ active draft streams for both target-side selection and acceptance.
Inactive drafts do not participate in the current-step minimum and therefore do not affect the
current-step coupling.

Hence, after conditioning on $(c_t,S_t)$, the current step is identical to the setting of
Theorem~\ref{thm::PML_accept} with the number of drafts equal to $J_t$.
Applying Theorem~\ref{thm::PML_accept} with $B$ replaced by $J_t$ gives
\[
\mathbf{P}(\mathcal{A}_t  | c_t,S_t)
\ge
1-\sum_{i\in\mathcal{V}}
p_i^{(t)}
\left(\frac{p_i^{(t)}}{p_i^{(t)}+q_i^{(t)}}\right)^{J_t}.
\]
At the first step of a speculative block, all drafts are active, so $J_t=B$, which yields
the final claim.
\end{proof}




\subsection{Extension of Theorem \ref{thm::PML_accept_exact} to non-identically distributed proposals}
\label{app::non_id_prop}

\begin{proposition}
\label{prop::heterogeneous_PML_bound}
Let $P$ be the target distribution on a finite vocabulary $\mathcal V$, and let
$Q_1,\ldots,Q_B$ be $B$ proposal distributions. Write
\[
    p_i:=P(i),\qquad q_{b,i}:=Q_b(i),
    \qquad i\in\mathcal V,\ b\in[B].
\]
Assume first that $p_i,q_{b,i}>0$ for all $i$ and $b$. For each $i\in\mathcal V$
and $b\in[B]$, define
\[
    \alpha_i^{(b)}
    :=
    p_i\sum_{j\in\mathcal V}
    \left(
        \frac{q_{b,j}}{q_{b,i}}
        -
        \frac{p_j}{p_i}
    \right)_+ .
\]
Then there exists a common-randomness coupling such that $Y\sim P$,
$X^{(b)}\sim Q_b$ for each $b\in[B]$, the proposals
$X^{(1)},\ldots,X^{(B)}$ are mutually independent, and
\begin{equation}
\label{eq::heterogeneous_PML_conditional}
    \Pr\!\left[
        Y\in\{X^{(1)},\ldots,X^{(B)}\}
        \,\middle|\,Y=i
    \right]
    \ge
    \sum_{b=1}^B
    \frac{1}{B+\alpha_i^{(b)}} .
\end{equation}
Consequently,
\begin{equation}
\label{eq::heterogeneous_PML_unconditional}
    \Pr\!\left[
        Y\in\{X^{(1)},\ldots,X^{(B)}\}
    \right]
    \ge
    \sum_{i\in\mathcal V}
    p_i
    \sum_{b=1}^B
    \frac{1}{B+\alpha_i^{(b)}} .
\end{equation}
\end{proposition}

\begin{proof}
The proof is based on the Poisson matching lemma~\citep{li2021unified}. 
The main difference from the identical-proposal case is that there is no single
proposal ordering when the proposal laws are $Q_1,\ldots,Q_B$. We therefore
augment the sample space by a block label. Define
\[
    \widetilde{\mathcal V}:=\mathcal V\times[B],
\]
and introduce the augmented target and proposal distributions
\[
    \widetilde P(i,b):=\frac{p_i}{B},
    \qquad
    \widetilde Q_b(i,\ell):=\mathbf 1\{\ell=b\}q_{b,i}.
\]
Let
\[
    (Y,J):=\tilde Z_{\widetilde P}(1),
    \qquad
    (X^{(b)},b):=\tilde Z_{\widetilde Q_b}(1),
    \quad b=1,\ldots,B,
\]
be generated from the same underlying Poisson process on
$\widetilde{\mathcal V}\times\mathbb R_{\ge 0}$. By the mapping theorem,
$(Y,J)\sim\widetilde P$, and hence $Y\sim P$ and $J$ is uniform on $[B]$ and
independent of $Y$. Also, for each $b$, $(X^{(b)},b)\sim\widetilde Q_b$, so
$X^{(b)}\sim Q_b$. Since the laws $\widetilde Q_b$ are supported on disjoint
blocks, the corresponding proposal samples are mutually independent.

Fix $i\in\mathcal V$. Define the acceptance event
\[
    \mathcal A
    :=
    \left\{
        Y\in\{X^{(1)},\ldots,X^{(B)}\}
    \right\}.
\]
For each $b\in[B]$, let
\[
    \mathcal A_b
    :=
    \left\{
        (Y,J)=(X^{(b)},b)
    \right\}.
\]
Then $\mathcal A\supseteq \bigcup_{b=1}^B\mathcal A_b$, and the events
$\mathcal A_1,\ldots,\mathcal A_B$ are pairwise disjoint because they occur in
different blocks. Therefore,
\begin{align}
    \mathbf{P}(\mathcal A| Y=i)
    &\ge
    \sum_{b=1}^B
    \mathbf{P}(\mathcal A_b| Y=i) \notag\\
    &=
    \sum_{b=1}^B
    \mathbf{P}(J=b| Y=i)
    \mathbf{P}(\mathcal A_b| Y=i,J=b) \notag\\
    &=
    \frac{1}{B}
    \sum_{b=1}^B
    \Pr\!\left[
        \tilde Z_{\widetilde P}(1)=\tilde Z_{\widetilde Q_b}(1)
        \,\middle|\,
        \tilde Z_{\widetilde P}(1)=(i,b)
    \right].
\label{eq::heterogeneous_reduction}
\end{align}

We now apply the exact $j=1$ tail formula in the generalized Poisson matching
lemma to the pair $(\widetilde P,\widetilde Q_b)$. Conditioned on
$\tilde Z_{\widetilde P}(1)=(i,b)$, the local PML parameter is
\[
\begin{aligned}
    \widetilde\alpha_i^{(b)}
    &:=
    \widetilde P(i,b)
    \sum_{\ell=1}^B
    \sum_{j\in\mathcal V}
    \left(
        \frac{\widetilde Q_b(j,\ell)}{\widetilde Q_b(i,b)}
        -
        \frac{\widetilde P(j,\ell)}{\widetilde P(i,b)}
    \right)_+  \\
    &=
    \frac{p_i}{B}
    \sum_{j\in\mathcal V}
    \left(
        \frac{q_{b,j}}{q_{b,i}}
        -
        \frac{p_j}{p_i}
    \right)_+  \\
    &=
    \frac{\alpha_i^{(b)}}{B}.
\end{aligned}
\]
Hence
\[
    \Pr\!\left[
        \tilde Z_{\widetilde P}(1)=\tilde Z_{\widetilde Q_b}(1)
        \,\middle|\,
        \tilde Z_{\widetilde P}(1)=(i,b)
    \right]
    \ge
    1-\frac{\widetilde\alpha_i^{(b)}}{1+\widetilde\alpha_i^{(b)}}
    =
    \frac{B}{B+\alpha_i^{(b)}}.
\]
Substituting this into~\eqref{eq::heterogeneous_reduction} gives
\[
    \mathbf{P}(\mathcal A| Y=i)
    \ge
    \frac1B
    \sum_{b=1}^B
    \frac{B}{B+\alpha_i^{(b)}}
    =
    \sum_{b=1}^B
    \frac{1}{B+\alpha_i^{(b)}}.
\]
This proves~\eqref{eq::heterogeneous_PML_conditional}. Averaging over
$Y\sim P$ gives~\eqref{eq::heterogeneous_PML_unconditional}.
\end{proof}


\section{Drafter Invariance}


\subsection{Definitions}

For a speculative sampling algorithm $\mathcal{A}$ with target model $\mathcal{M}_t$, draft model $\mathcal{M}_d$, shared randomness $\mathcal{R}$ and context $x_{:t}$, let $Y_{1:\tau}$ denote the emitted output sequence of $\mathcal{A}$ in the current speculative block.


\begin{definition}[\citep{daliri2025coupling}]
\label{def:strong_drafter_invariance_single}
The algorithm $\mathcal{A}$ is said to be \emph{strongly drafter invariant} if for every $1\le j\le \tau$,
\begin{equation*}
\mathbf{P}\big\{
Y_{1:j}=y_{1:j}\,\big|\,\mathcal{R},\,x_{:t}
\big\} =
\mathbf{P}\big\{
\widetilde{Y}_{1:j}=y_{1:j}\,\Big|\,\mathcal{R},\,x_{:t}
\big\}. 
\end{equation*}
\end{definition}


\begin{definition}[\citep{rowan2025list}]
\label{def:conditional_drafter_invariance_single}
The algorithm $\mathcal{A}$ is said to be \emph{conditional drafter invariant} if for every $1\le j\le \tau$, letting $X_{1:L} := X_{1:L}(\mathcal{M}_d)$ denote the length-$L$ draft sequence by $\mathcal{M}_d$, we have
\begin{equation*}
\mathbf{P}\big\{
Y_{1:j}=y_{1:j}\,\big|\,\mathcal{R},\,x_{:t},\,X_{1:L}=x_{1:L}
\big\}  =
\mathbf{P}\big\{
\widetilde{Y}_{1:j}=y_{1:j}\,\big|\,\mathcal{R},\,x_{:t},\,\widetilde{X}_{1:L}=x_{1:L}
\big\}. 
\end{equation*}
\end{definition}



\subsubsection{Proof of Proposition~\ref{prop::single_dft_st_tm_inv}}
\label{subsubapp::proof_single_dft_st_tm_inv}

We rephrase Proposition~\ref{prop::single_dft_st_tm_inv} to a more detailed version as follows. 


\begin{proposition}[Stopping-time drafter invariance]
\label{prop:single_pfr_stopping_time_invariance}
Let $\tau$ denote the number of tokens emitted in one speculative block of
Algorithm~\ref{alg:single_PFR_Spec_Water_simplified}. 
Fix the initial prompt $h:=w_{1:n}$ and the shared randomness
\[
\mathcal{R}:=\{\,G(\mathsf{k},c): c \text{ is any context}\,\}.
\]
Define the \emph{target-side PFR chain} recursively by
\[
Z_1:=\textsc{PFR}\bigl(P(\cdot | h),\,G(\mathsf{k},h)\bigr),
\]
and for $s\ge 2$,
\[
Z_s:=\textsc{PFR}\bigl(P(\cdot | h\|Z_{1:s-1}),\,G(\mathsf{k},\,h\|Z_{1:s-1})\bigr).
\]
Then, for any stopping-time value $\tau_0$ and every $1\le j\le \tau_0$,
\[
Y_{1:j}=Z_{1:j}
\qquad\text{almost surely on the event }\{\tau=\tau_0\}.
\]
Equivalently, for every output prefix $y_{1:j}$,
\[
\mathbf{P}\bigl\{
Y_{1:j}=y_{1:j}\,\big|\,\mathcal{R},\,h,\,\tau=\tau_0
\bigr\}
=
\mathds{1}\{Z_{1:j}=y_{1:j}\}.
\]
In particular, for any two draft models $Q$ and $\widetilde Q$,
\[
\mathbf{P}\bigl\{
Y_{1:j}=y_{1:j}\,\big|\,\mathcal{R},\,h,\,\tau=\tau_0
\bigr\}
=
\mathbf{P}\bigl\{
\widetilde Y_{1:j}=y_{1:j}\,\big|\,\mathcal{R},\,h,\,\widetilde\tau=\tau_0
\bigr\},
\]
so Algorithm~\ref{alg:single_PFR_Spec_Water_simplified} is stopping-time drafter invariant in the sense of Definition~\ref{def:st_tm_drafter_invariance_single}.
\end{proposition}

\begin{proof}
Fix $(\mathcal{R},h)$ and a stopping-time value $\tau_0$.
We first show by induction on $s$ that, on the event $\{\tau\ge s\}$,
\[
Y_s=Z_s.
\]

For $s=1$, the first context is
\[
c_1=h.
\]
Hence
\[
Y_1
=
\textsc{PFR}\bigl(P(\cdot | c_1),\,G(\mathsf{k},c_1)\bigr)
=
\textsc{PFR}\bigl(P(\cdot | h),\,G(\mathsf{k},h)\bigr)
=
Z_1.
\]

Now suppose $s\ge 2$ and that $Y_r=Z_r$ holds on $\{\tau\ge r\}$ for all
$r=1,\ldots,s-1$. On the event $\{\tau\ge s\}$, the block has reached step $s$,
so the first $s-1$ draft tokens were accepted. Therefore
\[
\tilde w_{1:s-1}=Y_{1:s-1}.
\]
By the induction hypothesis, on $\{\tau\ge s\}$ we also have
\[
Y_{1:s-1}=Z_{1:s-1},
\]
and hence
\[
c_s
=
h\|\tilde w_{1:s-1}
=
h\|Y_{1:s-1}
=
h\|Z_{1:s-1}.
\]
Thus,
\[
Y_s
=
\textsc{PFR}\bigl(P(\cdot | c_s),\,G(\mathsf{k},c_s)\bigr)
=
\textsc{PFR}\bigl(P(\cdot | h\|Z_{1:s-1}),\,G(\mathsf{k},\,h\|Z_{1:s-1})\bigr)
=
Z_s.
\]
This completes the induction.

Now let $1\le j\le \tau_0$. On the event $\{\tau=\tau_0\}$ we have
$\tau\ge j$, so the above induction gives
\[
Y_{1:j}=Z_{1:j}
\qquad\text{almost surely on }\{\tau=\tau_0\}.
\]
Therefore,
\[
\mathbf{P}\bigl\{
Y_{1:j}=y_{1:j}\,\big|\,\mathcal{R},\,h,\,\tau=\tau_0
\bigr\}
=
\mathds{1}\{Z_{1:j}=y_{1:j}\}.
\]
Since the right-hand side depends only on $(\mathcal{R},h,\tau_0)$ and not on
the draft model, the same conditional law holds for any other draft model
$\widetilde Q$ conditioned on $\widetilde\tau=\tau_0$. This is exactly
stopping-time drafter invariance.
\end{proof}


\subsection{Multi-Draft Drafter invariance}
\label{subapp::multi_draft_invariance}

For multi-draft speculative sampling scheme $\mathcal{A}$ with target model $\mathcal{M}_t$, draft models $\mathcal{M}_d^{(1)},\ldots,\mathcal{M}_d^{(B)}$, context $x_{:t}$, shared randomness $\mathcal{R}$ and $b\in[B]$, let $X_{1:L}^{(b)} := X_{1:L}(\mathcal{M}_d^{(b)})$ denote the length-$L$ draft sequence by $\mathcal{M}_d^{(b)}$ and $Y_{1:\tau}$ denote the emitted output sequence of $\mathcal{A}$ in the current speculative block.


\begin{definition}[Strong drafter invariance]
\label{def:drafter_invariance}
The algorithm $\mathcal{A}$ is \emph{drafter invariant} if for every $1\le j\le \tau$,
every initial context $x_{:t}$, every realization of the shared randomness $\mathcal{R}$,
and every output prefix $y_{1:j}$,
\begin{align*}
&\mathbf{P}\Bigl\{
Y_{1:j}=y_{1:j}\,\Big|\,\mathcal{R},\,x_{:t}
\Bigr\} =
\mathbf{P}\Bigl\{
\widetilde{Y}_{1:j}=y_{1:j}\,\Big|\,\mathcal{R},\,x_{:t}
\Bigr\},
\end{align*}
for any two choices of draft models
\[
\mathcal{M}_d^{(1)},\ldots,\mathcal{M}_d^{(B)}
\quad\text{and}\quad
\widetilde{\mathcal{M}}_d^{(1)},\ldots,\widetilde{\mathcal{M}}_d^{(B)}. 
\]
\end{definition}


\begin{definition}[Conditional drafter invariance]
\label{def:conditional_drafter_invariance}
We say that the algorithm is \emph{conditionally drafter invariant} if, for every $1\le j\le \tau$,
every initial context $x_{:t}$, every realization of the shared randomness $\mathcal{R}$,
every draft sequences $x_{1:L}^{(1)},\ldots,x_{1:L}^{(B)}$, and every output prefix $y_{1:j}$,
\begin{align*}
&\mathbf{P}\Bigl\{
Y_{1:j}=y_{1:j}\,\Big|\,
\mathcal{R},\,x_{:t},\,
X_{1:L}^{(1)}=x_{1:L}^{(1)},\ldots,
X_{1:L}^{(B)}=x_{1:L}^{(B)}
\Bigr\} \\
&=
\mathbf{P}\Bigl\{
\widetilde{Y}_{1:j}=y_{1:j}\,\Big|\,
\mathcal{R},\,x_{:t},\,
\widetilde{X}_{1:L}^{(1)}=x_{1:L}^{(1)},\ldots,
\widetilde{X}_{1:L}^{(B)}=x_{1:L}^{(B)}
\Bigr\},
\end{align*}
for any two choices of draft models
\[
\mathcal{M}_d^{(1)},\ldots,\mathcal{M}_d^{(B)}
\quad\text{and}\quad
\widetilde{\mathcal{M}}_d^{(1)},\ldots,\widetilde{\mathcal{M}}_d^{(B)}.
\]
That is, conditioned on the shared randomness, the initial context, and the realized draft sequences, the conditional law of the emitted output prefix does not depend on which draft models generated those sequences.
\end{definition}




\begin{definition}[stopping time drafter invariance]
\label{def:block_length_drafter_invariance}
Consider a speculative sampling algorithm with target model $\mathcal{M}_t$,
draft models $\mathcal{M}_d^{(1)},\ldots,\mathcal{M}_d^{(B)}$, initial context $x_{:t}$,
shared randomness $\mathcal{R}$, and block output $Y_{1:\tau}$, where $\tau$ is the number
of emitted tokens in the current speculative block.

We say that the algorithm is \emph{stopping time drafter invariant} if, for every $1\le j\le \tau$,
every initial context $x_{:t}$, every realization of the shared randomness $\mathcal{R}$,
every block length value $\tau_0$, and every output prefix $y_{1:j}$,
\begin{align*}
&\mathbf{P}\Bigl\{
Y_{1:j}=y_{1:j}\,\Big|\,\mathcal{R},\,x_{:t},\,\tau=\tau_0
\Bigr\} \\
&=
\mathbf{P}\Bigl\{
\widetilde{Y}_{1:j}=y_{1:j}\,\Big|\,\mathcal{R},\,x_{:t},\,\widetilde{\tau}=\tau_0
\Bigr\},
\end{align*}
for any two choices of draft models
\[
\mathcal{M}_d^{(1)},\ldots,\mathcal{M}_d^{(B)}
\quad\text{and}\quad
\widetilde{\mathcal{M}}_d^{(1)},\ldots,\widetilde{\mathcal{M}}_d^{(B)}.
\]
That is, conditioned on the shared randomness, the initial context, and the realized block length,
the conditional law of the emitted output prefix does not depend on the draft models.
\end{definition}



\section{Experiments}
\label{sec:experiments}

We evaluate the algorithm along the two axes that mirror the contributions
claimed in Section~\ref{sec::intro}: (i) the \textbf{single-draft regime}
($B=1$), where the goal is to certify that our keyed Poisson-process sampler
preserves acceptance and quality while carrying watermark signal -
contradicting the Pareto frontier identified by~\citet{hu2024inevitable};
and (ii) the \textbf{multi-draft regime} ($B>1$), where we both
extend the unbiasedness/strength guarantees to larger $B$ and benchmark
sampling efficiency against the list-coupling scheme
of~\citet{rowan2025list}.

\subsection{Setup}
\label{sec:setup}

\paragraph{Models.}
We use two target/drafter pairs that span the two model families used
in the prior watermark and list-coupling literature.
Our headline configuration follows~\citet{rowan2025list}: target
\texttt{Qwen2.5-7B-Instruct} with drafter \texttt{Qwen2.5-0.5B-Instruct}.
For instruction-following diagnostics we additionally use
\texttt{Vicuna-7B-v1.5} paired with a same-family small drafter,
with the FastChat chat template injected explicitly since the released
checkpoint does not ship one.
All models run in \texttt{float16} on a single GPU.

\paragraph{Datasets.}
We evaluate on three tasks chosen to match the coverage of prior work:
\textsc{CNN/DailyMail} (summarization, prompted with the
\texttt{System:/INPUT:/OUTPUT:} template of~\citet{hu2024inevitable}),
\textsc{ELI5} (open-ended explanation).

\paragraph{Sampling parameters.}
Unless otherwise stated, we use $\mathrm{top}_k = 50$, $T = 1.0$, and \texttt{max\_new\_tokens} $= 128$. The watermark key
is fixed across all runs so that detection is deterministic.

\paragraph{Sweeps.}
For single draft we sweep lookahead $L \in \{1,2,3,4\}$, matching
\citet{hu2024inevitable}.
For multi-draft we fix $L = 4$ and sweep
$B \in \{2, 4, 6, 8\}$, matching~\citet{rowan2025list}.
Each $(\textsc{model}, \textsc{dataset}, \textsc{decoder}, L, B)$
configuration is run with multiple independent seeds and reported as
$\textrm{mean} \pm 95\%$ CI.

\paragraph{Baselines.}
We organize baselines by the role they play in our claim:
\begin{itemize}
  \item \textbf{Quality anchor} - \emph{Basic Sampling}: temperature-sampled
    decoding with no speculation and no watermark. Provides the LOGPPL
    reference value.
  \item \textbf{Watermark-only} - \emph{VUW}~\citep{hu2024inevitable}: the
    vanilla unbiased watermark with no speculative decoding. Establishes
    the detection signal achievable without speculation.
  \item \textbf{Speculation-only} - \emph{Standard
    SpecSampling}~\citep{leviathan2023fast} at $B=1$ (denoted \texttt{VSps}
    in our tables) and the
    \emph{InvariantMultiDraftStrategy} of~\citet{rowan2025list} at $B>1$:
    speculative decoding with no watermark. These set the AATPS / token-rate
    ceilings that any watermarked scheme must approach.
  \item \textbf{Trade-off frontier} - \emph{MWS}, \emph{MSE}, and
    \emph{MSE\_PSEUDO}~\citep{hu2024inevitable}: prior schemes that explicitly
    trade efficiency for watermark strength or vice versa
    (\texttt{MSE\_PSEUDO} is the Algorithm 1 dual-key + pseudorandom-$r$
    variant). They mark the Pareto frontier we aim to circumvent.
  \item \textbf{Multi-draft list-coupling} -~\citet{rowan2025list}:
    the closest unwatermarked analogue of our multi-draft sampler.
    We benchmark Block Efficiency (BE) and relative token rate following
    their reporting protocol.
\end{itemize}

\paragraph{Metrics.}
We separate the three concerns explicitly, following~\citet{hu2024inevitable}
and~\citet{he2026improving}:
\begin{itemize}
  \item \emph{Efficiency.} \textbf{AATPS} (average accepted tokens per step)
    and \textbf{Block Efficiency} (mean accepted draft tokens per verifier
    call, following~\citet{rowan2025list}). Token rate is reported in
    \emph{two} forms: absolute \texttt{tok/s} (with hardware footnote) and
    \% speedup over Standard SpecSampling at $B=1$, so that hardware
    differences are neutralized.
  \item \emph{Detectability.} \textbf{ANLPPT}
    - the average negative log-$p$-value per token under three score
    variants (U-score, Li, PL), computed by Chernoff-bound inversion of
    the per-token log-MGF as in~\citet{hu2024inevitable}.

  \item \emph{Quality.} \textbf{LOGPPL} (per-token log-perplexity), reported
    as a separate audit column. We require LOGPPL parity with
    \emph{Basic Sampling} to certify unbiasedness; we do not optimize
    against this metric.
\end{itemize}

\paragraph{Reproducibility.}
The exact $(seed, \textsc{model}, \textsc{dataset}, \textsc{decoder}, L, B)$
configurations, software versions (PyTorch, \texttt{transformers}, attention
kernel), GPU model, and per-seed result files are released with the paper.











\subsection{Single-draft experiments: circumventing the inevitable trade-off}
\label{sec:single_draft}

\paragraph{Goal.}
\citet{hu2024inevitable} prove that, under speculative sampling,
any unbiased reweighting scheme must give up either watermark strength
(MSE-style) or acceptance rate (MWS-style); their Theorem~3.1 forbids
schemes that strictly improve both axes simultaneously.
We claim that our Poisson-process construction violates this dichotomy:
it is provably unbiased on the actually-sampled distribution while
preserving acceptance.

\paragraph{Protocol.}
We follow the protocol of~\citet{hu2024inevitable}: the same metric
battery (AATPS, ANLPPT, LOGPPL) and lookahead range $L \in \{1,2,3,4\}$.
We run on \texttt{Qwen2.5-7B-Instruct}/\texttt{0.5B-Instruct} and
\texttt{Vicuna-7B-v1.5} on \textsc{CNN/DailyMail}.
On top of the Hu-Huang decoder set we add (i) our PFR sampler
and (ii) \emph{Basic Sampling} as the unbiasedness anchor.

\paragraph{Headline figure.}
The headline plot is a panel grid in the style
of~\citet[Fig.~2]{hu2024inevitable}.
Rows correspond to score variants (U-score / maximin-LLR);
columns correspond to dataset.
Within each panel, every method is drawn as a curve in
$(\mathrm{AATPS}, \mathrm{ANLPPT})$ space sweeping $L$.
Our scheme dominates the MWS/MSE Pareto frontier on every panel:
at every $L$ it sits both higher (greater detection signal) and to the
right (greater acceptance) than the trade-off baselines, providing visual
evidence that we have crossed the frontier of~\citet{hu2024inevitable}.

\paragraph{Numeric appendix table.}
Mirroring the visual main result we report a numeric table with rows
$=(L, \textsc{decoder})$ and columns $=$ AATPS, per-token time, LOGPPL,
ANLPPT(U-score), and ANLPPT(maximin-LLR), each as
$\textrm{mean} \pm 95\%$ CI.

\paragraph{Unbiasedness audit.}
We require our scheme to match \emph{Basic Sampling} in LOGPPL within
statistical noise; this is the strongest empirical signature of breaking
the Hu-Huang trade-off, since any biased scheme would shift LOGPPL.
The audit is reported as a separate column in the appendix table.

\begin{figure}[htpb]
    \centering
    \includegraphics[scale=0.35]{images/qwen_eli5_figure1.png}
    \caption{
    The tradeoff between watermark strength and sampling efficiency applies to the algorithms in~\citep{hu2024inevitable}: Vanilla Speculative Sampling (VSpS), Maintain Watermark Strength (MWS), and Maintain Sampling Efficiency (MSE), but \textbf{not} to our algorithm based on the PFR. 
    Even only using the single draft version, our algorithm achieves better performance on both the Average Accepted Tokens Per Step (AATPS) and the watermark score, and attaching watermark does not decrease AATPS. 
    The figure is from using Qwen2.5-7B and 0.5B-Instruct as the model pair on dataset \textsc{eli5}~\citep{fan2019eli5}.
    }
    \label{fig::temp2}
\end{figure}


\begin{figure}[htpb]
    \centering
    \includegraphics[scale=0.35]{images/vicuna_cnn_figure1.png}
    \caption{
    The tradeoff between watermark strength and sampling efficiency applies to the algorithms in~\citep{hu2024inevitable}: Vanilla Speculative Sampling (VSpS), Maintain Watermark Strength (MWS), and Maintain Sampling Efficiency (MSE), but \textbf{not} to our algorithm based on the PFR. 
    Even only using the single draft version, our algorithm achieves better performance on both the Average Accepted Tokens Per Step (AATPS) and the watermark score, and attaching watermark does not decrease AATPS. 
    The figure is from using Vicuna-7B-v1.5 and 68m as the model pair on dataset \textsc{eli5}~\citep{fan2019eli5}.
    }
    \label{fig::temp2}
\end{figure}


\begin{figure}[htpb]
    \centering
    \includegraphics[scale=0.35]{images/vicuna_eli5_figure1.png}
    \caption{
    The tradeoff between watermark strength and sampling efficiency applies to the algorithms in~\citep{hu2024inevitable}: Vanilla Speculative Sampling (VSpS), Maintain Watermark Strength (MWS), and Maintain Sampling Efficiency (MSE), but \textbf{not} to our algorithm based on the PFR. 
    Even only using the single draft version, our algorithm achieves better performance on both the Average Accepted Tokens Per Step (AATPS) and the watermark score, and attaching watermark does not decrease AATPS. 
    The figure is from using Qwen2.5-7B-Instruct as the target model and Qwen2.5-0.5B-Instruct as the draft model on dataset \textsc{cnn\_dailymail} dataset.
    }
    \label{fig::temp2}
\end{figure}

\subsubsection{Tables}


\begin{table*}[!htbp]
\centering
\caption{Single-draft results with across-prompt standard deviations. Each cell reports mean $\pm$ standard deviation across 1000 prompts. The figure uses Qwen2.5-7B-Instruct as the target model and Qwen2.5-0.5B-Instruct as the draft model on the \textsc{cnn\_dailymail} dataset. }
\label{tab:single_draft_watermark_std}
\resizebox{\textwidth}{!}{
\begin{tabular}{llcccccc}
\toprule
$L$ & Method & AATPS & Token rate & ANLPPT-$U$ & ANLPPT-Li & ANLPPT-PL & LPPL \\
\midrule
1 & \texttt{basic\_uwm L=1} & $1.000 \pm 0.000$ & $33.353 \pm 2.815$ & $0.048 \pm 0.035$ & $0.035 \pm 0.029$ & $0.051 \pm 0.036$ & $0.473 \pm 0.134$ \\
 & \texttt{VSps} & $1.641 \pm 0.070$ & $33.229 \pm 1.944$ & $0.003 \pm 0.005$ & $0.003 \pm 0.006$ & $0.003 \pm 0.006$ & $0.477 \pm 0.136$ \\
 & \texttt{pfr\_no\_watermark} & $1.623 \pm 0.075$ & $32.147 \pm 2.161$ & $0.002 \pm 0.006$ & $0.003 \pm 0.006$ & $0.002 \pm 0.006$ & $0.474 \pm 0.140$ \\
 & \texttt{pfr} & $1.625 \pm 0.072$ & $33.321 \pm 1.706$ & $0.051 \pm 0.036$ & $0.036 \pm 0.030$ & $0.053 \pm 0.038$ & $0.470 \pm 0.143$ \\
 & \texttt{MSE} & $1.642 \pm 0.070$ & $20.578 \pm 2.645$ & $0.023 \pm 0.021$ & $0.016 \pm 0.017$ & $0.025 \pm 0.023$ & $0.472 \pm 0.136$ \\
 & \texttt{MWS} & $1.621 \pm 0.072$ & $20.136 \pm 2.952$ & $0.049 \pm 0.035$ & $0.034 \pm 0.029$ & $0.052 \pm 0.036$ & $0.475 \pm 0.133$ \\
 & \texttt{MSE\_PSEUDO} & $1.641 \pm 0.072$ & $20.703 \pm 2.420$ & $0.033 \pm 0.028$ & $0.023 \pm 0.023$ & $0.035 \pm 0.029$ & $0.475 \pm 0.138$ \\
\midrule
2 & \texttt{basic\_uwm L=1} & $1.000 \pm 0.000$ & $33.353 \pm 2.815$ & $0.048 \pm 0.035$ & $0.035 \pm 0.029$ & $0.051 \pm 0.036$ & $0.473 \pm 0.134$ \\
 & \texttt{VSps} & $2.055 \pm 0.147$ & $29.399 \pm 2.693$ & $0.002 \pm 0.006$ & $0.002 \pm 0.006$ & $0.002 \pm 0.006$ & $0.482 \pm 0.139$ \\
 & \texttt{pfr\_no\_watermark} & $2.014 \pm 0.151$ & $30.119 \pm 2.438$ & $0.002 \pm 0.005$ & $0.003 \pm 0.005$ & $0.002 \pm 0.005$ & $0.474 \pm 0.140$ \\
 & \texttt{pfr} & $2.016 \pm 0.154$ & $29.746 \pm 2.374$ & $0.051 \pm 0.036$ & $0.036 \pm 0.030$ & $0.053 \pm 0.038$ & $0.471 \pm 0.143$ \\
 & \texttt{MSE} & $2.056 \pm 0.149$ & $16.954 \pm 2.182$ & $0.016 \pm 0.017$ & $0.012 \pm 0.014$ & $0.018 \pm 0.019$ & $0.475 \pm 0.139$ \\
 & \texttt{MWS} & $2.012 \pm 0.155$ & $17.915 \pm 2.679$ & $0.048 \pm 0.035$ & $0.034 \pm 0.028$ & $0.052 \pm 0.036$ & $0.473 \pm 0.134$ \\
 & \texttt{MSE\_PSEUDO} & $2.053 \pm 0.150$ & $17.768 \pm 2.415$ & $0.030 \pm 0.026$ & $0.022 \pm 0.021$ & $0.033 \pm 0.028$ & $0.481 \pm 0.143$ \\
\midrule
3 & \texttt{basic\_uwm L=1} & $1.000 \pm 0.000$ & $33.353 \pm 2.815$ & $0.048 \pm 0.035$ & $0.035 \pm 0.029$ & $0.051 \pm 0.036$ & $0.473 \pm 0.134$ \\
 & \texttt{VSps} & $2.327 \pm 0.228$ & $26.272 \pm 2.968$ & $0.002 \pm 0.006$ & $0.002 \pm 0.006$ & $0.003 \pm 0.006$ & $0.481 \pm 0.141$ \\
 & \texttt{pfr\_no\_watermark} & $2.270 \pm 0.228$ & $26.709 \pm 2.682$ & $0.002 \pm 0.006$ & $0.002 \pm 0.006$ & $0.002 \pm 0.005$ & $0.474 \pm 0.140$ \\
 & \texttt{pfr} & $2.272 \pm 0.225$ & $25.999 \pm 2.742$ & $0.051 \pm 0.036$ & $0.037 \pm 0.030$ & $0.053 \pm 0.038$ & $0.472 \pm 0.144$ \\
 & \texttt{MSE} & $2.330 \pm 0.228$ & $13.453 \pm 1.603$ & $0.015 \pm 0.017$ & $0.011 \pm 0.014$ & $0.017 \pm 0.018$ & $0.476 \pm 0.140$ \\
 & \texttt{MWS} & $2.262 \pm 0.222$ & $16.187 \pm 2.509$ & $0.048 \pm 0.035$ & $0.034 \pm 0.028$ & $0.051 \pm 0.036$ & $0.474 \pm 0.134$ \\
 & \texttt{MSE\_PSEUDO} & $2.337 \pm 0.232$ & $14.933 \pm 2.118$ & $0.029 \pm 0.026$ & $0.021 \pm 0.022$ & $0.032 \pm 0.027$ & $0.480 \pm 0.138$ \\
\midrule
4 & \texttt{basic\_uwm L=1} & $1.000 \pm 0.000$ & $33.353 \pm 2.815$ & $0.048 \pm 0.035$ & $0.035 \pm 0.029$ & $0.051 \pm 0.036$ & $0.473 \pm 0.134$ \\
 & \texttt{VSps} & $2.513 \pm 0.288$ & $23.229 \pm 3.176$ & $0.002 \pm 0.006$ & $0.002 \pm 0.005$ & $0.002 \pm 0.006$ & $0.475 \pm 0.134$ \\
 & \texttt{pfr\_no\_watermark} & $2.439 \pm 0.292$ & $23.285 \pm 2.840$ & $0.002 \pm 0.005$ & $0.002 \pm 0.005$ & $0.002 \pm 0.006$ & $0.474 \pm 0.142$ \\
 & \texttt{pfr} & $2.446 \pm 0.286$ & $23.075 \pm 2.755$ & $0.050 \pm 0.036$ & $0.036 \pm 0.030$ & $0.053 \pm 0.038$ & $0.471 \pm 0.144$ \\
 & \texttt{MSE} & $2.525 \pm 0.298$ & $12.623 \pm 1.893$ & $0.014 \pm 0.016$ & $0.010 \pm 0.013$ & $0.016 \pm 0.017$ & $0.479 \pm 0.145$ \\
 & \texttt{MWS} & $2.432 \pm 0.287$ & $13.592 \pm 2.213$ & $0.048 \pm 0.035$ & $0.034 \pm 0.028$ & $0.052 \pm 0.036$ & $0.475 \pm 0.132$ \\
 & \texttt{MSE\_PSEUDO} & $2.499 \pm 0.287$ & $13.494 \pm 1.934$ & $0.029 \pm 0.026$ & $0.021 \pm 0.021$ & $0.031 \pm 0.027$ & $0.481 \pm 0.141$ \\
\bottomrule
\end{tabular}
}
\end{table*}


\begin{table*}[!htbp]
\centering
\caption{Single-draft results with across-prompt standard deviations. Each cell reports mean $\pm$ standard deviation across 1000 prompts. The results use Qwen2.5-7B-Instruct as the target model and Qwen2.5-0.5B-Instruct as the draft model on the \textsc{eli5} dataset.}
\label{tab:single_draft_qwen_eli5_watermark_std}
\resizebox{\textwidth}{!}{
\begin{tabular}{llcccccc}
\toprule
$L$ & Method & AATPS & Token rate & ANLPPT-$U$ & ANLPPT-Li & ANLPPT-PL & LPPL \\
\midrule
1 & \texttt{basic\_uwm L=1} & $1.000 \pm 0.000$ & $26.132 \pm 1.924$ & $0.100 \pm 0.045$ & $0.072 \pm 0.039$ & $0.104 \pm 0.046$ & $0.696 \pm 0.141$ \\
 & \texttt{VSps} & $1.593 \pm 0.060$ & $32.090 \pm 1.889$ & $0.002 \pm 0.005$ & $0.002 \pm 0.005$ & $0.002 \pm 0.005$ & $0.701 \pm 0.133$ \\
 & \texttt{pfr\_no\_watermark} & $1.567 \pm 0.061$ & $33.046 \pm 1.405$ & $0.002 \pm 0.004$ & $0.002 \pm 0.004$ & $0.002 \pm 0.004$ & $0.703 \pm 0.139$ \\
 & \texttt{pfr} & $1.570 \pm 0.062$ & $33.213 \pm 1.487$ & $0.104 \pm 0.048$ & $0.075 \pm 0.039$ & $0.107 \pm 0.049$ & $0.699 \pm 0.134$ \\
 & \texttt{MSE} & $1.595 \pm 0.060$ & $18.218 \pm 2.022$ & $0.048 \pm 0.030$ & $0.032 \pm 0.024$ & $0.053 \pm 0.032$ & $0.704 \pm 0.135$ \\
 & \texttt{MWS} & $1.565 \pm 0.063$ & $20.177 \pm 2.688$ & $0.102 \pm 0.047$ & $0.075 \pm 0.040$ & $0.106 \pm 0.048$ & $0.700 \pm 0.137$ \\
 & \texttt{MSE\_PSEUDO} & $1.593 \pm 0.062$ & $19.553 \pm 2.107$ & $0.068 \pm 0.038$ & $0.049 \pm 0.032$ & $0.071 \pm 0.039$ & $0.703 \pm 0.145$ \\
\midrule
2 & \texttt{basic\_uwm L=1} & $1.000 \pm 0.000$ & $26.132 \pm 1.924$ & $0.100 \pm 0.045$ & $0.072 \pm 0.039$ & $0.104 \pm 0.046$ & $0.696 \pm 0.141$ \\
 & \texttt{VSps} & $1.948 \pm 0.116$ & $27.350 \pm 2.434$ & $0.002 \pm 0.005$ & $0.002 \pm 0.005$ & $0.002 \pm 0.004$ & $0.701 \pm 0.136$ \\
 & \texttt{pfr\_no\_watermark} & $1.890 \pm 0.118$ & $28.719 \pm 1.782$ & $0.002 \pm 0.005$ & $0.002 \pm 0.004$ & $0.002 \pm 0.004$ & $0.698 \pm 0.137$ \\
 & \texttt{pfr} & $1.895 \pm 0.125$ & $28.851 \pm 1.937$ & $0.104 \pm 0.048$ & $0.075 \pm 0.039$ & $0.107 \pm 0.049$ & $0.699 \pm 0.134$ \\
 & \texttt{MSE} & $1.956 \pm 0.126$ & $14.867 \pm 1.179$ & $0.036 \pm 0.024$ & $0.024 \pm 0.020$ & $0.040 \pm 0.026$ & $0.704 \pm 0.140$ \\
 & \texttt{MWS} & $1.888 \pm 0.121$ & $16.985 \pm 2.208$ & $0.102 \pm 0.046$ & $0.074 \pm 0.039$ & $0.105 \pm 0.047$ & $0.700 \pm 0.142$ \\
 & \texttt{MSE\_PSEUDO} & $1.951 \pm 0.118$ & $24.767 \pm 2.016$ & $0.062 \pm 0.034$ & $0.043 \pm 0.028$ & $0.066 \pm 0.036$ & $0.706 \pm 0.136$ \\
\midrule
3 & \texttt{basic\_uwm L=1} & $1.000 \pm 0.000$ & $26.132 \pm 1.924$ & $0.100 \pm 0.045$ & $0.072 \pm 0.039$ & $0.104 \pm 0.046$ & $0.696 \pm 0.141$ \\
 & \texttt{VSps} & $2.159 \pm 0.175$ & $23.619 \pm 2.707$ & $0.002 \pm 0.004$ & $0.002 \pm 0.004$ & $0.002 \pm 0.003$ & $0.706 \pm 0.139$ \\
 & \texttt{pfr\_no\_watermark} & $2.074 \pm 0.167$ & $24.315 \pm 2.065$ & $0.002 \pm 0.004$ & $0.002 \pm 0.005$ & $0.002 \pm 0.004$ & $0.710 \pm 0.139$ \\
 & \texttt{pfr} & $2.084 \pm 0.174$ & $24.675 \pm 2.059$ & $0.104 \pm 0.049$ & $0.075 \pm 0.041$ & $0.107 \pm 0.049$ & $0.697 \pm 0.135$ \\
 & \texttt{MSE} & $2.164 \pm 0.175$ & $13.359 \pm 2.005$ & $0.030 \pm 0.023$ & $0.020 \pm 0.019$ & $0.035 \pm 0.025$ & $0.705 \pm 0.135$ \\
 & \texttt{MWS} & $2.071 \pm 0.169$ & $14.243 \pm 1.924$ & $0.100 \pm 0.046$ & $0.073 \pm 0.040$ & $0.104 \pm 0.047$ & $0.698 \pm 0.138$ \\
 & \texttt{MSE\_PSEUDO} & $2.169 \pm 0.172$ & $21.339 \pm 2.086$ & $0.061 \pm 0.034$ & $0.043 \pm 0.028$ & $0.064 \pm 0.035$ & $0.709 \pm 0.133$ \\
\midrule
4 & \texttt{basic\_uwm L=1} & $1.000 \pm 0.000$ & $26.132 \pm 1.924$ & $0.100 \pm 0.045$ & $0.072 \pm 0.039$ & $0.104 \pm 0.046$ & $0.696 \pm 0.141$ \\
 & \texttt{VSps} & $2.294 \pm 0.220$ & $17.142 \pm 3.485$ & $0.002 \pm 0.005$ & $0.002 \pm 0.005$ & $0.002 \pm 0.004$ & $0.705 \pm 0.137$ \\
 & \texttt{pfr\_no\_watermark} & $2.187 \pm 0.204$ & $21.512 \pm 2.145$ & $0.002 \pm 0.004$ & $0.002 \pm 0.005$ & $0.002 \pm 0.004$ & $0.703 \pm 0.138$ \\
 & \texttt{pfr} & $2.196 \pm 0.214$ & $21.319 \pm 2.212$ & $0.104 \pm 0.048$ & $0.075 \pm 0.041$ & $0.107 \pm 0.049$ & $0.698 \pm 0.134$ \\
 & \texttt{MSE} & $2.292 \pm 0.220$ & $13.200 \pm 2.021$ & $0.029 \pm 0.021$ & $0.019 \pm 0.017$ & $0.033 \pm 0.024$ & $0.701 \pm 0.136$ \\
 & \texttt{MWS} & $2.184 \pm 0.208$ & $11.570 \pm 1.496$ & $0.101 \pm 0.046$ & $0.073 \pm 0.039$ & $0.105 \pm 0.047$ & $0.696 \pm 0.141$ \\
 & \texttt{MSE\_PSEUDO} & $2.302 \pm 0.214$ & $18.298 \pm 1.942$ & $0.060 \pm 0.033$ & $0.042 \pm 0.028$ & $0.063 \pm 0.035$ & $0.705 \pm 0.132$ \\
\bottomrule
\end{tabular}
}
\end{table*}

\begin{table*}[!htbp]
\centering
\caption{Single-draft results with across-prompt standard deviations. Each cell reports mean $\pm$ standard deviation across 1000 prompts. The results use Vicuna-7B-v1.5 as the target model and a Vicuna-68m drafter on the \textsc{cnn\_dailymail} dataset.}
\label{tab:single_draft_vicuna_cnn_watermark_std}
\resizebox{\textwidth}{!}{
\begin{tabular}{llcccccc}
\toprule
$L$ & Method & AATPS & Token rate & ANLPPT-$U$ & ANLPPT-Li & ANLPPT-PL & LPPL \\
\midrule
1 & \texttt{basic\_uwm L=1} & $1.000 \pm 0.000$ & $37.058 \pm 0.848$ & $0.021 \pm 0.024$ & $0.016 \pm 0.019$ & $0.022 \pm 0.025$ & $0.274 \pm 0.190$ \\
 & \texttt{VSps} & $1.557 \pm 0.086$ & $48.890 \pm 3.439$ & $0.002 \pm 0.006$ & $0.003 \pm 0.006$ & $0.002 \pm 0.006$ & $0.275 \pm 0.167$ \\
 & \texttt{pfr\_no\_watermark} & $1.544 \pm 0.088$ & $37.027 \pm 5.064$ & $0.002 \pm 0.006$ & $0.003 \pm 0.006$ & $0.002 \pm 0.005$ & $0.269 \pm 0.121$ \\
 & \texttt{pfr} & $1.550 \pm 0.090$ & $38.256 \pm 6.092$ & $0.022 \pm 0.024$ & $0.017 \pm 0.019$ & $0.023 \pm 0.025$ & $0.272 \pm 0.128$ \\
 & \texttt{MSE} & $1.557 \pm 0.088$ & $45.513 \pm 4.490$ & $0.009 \pm 0.013$ & $0.007 \pm 0.011$ & $0.010 \pm 0.015$ & $0.269 \pm 0.123$ \\
 & \texttt{MWS} & $1.547 \pm 0.091$ & $30.814 \pm 4.078$ & $0.022 \pm 0.024$ & $0.017 \pm 0.020$ & $0.023 \pm 0.026$ & $0.276 \pm 0.192$ \\
 & \texttt{MSE\_PSEUDO} & $1.555 \pm 0.088$ & $27.597 \pm 2.405$ & $0.015 \pm 0.018$ & $0.012 \pm 0.016$ & $0.016 \pm 0.018$ & $0.279 \pm 0.118$ \\
\midrule
2 & \texttt{VSps} & $1.888 \pm 0.179$ & $53.865 \pm 5.770$ & $0.002 \pm 0.006$ & $0.002 \pm 0.006$ & $0.002 \pm 0.006$ & $0.268 \pm 0.129$ \\
 & \texttt{pfr\_no\_watermark} & $1.857 \pm 0.178$ & $40.154 \pm 5.953$ & $0.003 \pm 0.005$ & $0.003 \pm 0.006$ & $0.003 \pm 0.005$ & $0.273 \pm 0.123$ \\
 & \texttt{pfr} & $1.866 \pm 0.183$ & $42.639 \pm 7.272$ & $0.022 \pm 0.024$ & $0.017 \pm 0.019$ & $0.023 \pm 0.025$ & $0.272 \pm 0.127$ \\
 & \texttt{MSE} & $1.876 \pm 0.186$ & $46.276 \pm 7.794$ & $0.007 \pm 0.010$ & $0.006 \pm 0.009$ & $0.007 \pm 0.011$ & $0.268 \pm 0.131$ \\
 & \texttt{MWS} & $1.858 \pm 0.181$ & $31.588 \pm 5.453$ & $0.022 \pm 0.024$ & $0.016 \pm 0.020$ & $0.023 \pm 0.025$ & $0.278 \pm 0.193$ \\
 & \texttt{MSE\_PSEUDO} & $1.877 \pm 0.182$ & $29.023 \pm 4.118$ & $0.015 \pm 0.018$ & $0.012 \pm 0.015$ & $0.015 \pm 0.020$ & $0.273 \pm 0.127$ \\
\midrule
3  & \texttt{VSps} & $2.076 \pm 0.254$ & $53.835 \pm 7.463$ & $0.002 \pm 0.006$ & $0.003 \pm 0.006$ & $0.002 \pm 0.005$ & $0.276 \pm 0.212$ \\
 & \texttt{pfr\_no\_watermark} & $2.044 \pm 0.259$ & $42.513 \pm 7.326$ & $0.003 \pm 0.006$ & $0.003 \pm 0.006$ & $0.002 \pm 0.006$ & $0.274 \pm 0.128$ \\
 & \texttt{pfr} & $2.056 \pm 0.263$ & $41.813 \pm 7.586$ & $0.022 \pm 0.024$ & $0.017 \pm 0.019$ & $0.023 \pm 0.025$ & $0.272 \pm 0.127$ \\
 & \texttt{MSE} & $2.076 \pm 0.260$ & $38.713 \pm 8.338$ & $0.006 \pm 0.009$ & $0.005 \pm 0.008$ & $0.006 \pm 0.010$ & $0.266 \pm 0.118$ \\
 & \texttt{MWS} & $2.043 \pm 0.258$ & $28.960 \pm 4.759$ & $0.021 \pm 0.024$ & $0.016 \pm 0.020$ & $0.022 \pm 0.026$ & $0.277 \pm 0.191$ \\
 & \texttt{MSE\_PSEUDO} & $2.062 \pm 0.257$ & $32.256 \pm 6.945$ & $0.015 \pm 0.019$ & $0.012 \pm 0.016$ & $0.015 \pm 0.020$ & $0.277 \pm 0.125$ \\
\midrule
4 & \texttt{VSps} & $2.194 \pm 0.324$ & $53.940 \pm 8.237$ & $0.002 \pm 0.005$ & $0.002 \pm 0.005$ & $0.002 \pm 0.005$ & $0.271 \pm 0.120$ \\
 & \texttt{pfr\_no\_watermark} & $2.159 \pm 0.326$ & $40.083 \pm 7.877$ & $0.002 \pm 0.005$ & $0.002 \pm 0.005$ & $0.002 \pm 0.005$ & $0.268 \pm 0.117$ \\
 & \texttt{pfr} & $2.177 \pm 0.323$ & $41.800 \pm 9.355$ & $0.022 \pm 0.023$ & $0.017 \pm 0.019$ & $0.023 \pm 0.025$ & $0.272 \pm 0.127$ \\
 & \texttt{MSE} & $2.220 \pm 0.330$ & $32.730 \pm 6.784$ & $0.005 \pm 0.008$ & $0.005 \pm 0.008$ & $0.006 \pm 0.009$ & $0.264 \pm 0.119$ \\
 & \texttt{MWS} & $2.163 \pm 0.319$ & $27.253 \pm 4.410$ & $0.021 \pm 0.024$ & $0.016 \pm 0.019$ & $0.023 \pm 0.025$ & $0.277 \pm 0.191$ \\
 & \texttt{MSE\_PSEUDO} & $2.183 \pm 0.317$ & $33.007 \pm 6.806$ & $0.014 \pm 0.018$ & $0.011 \pm 0.015$ & $0.014 \pm 0.018$ & $0.277 \pm 0.121$ \\
\bottomrule
\end{tabular}
}
\end{table*}

\begin{table*}[!htbp]
\centering
\caption{Single-draft results with across-prompt standard deviations. Each cell reports mean $\pm$ standard deviation across 1000 prompts. The results use Vicuna-7B-v1.5 as the target model and a Vicuna-68m drafter on the \textsc{eli5} dataset.}
\label{tab:single_draft_vicuna_eli5_watermark_std}
\resizebox{\textwidth}{!}{
\begin{tabular}{llcccccc}
\toprule
$L$ & Method & AATPS & Token rate & ANLPPT-$U$ & ANLPPT-Li & ANLPPT-PL & LPPL \\
\midrule
1 & \texttt{basic\_uwm L=1} & $1.000 \pm 0.000$ & $38.536 \pm 1.896$ & $0.060 \pm 0.040$ & $0.044 \pm 0.033$ & $0.062 \pm 0.042$ & $0.507 \pm 0.152$ \\
 & \texttt{VSps} & $1.410 \pm 0.063$ & $46.306 \pm 3.178$ & $0.002 \pm 0.005$ & $0.003 \pm 0.005$ & $0.002 \pm 0.005$ & $0.511 \pm 0.150$ \\
 & \texttt{pfr\_no\_watermark} & $1.399 \pm 0.063$ & $33.636 \pm 4.757$ & $0.003 \pm 0.006$ & $0.003 \pm 0.007$ & $0.002 \pm 0.005$ & $0.509 \pm 0.152$ \\
 & \texttt{pfr} & $1.397 \pm 0.065$ & $35.231 \pm 5.397$ & $0.061 \pm 0.043$ & $0.045 \pm 0.038$ & $0.063 \pm 0.042$ & $0.513 \pm 0.162$ \\
 & \texttt{MSE} & $1.415 \pm 0.062$ & $44.719 \pm 3.192$ & $0.020 \pm 0.020$ & $0.014 \pm 0.017$ & $0.022 \pm 0.022$ & $0.509 \pm 0.147$ \\
 & \texttt{MWS} & $1.399 \pm 0.064$ & $26.309 \pm 3.816$ & $0.060 \pm 0.041$ & $0.044 \pm 0.034$ & $0.062 \pm 0.042$ & $0.507 \pm 0.151$ \\
 & \texttt{MSE\_PSEUDO} & $1.416 \pm 0.063$ & $28.987 \pm 4.863$ & $0.040 \pm 0.030$ & $0.029 \pm 0.025$ & $0.042 \pm 0.031$ & $0.504 \pm 0.152$ \\
\midrule
2  & \texttt{VSps} & $1.589 \pm 0.104$ & $47.569 \pm 3.837$ & $0.003 \pm 0.006$ & $0.003 \pm 0.006$ & $0.003 \pm 0.006$ & $0.509 \pm 0.146$ \\
 & \texttt{pfr\_no\_watermark} & $1.562 \pm 0.103$ & $34.678 \pm 4.452$ & $0.002 \pm 0.005$ & $0.002 \pm 0.005$ & $0.002 \pm 0.006$ & $0.505 \pm 0.147$ \\
 & \texttt{pfr} & $1.558 \pm 0.106$ & $34.425 \pm 4.996$ & $0.061 \pm 0.043$ & $0.044 \pm 0.038$ & $0.063 \pm 0.042$ & $0.513 \pm 0.162$ \\
 & \texttt{MSE} & $1.587 \pm 0.105$ & $33.648 \pm 4.994$ & $0.012 \pm 0.014$ & $0.009 \pm 0.012$ & $0.014 \pm 0.016$ & $0.509 \pm 0.147$ \\
 & \texttt{MWS} & $1.560 \pm 0.109$ & $23.481 \pm 2.014$ & $0.060 \pm 0.040$ & $0.043 \pm 0.034$ & $0.062 \pm 0.042$ & $0.510 \pm 0.157$ \\
 & \texttt{MSE\_PSEUDO} & $1.593 \pm 0.110$ & $28.595 \pm 4.733$ & $0.040 \pm 0.032$ & $0.029 \pm 0.027$ & $0.041 \pm 0.032$ & $0.509 \pm 0.152$ \\
\m& \texttt{VSps} & $1.670 \pm 0.134$ & $45.003 \pm 5.451$ & $0.003 \pm 0.006$ & $0.002 \pm 0.005$ & $0.002 \pm 0.006$ & $0.513 \pm 0.157$ \\
 & \texttt{pfr\_no\_watermark} & $1.624 \pm 0.132$ & $31.214 \pm 4.662$ & $0.002 \pm 0.005$ & $0.002 \pm 0.005$ & $0.002 \pm 0.005$ & $0.506 \pm 0.149$ \\
 & \texttt{pfr} & $1.626 \pm 0.130$ & $35.694 \pm 5.835$ & $0.062 \pm 0.043$ & $0.045 \pm 0.038$ & $0.063 \pm 0.042$ & $0.515 \pm 0.162$ \\
 & \texttt{MSE} & $1.671 \pm 0.136$ & $26.134 \pm 3.956$ & $0.011 \pm 0.014$ & $0.008 \pm 0.012$ & $0.012 \pm 0.015$ & $0.499 \pm 0.145$ \\
 & \texttt{MWS} & $1.626 \pm 0.136$ & $22.536 \pm 2.127$ & $0.059 \pm 0.040$ & $0.043 \pm 0.033$ & $0.061 \pm 0.041$ & $0.508 \pm 0.152$ \\
 & \texttt{MSE\_PSEUDO} & $1.666 \pm 0.134$ & $28.611 \pm 4.748$ & $0.037 \pm 0.028$ & $0.027 \pm 0.023$ & $0.039 \pm 0.029$ & $0.505 \pm 0.148$ \\
\midrule
4  & \texttt{VSps} & $1.697 \pm 0.152$ & $43.287 \pm 5.333$ & $0.002 \pm 0.005$ & $0.002 \pm 0.006$ & $0.002 \pm 0.005$ & $0.510 \pm 0.152$ \\
 & \texttt{pfr\_no\_watermark} & $1.656 \pm 0.145$ & $28.914 \pm 4.147$ & $0.002 \pm 0.005$ & $0.002 \pm 0.005$ & $0.002 \pm 0.005$ & $0.503 \pm 0.149$ \\
 & \texttt{pfr} & $1.656 \pm 0.146$ & $32.316 \pm 5.921$ & $0.062 \pm 0.043$ & $0.045 \pm 0.038$ & $0.063 \pm 0.042$ & $0.515 \pm 0.162$ \\
 & \texttt{MSE} & $1.700 \pm 0.154$ & $23.993 \pm 3.469$ & $0.011 \pm 0.014$ & $0.008 \pm 0.012$ & $0.012 \pm 0.015$ & $0.514 \pm 0.150$ \\
 & \texttt{MWS} & $1.651 \pm 0.149$ & $22.198 \pm 3.679$ & $0.059 \pm 0.039$ & $0.043 \pm 0.033$ & $0.062 \pm 0.041$ & $0.509 \pm 0.153$ \\
 & \texttt{MSE\_PSEUDO} & $1.703 \pm 0.155$ & $26.004 \pm 4.603$ & $0.037 \pm 0.030$ & $0.027 \pm 0.024$ & $0.039 \pm 0.031$ & $0.508 \pm 0.147$ \\
\bottomrule
\end{tabular}
}
\end{table*}




\subsection{Multi-draft experiments: full per-cell results}
\label{sec:multi_draft}

\paragraph{Setup.}
The main-text multi-draft evaluation
(Section~\ref{sec::experiments}, Table~\ref{tab:headline}) reports a
single representative cell.  Here we extend the evaluation to the
$2{\times}2$ matrix of \{Qwen2.5-7B-Instruct, Vicuna-7b-v1.5\}
$\times$ \{CNN/DailyMail, ELi5\}, using target/drafter pairs
\texttt{Qwen2.5-7B-Instruct} / \texttt{Qwen2.5-0.5B-Instruct} and
\texttt{lmsys/vicuna-7b-v1.5} / \texttt{double7/vicuna-68m}.  Each cell
uses $n{=}1000$ prompts, lookahead $L{=}4$, draft counts
$B\in\{2,4,6,8\}$, and $\mathrm{top\_k}=50,\ \mathrm{top\_p}=1.0,\
T_{\mathrm{eval}}=1.0$.

\paragraph{Decoders.}
\textsc{MPFR} denotes our keyed multi-draft Poisson sampler;
\textsc{Invariant} denotes the unwatermarked
\emph{InvariantMultiDraftStrategy} of \citet{rowan2025list} run in the
same harness on the same prompt set.  The two decoders share the
draft/verify mechanism and differ only in whether the per-context
watermark key is applied.

\paragraph{Per-cell tables.}
Tables~\ref{tab:full-qwen-cnn}--\ref{tab:full-vicuna-eli5} report all
quality and watermark-strength metrics per (decoder, $B$).  Columns:
average accepted tokens per step (AATPS); token rate in tokens/sec
(TR); ANLPPT-\{U,\,Li,\,PL\} (per-token recovery scores under the
three score variants); LPPL (per-token log-perplexity of the generated
tokens under the target model).  TPR-vs-$T$ curves are reported as a
figure in Section~\ref{sec:tpr-figure}.

\begin{table}[h]
\centering
\small
\caption{Qwen2.5-7B-Instruct $\times$ CNN/DailyMail.}
\label{tab:full-qwen-cnn}
\begin{tabular}{l c c c c c c c}
\toprule
Decoder & $B$ & AATPS & TR & U & Li & PL & LPPL \\
\midrule
MPFR      & 2 & 2.799 & 23.18 & 0.052 & 0.037 & 0.055 & 0.475 \\
MPFR      & 4 & 3.124 & 24.89 & 0.052 & 0.037 & 0.055 & 0.474 \\
MPFR      & 6 & 3.298 & 25.39 & 0.052 & 0.037 & 0.055 & 0.476 \\
MPFR      & 8 & 3.410 & 25.17 & 0.053 & 0.037 & 0.055 & 0.476 \\
Invariant & 2 & 2.804 & 25.03 & 0.003 & 0.003 & 0.002 & 0.487 \\
Invariant & 4 & 3.114 & 26.63 & 0.002 & 0.002 & 0.002 & 0.481 \\
Invariant & 6 & 3.275 & 26.79 & 0.002 & 0.003 & 0.002 & 0.482 \\
Invariant & 8 & 3.387 & 26.38 & 0.002 & 0.003 & 0.002 & 0.479 \\
\bottomrule
\end{tabular}
\end{table}

\begin{table}[h]
\centering
\small
\caption{Qwen2.5-7B-Instruct $\times$ ELi5.}
\label{tab:full-qwen-eli5}
\begin{tabular}{l c c c c c c c}
\toprule
Decoder & $B$ & AATPS & TR & U & Li & PL & LPPL \\
\midrule
MPFR      & 2 & 2.563 & 22.38 & 0.104 & 0.076 & 0.108 & 0.704 \\
MPFR      & 4 & 2.919 & 24.48 & 0.104 & 0.076 & 0.107 & 0.703 \\
MPFR      & 6 & 3.111 & 25.54 & 0.105 & 0.076 & 0.108 & 0.706 \\
MPFR      & 8 & 3.237 & 25.83 & 0.105 & 0.076 & 0.108 & 0.704 \\
Invariant & 2 & 2.519 & 23.18 & 0.002 & 0.002 & 0.002 & 0.750 \\
Invariant & 4 & 2.852 & 25.75 & 0.002 & 0.002 & 0.002 & 0.733 \\
Invariant & 6 & 3.043 & 27.11 & 0.002 & 0.002 & 0.002 & 0.730 \\
Invariant & 8 & 3.158 & 27.96 & 0.002 & 0.002 & 0.002 & 0.725 \\
\bottomrule
\end{tabular}
\end{table}

\begin{table}[h]
\centering
\small
\caption{Vicuna-7b-v1.5 $\times$ CNN/DailyMail.}
\label{tab:full-vicuna-cnn}
\begin{tabular}{l c c c c c c c}
\toprule
Decoder & $B$ & AATPS & TR & U & Li & PL & LPPL \\
\midrule
MPFR      & 2 & 2.422 & 52.10 & 0.022 & 0.017 & 0.022 & 0.270 \\
MPFR      & 4 & 2.664 & 48.28 & 0.022 & 0.017 & 0.023 & 0.271 \\
MPFR      & 6 & 2.803 & 44.44 & 0.022 & 0.017 & 0.023 & 0.270 \\
MPFR      & 8 & 2.899 & 40.76 & 0.022 & 0.017 & 0.023 & 0.270 \\
Invariant & 2 & 2.478 & 58.15 & 0.002 & 0.002 & 0.002 & 0.287 \\
Invariant & 4 & 2.690 & 53.12 & 0.002 & 0.002 & 0.002 & 0.286 \\
Invariant & 6 & 2.831 & 48.41 & 0.002 & 0.003 & 0.002 & 0.283 \\
Invariant & 8 & 2.925 & 44.01 & 0.002 & 0.002 & 0.002 & 0.288 \\
\bottomrule
\end{tabular}
\end{table}

\begin{table}[!htbp]
\centering
\small
\caption{Vicuna-7b-v1.5 $\times$ ELi5.}
\label{tab:full-vicuna-eli5}
\begin{tabular}{l c c c c c c c}
\toprule
Decoder & $B$ & AATPS & TR & U & Li & PL & LPPL \\
\midrule
MPFR      & 2 & 1.889 & 45.89 & 0.058 & 0.043 & 0.061 & 0.504 \\
MPFR      & 4 & 2.133 & 47.43 & 0.058 & 0.042 & 0.061 & 0.504 \\
MPFR      & 6 & 2.278 & 47.42 & 0.058 & 0.042 & 0.062 & 0.503 \\
MPFR      & 8 & 2.377 & 46.42 & 0.059 & 0.043 & 0.062 & 0.504 \\
Invariant & 2 & 1.884 & 50.29 & 0.002 & 0.002 & 0.002 & 0.555 \\
Invariant & 4 & 2.108 & 53.42 & 0.002 & 0.002 & 0.002 & 0.552 \\
Invariant & 6 & 2.248 & 55.08 & 0.002 & 0.003 & 0.002 & 0.546 \\
Invariant & 8 & 2.361 & 56.00 & 0.003 & 0.003 & 0.002 & 0.539 \\
\bottomrule
\end{tabular}
\end{table}


\paragraph{Speculative-acceptance parity.}
Within each cell, \textsc{MPFR}'s AATPS matches \textsc{Invariant}'s
within $\pm 1\%$ at every $B \in \{2,4,6,8\}$, indicating that the
target-side keyed-race construction adds no per-step acceptance
penalty.  Realized token rate (TR) tracks AATPS up to a fixed
Python-side label-hashing overhead whose relative cost depends on the
target/drafter compute ratio: with the small \texttt{vicuna-68m}
drafter the overhead is a larger fraction of per-step time, and TR is
$\sim\!10\%$ below \textsc{Invariant}, while on the
\texttt{Qwen2.5-7B/0.5B} pair the gap narrows to $\sim\!5\%$.

\paragraph{Stability of the watermark signal across $B$.}
$\mathrm{ANLPPT}_U$, $\mathrm{ANLPPT}_{Li}$, and $\mathrm{ANLPPT}_{PL}$
for \textsc{MPFR} remain within $\pm 0.001$ across $B\in\{2,4,6,8\}$ on
every cell of Tables~\ref{tab:full-qwen-cnn}--\ref{tab:full-vicuna-eli5}.
The per-token detection signal therefore does not dilute with the
number of drafts, in contrast to naive multi-draft watermarking, which
would re-introduce the efficiency--detectability trade-off of
\citet{hu2024inevitable} at the multi-draft level.

\paragraph{Distortion-free audit.}
LPPL of \textsc{MPFR} outputs falls within $\pm 0.013$ nats/token of
\textsc{Invariant} on every cell.  This is consistent with per-token
marginal preservation and confirms that the watermark does not
introduce a measurable shift in target-model log-perplexity in the
multi-draft regime.

\subsection{Prompt templates}\label{app:prompts}
For all instruction-tuned targets we apply each model's native chat
template via the HuggingFace \texttt{apply\_chat\_template} interface;  The Vicuna v1.5 release ships without a
\texttt{chat\_template} attribute, so we explicitly inject the FastChat
v1.1 template before tokenization.

\paragraph{CNN/DailyMail (summarization).}
The message list passed to the chat template is
\begin{verbatim}
[
  {"role": "system",
   "content": "You are a helpful assistant."},
  {"role": "user",
   "content": "Summarize the following article in 3-5 sentences:\n\n"
              + article[:1500]}
]
\end{verbatim}
The article is truncated to 1500 characters before tokenization.

\paragraph{ELI5 (open-ended QA).}
\begin{verbatim}
[
  {"role": "system",
   "content": "You are a helpful assistant."},
  {"role": "user",
   "content": "Please explain like I'm five:\n\n" + question}
]
\end{verbatim}

For each prompt the model is asked to generate up to 128 new tokens
under \texttt{top\_k}\,$=50$, \texttt{top\_p}\,$=1.0$, and temperature
$T=1.0$.  The watermark key is fixed across all runs.



\subsection{Ablation Details}
\label{app:ablation}
\paragraph{Drafter conditions.}
$D_0$: Qwen2.5-0.5B-Instruct at $T=1.0$ (default).
$D_1$: Qwen2.5-1.5B-Instruct at $T=1.0$ ($3\!\times$ scale, model swap).
$D_2$: Qwen2.5-0.5B-Instruct at $T=0.5$ (sharper drafter).
$D_3$: Qwen2.5-0.5B-Instruct at $T=1.5$ (more diffuse drafter).
Target temperature, watermark key, and all other parameters are held
fixed across $D_0$--$D_3$.
\paragraph{Detection statistic.}
For Aaronson Gamma-tail , per-token score
$s_t = -\log(1 - U_t)$ where $U_t$ is the per-token uniform recovered
from the watermark code. Under the no-watermark null, the cumulative
score $S_T = \sum_{t \le T} s_t$ is $\mathrm{Gamma}(T, 1)$; we declare
detection at FPR $=\alpha$ when $\Pr_{H_0}[\mathrm{Gamma}(T,1)\!>\!S_T]
\le \alpha$, computed via the regularised upper incomplete gamma. In this paper, we report as $\alpha =1\%$

\paragraph{Full detection across models and datasets.}
Figure~\ref{fig:full_detection_4cells} extends the default-drafter
detection comparison (Figure~\ref{fig:drafter_invariance}, left) to all
four $(\textsc{model}, \textsc{dataset})$ cells of the main 2$\times$2
result: Qwen2.5-7B and Vicuna-7B targets crossed with
\textsc{cnn\_dailymail} and \textsc{eli5}. In every cell,
\textsc{pfr}/\textsc{mpfr} (ours) coincide with the strong-watermark
ceiling defined by autoregressive \textsc{basic\_uwm} and \textsc{mws},
while \textsc{mse} and \textsc{mse\_pseudo}~\citep{he2026improving} are
uniformly weaker. The size of this gap is dataset-dependent—largest on
\textsc{cnn\_dailymail}, smallest on the longer-form \textsc{eli5} where
every scheme approaches saturation by $T_{\mathrm{eval}}\!=\!128$—but
the ranking is invariant. This confirms that the detection-strength
conclusion drawn from the anchor cell in
Figure~\ref{fig:drafter_invariance} (left) is not specific to a single
$(\textsc{target}, \textsc{dataset})$ pair.

\begin{figure}[h]
\centering
\includegraphics[width=\linewidth]{images/tpr_comparison_4cells.png}
\caption{\textbf{Full detection across $(\textsc{target}, \textsc{dataset})$
cells.}
TPR@1\%FPR (Aaronson Gamma-tail) vs.\ detection token budget
$T_{\mathrm{eval}}$ at the default drafter $D_0$, for the seven schemes
plotted in Figure~\ref{fig:drafter_invariance} (left). Each panel is one
cell of the main 2$\times$2 result: target~$\in$
\{Qwen2.5-7B-Instruct, Vicuna-7B-v1.5\} $\times$ dataset~$\in$
\{\textsc{cnn\_dailymail}, \textsc{eli5}\}. \textsc{pfr}/\textsc{mpfr}
(ours) match the strong-watermark ceiling (\textsc{basic\_uwm},
\textsc{mws}) in every cell, while \textsc{mse} and \textsc{mse\_pseudo}
are uniformly weaker. The empirical $H_0$ (no watermark) floor sits
near the nominal $1\%$ FPR, calibrating the test.}
\label{fig:full_detection_4cells}
\end{figure}


\paragraph{Pairwise ROUGE-L.}
For each prompt $p$ and decoder $d$ we have one realised output token
sequence per drafter condition. We compute ROUGE-L F1 over each pair
$(D_i, D_j)$ of drafter conditions on these token sequences, then
report the per-prompt minimum over the six pairs (worst pair) averaged
across prompts (Table~\ref{tab:exp2_pairwise_rouge}). An unbiased
verify rule that depends only on target logits and key produces
identical sequences across all drafters, so this quantity should
approach $1$.

\begin{table}[h]
\centering
\caption{\textbf{Output quality.}
\textsc{lppl}$\downarrow$ under target.
ROUGE-L$\uparrow$ vs.\ gold \texttt{highlights}.
AATPS$\uparrow$ accepted target tokens per step.}
\label{tab:exp1_quality}
\small
\setlength{\tabcolsep}{4pt}
\resizebox{\textwidth}{!}{%
\begin{tabular}{l ccccccc | cc}
\toprule
 & \multicolumn{7}{c|}{\emph{Single-draft baselines}}
 & \multicolumn{2}{c}{\emph{Multi-draft ($B=4$)}} \\
\cmidrule(lr){2-8} \cmidrule(lr){9-10}
Metric
 & \textsc{vsps}
 & \textsc{pfr\_no\_wm}
 & \textsc{basic\_uwm}
 & \textsc{mse}
 & \textsc{mws}
 & \textsc{mse\_pseudo}
 & \textbf{\textsc{pfr}}
 & \textsc{invariant\_multi}
 & \textbf{\textsc{mpfr}} \\
\midrule
\textsc{lppl}$\downarrow$
 & 0.474 & 0.469 & 0.472 & 0.470 & 0.475 & 0.477 & \textbf{0.465}
 & 0.484 & \textbf{0.474} \\
ROUGE-L$\uparrow$
 & 0.219 & 0.222 & 0.223 & 0.221 & 0.223 & 0.219 & \textbf{0.223}
 & 0.222 & \textbf{0.222} \\
AATPS$\uparrow$
 & 2.51  & 2.46  & 1.00  & 2.53  & 2.43  & 2.52  & \textbf{2.46}
 & 3.12  & \textbf{3.13} \\
\bottomrule
\end{tabular}%
}
\end{table}

\begin{table}[h]
\centering
\caption{\textbf{Output stability under drafter substitution.}
Pairwise ROUGE-L F1 between the realised token sequences across each
$(D_i,D_j)$ pair, per-prompt minimum averaged over prompts. An
unbiased verify rule produces identical sequences across all
drafters, so ROUGE-L $\to 1$.}
\label{tab:exp2_pairwise_rouge}
\small
\setlength{\tabcolsep}{8pt}
\begin{tabular}{lc}
\toprule
Decoder & min-pair mean ROUGE-L \\
\midrule
\textsc{basic\_uwm} (no drafter, vacuous bound) & 1.000 \\
\textbf{\textsc{pfr}} (ours)                    & \textbf{0.990} \\
\textsc{mws}                                    & 0.922 \\
\midrule
\textsc{pfr\_no\_wm} (no wm)                    & 0.436 \\
\textsc{mse\_pseudo}                            & 0.426 \\
\textsc{mse}                                    & 0.406 \\
\textsc{vsps} (no wm)                           & 0.394 \\
\bottomrule
\end{tabular}
\end{table}






\section{More on Watermark}
\label{app::water}
% \textcolor{red}{Note in implementation, to compare to~\citep{hu2024inevitable}, we also need to only include positions whose seed/context code hasn't appeared before, i.e. ``skipping repeats'' just like them so that $\{U_t\}$ are (approximately) independent across $t$. }




\subsection{Gumbel Watermark}
\label{subapp::gumbel_water}

By the Gumbel watermark~\citep{aaronson2023watermarking}: 
\begin{itemize}
    \item If $y_{1:n}$ is \emph{not} generated from the watermarked model, the $\mathrm{TestScore}$ is simply a sum of exponential random variable and hence 
    \[
    \mathbf{E}\left[ \mathrm{TestScore} (y_{1:n}) \right] = n - m.
    \]

    \item If $y_{1:n}$ \emph{is} generated from the watermarked model, we can find 
    \[
    \mathbf{E}\left[ \mathrm{TestScore} (y_{1:n}) \right] \geq (n - m) + \left(\frac{\pi^2}{6} - 1 \right) \sum^n_{t = m + 1} 
    \mathbf{E}\left[
    \mathrm{Entropy}[p_t(\cdot)]
     \right]
    \]
    where $p_t = \mathrm{Softmax}(u_t/T)$ and the entropy is in \emph{nats} for a discrete distribution. 
\end{itemize}

In~\citep{hu2024inevitable}, the likelihood-agnostic $U$ score used for the DeltaGumbel reweighting is 
\[
U_t
:=
r_t(w_t)
\]
and hence the Aaronson’s original per-token score is a deterministic transform of the DeltaGumbel score 
\[
A_t
:=
-\log(1-r_t(w_t))
=
-\log(1-U_t)
\]
and therefore
\[
S_{\mathrm{A}}
=
\sum_{t\in\mathcal{T}} A_t
=
\sum_{t\in\mathcal{T}} -\log(1-U_t),
\]

Moreover, following the $U$-score framework~\citep{hu2024inevitable}, one may instead aggregate the keyed uniforms directly:
\[
S_{\mathrm{U}}
:=
\sum_{t\in\mathcal{T}} U_t.
\]
Since for $U\sim \mathrm{Unif}(0,1)$,
\[
\mathbf{E}[e^{\lambda U}]
=
\frac{e^\lambda-1}{\lambda},
\qquad \lambda\ge 0,
\]
we obtain
\[
\mathbf{E}[e^{\lambda S_{\mathrm{U}}}]
=
\left(\frac{e^\lambda-1}{\lambda}\right)^M.
\]
Therefore, Chernoff's bound yields
\[
p_{\mathrm{U}}
\le
\inf_{\lambda\ge 0}
\exp \left(
M\log\frac{e^\lambda-1}{\lambda}
-\lambda S_{\mathrm{U}}
\right),
\]
and the corresponding ANLPPT is
\[
\mathrm{ANLPPT}_{\mathrm{U}}
:=
-\frac{1}{M}\log p_{\mathrm{U}}.
\]

In both cases, one declares the sequence watermarked when the resulting $p$-value is below a prescribed significance level, or equivalently when the ANLPPT is sufficiently large. 
Crucially, because every realized output token in Algorithm~\ref{alg:single_PFR_Spec_Water_simplified} is the target-side PFR winner under the keyed source associated with its realized prefix, the detector depends only on the final output sequence and the secret key, and does not need access to the draft model $Q$ or the acceptance indicators.



\subsection{Connection to \citep{hu2024inevitable}}
\label{subapp::aaronson_inevitable_discuss}


For each scored position $t\in\mathcal{T}$, let
\[
P_t:=P(\cdot|w_{:t-1})
\]
and let
\[
g_t(v):=-\log(-\log r_t(v)), \qquad v\in\mathcal{V}.
\]
Then $(g_t(v))_{v\in\mathcal{V}}$ are i.i.d. standard Gumbel random variables.
The DeltaGumbel reweighting rule selects
\[
a_t^\star
:=
\arg\max_{v\in\mathcal{V}}
\left\{\log P_t(v)+g_t(v)\right\}.
\]
Equivalently, the watermarked distribution at position $t$ is
\[
Q_t:=R_{g_t}(P_t)=\delta_{a_t^\star}.
\]

The likelihood-agnostic $U$ score used for the DeltaGumbel reweighting is
\[
U_t
:=
\exp(-\exp(-g_t(w_t))).
\]
By the definition of $g_t$, this is exactly
\[
U_t=r_t(w_t).
\]
Therefore, Aaronson's original per-token score is a deterministic transform of the
DeltaGumbel $U$ score:
\[
A_t
:=
-\log(1-r_t(w_t))
=
-\log(1-U_t).
\]
Thus
\[
S_{\mathrm{A}}
=
\sum_{t\in\mathcal{T}} A_t
=
\sum_{t\in\mathcal{T}} -\log(1-U_t).
\]
In particular, if the entire vector $(U_t)_{t\in\mathcal{T}}$ is retained, then Aaronson's score can be
computed exactly from the $U$ scores. Conversely,
\[
U_t=1-\exp(-A_t).
\]
But the aggregate statistic $\sum_{t\in\mathcal{T}} U_t$ is not equivalent to $S_{\mathrm{A}}=\sum_{t\in\mathcal{T}}-\log(1-U_t)$. 





\subsection{Watermark Score by~\citep{li2025statistical}} 


The score of~\citet{li2025statistical} is also computed from the same pivot
$r_t(w_t)$. For a parameter $\Delta$, define
\[
k_\Delta:=\left\lfloor\frac{1}{1-\Delta}\right\rfloor,
\qquad
q_\Delta:=1-k_\Delta(1-\Delta).
\]
Their per-token score is
\[
h^\star_{\mathrm{gum},\Delta}(r)
=
\log\left(
k_\Delta r^{\frac{\Delta}{1-\Delta}}
+
\mathds{1}\{q_\Delta>0\}
r^{\frac{1-q_\Delta}{q_\Delta}}
\right),
\]
and the corresponding statistic is
\[
S_{\mathrm{Li},\Delta}
:=
\sum_{t\in\mathcal{T}}
h^\star_{\mathrm{gum},\Delta}(r_t(w_t)).
\]



We now discuss another watermark score for Aaronson's Gumbel Watermark, proposed by~\citep{li2025statistical}, which is a statistically optimal score function for the Gumbel watermark under a class-dependent efficiency criterion. 
We introduce their result using our notation as follows. 
Their Gumbel pivot is exactly
\[
Y_t^{\mathrm{gum}}
:=
r_t(w_t).
\]
Under the null hypothesis,
\[
Y_t^{\mathrm{gum}}\sim\mathrm{Unif}[0,1].
\]
Under the Gumbel-watermarked alternative, conditional on
\[
P_t:=P(\cdot|w_{:t-1}),
\]
the pivot satisfies
\[
\mathbf{P}_{H_1}\!\left(Y_t^{\mathrm{gum}}\leq r| P_t\right)
=
\sum_{v\in\mathcal{V}}P_t(v)r^{1/P_t(v)},
\qquad r\in[0,1],
\]
where terms with $P_t(v)=0$ are omitted.

Aaronson's original score corresponds to the per-token score function
\[
h_{\mathrm{A}}(r):=-\log(1-r),
\]
so that
\[
S_{\mathrm{A}}
=
\sum_{t\in\mathcal{T}}h_{\mathrm{A}}(r_t(w_t)).
\]
In contrast, Li et al.~\citep{li2025statistical} consider the distribution class
\[
\mathcal{P}_{\Delta}
:=
\left\{
P:\max_{v\in\mathcal{V}}P(v)\leq 1-\Delta
\right\}
\]
and derive the optimal Gumbel score
\[
h^\star_{\mathrm{gum},\Delta}(r)
=
\log\left(
k_\Delta r^{\frac{\Delta}{1-\Delta}}
+
\mathds{1}\{q_\Delta>0\}
r^{\frac{1-q_\Delta}{q_\Delta}}
\right),
\]
where
\[
k_\Delta
:=
\left\lfloor\frac{1}{1-\Delta}\right\rfloor,
\qquad
q_\Delta
:=
1-k_\Delta(1-\Delta).
\]
The corresponding statistic in our notation is
\[
S_{\mathrm{Li},\Delta}
:=
\sum_{t\in\mathcal{T}}
h^\star_{\mathrm{gum},\Delta}(r_t(w_t)).
\]
This score is the log-likelihood ratio associated with the least-favorable distribution
\[
P_\Delta^\star
=
(\underbrace{1-\Delta,\ldots,1-\Delta}_{k_\Delta\text{ times}},
q_\Delta,0,\ldots,0),
\]
with the $q_\Delta$ term omitted when $q_\Delta=0$.
Therefore, $h^\star_{\mathrm{gum},\Delta}$ is not the same as Aaronson's original score $h_{\mathrm{A}}$.
Indeed,
\[
h_{\mathrm{A}}(r)\to\infty
\qquad\text{as }r\uparrow 1,
\]
whereas
\[
h^\star_{\mathrm{gum},\Delta}(r)
\to
\log\left(k_\Delta+\mathds{1}\{q_\Delta>0\}\right)
\qquad\text{as }r\uparrow 1.
\]
Thus, Li et al.'s score and Aaronson's original score are based on the same Gumbel pivot
$r_t(w_t)$, but they lead to different cumulative detection statistics.



\subsection{Watermark Score by~\citep{lattimore2026refined}}


The score of Lattimore~\citep{lattimore2026refined} is the truncated power-law score
\[
h_{\mathrm{PL},\epsilon}(r)
=
\min\left\{
\epsilon^{-1/2},
(1-r)^{-1/2}
\right\}
-(2-\sqrt{\epsilon}),
\]
and the corresponding statistic is
\[
S_{\mathrm{PL},\epsilon}
:=
\sum_{t\in\mathcal{T}}
h_{\mathrm{PL},\epsilon}(r_t(w_t)).
\]
Thus, all four scores are computed from the same keyed Gumbel pivot
$r_t(w_t)$, but they use different per-token transformations.




Other score criteria can also be explored. For example, \citet{lattimore2026refined}
proposed a refined detector for the Gumbel watermark based on a truncated power-law score.
In our notation, the same Gumbel pivot is
\[
r_t(w_t),
\qquad t\in\mathcal{T}.
\]
Instead of Aaronson's original per-token score
\[
h_{\mathrm{A}}(r):=-\log(1-r),
\]
Lattimore~\citep{lattimore2026refined} considers
\[
h_{\mathrm{PL},\epsilon}(r)
=
\min\left\{
\epsilon^{-1/2},
(1-r)^{-1/2}
\right\}
-(2-\sqrt{\epsilon}),
\]
and the corresponding statistic is
\[
S_{\mathrm{PL},\epsilon}
:=
\sum_{t\in\mathcal{T}}
h_{\mathrm{PL},\epsilon}(r_t(w_t)).
\]
Here the centering term $2-\sqrt{\epsilon}$ is chosen so that, under the null hypothesis,
\[
\mathbf{E}\!\left[h_{\mathrm{PL},\epsilon}(U)\right]=0,
\qquad U\sim\mathrm{Unif}[0,1].
\]
Thus, under the null and assuming independence across the scored positions,
\[
\mathbf{E}[S_{\mathrm{PL},\epsilon}]=0.
\]
Both $S_{\mathrm{PL},\epsilon}$ and Aaronson's statistic
\[
S_{\mathrm{A}}
=
\sum_{t\in\mathcal{T}}h_{\mathrm{A}}(r_t(w_t))
\]
are model-agnostic statistics based on the same Gumbel pivot $r_t(w_t)$.
However, they are not equivalent. In particular,
\[
h_{\mathrm{A}}(r)\to\infty
\qquad\text{as }r\uparrow 1,
\]
whereas
\[
h_{\mathrm{PL},\epsilon}(r)
\leq
\epsilon^{-1/2}-(2-\sqrt{\epsilon})
\]
is bounded because of the truncation. Therefore, Lattimore's truncated power-law score and Aaronson's original score use the same keyed randomness, but lead to different cumulative detection statistics. 



\section{Algorithm Implementation}
\label{app::alg_imple}


\begin{algorithm}[tbp]
\SetAlgoLined
\DontPrintSemicolon
\SetKwInOut{Input}{Input}\SetKwInOut{Output}{Output}

\Input{output sequence $w_{1:N}$, prompt length $n_0$, labeler $\Lambda$, secret key $\mathsf{k}$, reference measure $\mu$}
\Output{Aaronson score $S_{\mathrm A}$, $p$-value $p_{\mathrm A}$, ANLPPT $\mathrm{ANLPPT}_{\mathrm A}$ and $M = |\mathcal{T}|$}

$S_{\mathrm A} \gets 0$\;
$M \gets 0$\;

\For{$t=n_0+1$ \KwTo $N$}{
    $c_t \gets w_{:t-1}$\;
    $\ell_t \gets \Lambda(c_t)$\;
    $\Pi_t \gets G(\mathsf{k}, \ell_t)$\;

    \tcp{recover keyed uniform value for the observed token $w_t$}
    $\tau_t(w_t) \gets \textsc{FirstArrival}(w_t,\Pi_t,\mu)$\;
    $r_t \gets \exp(-\mu(w_t)\tau_t(w_t))$\;

    $S_{\mathrm A} \gets S_{\mathrm A} - \log(1-r_t)$\;
    $M \gets M + 1$\;
}

\tcp{Under $H_0$, $S_{\mathrm A}\sim \mathrm{Gamma}(M,1)$}
$p_{\mathrm A} \gets \Gamma(M,S_{\mathrm A})/\Gamma(M)$\;

$\mathrm{ANLPPT}_{\mathrm A} \gets -\frac{1}{M}\log p_{\mathrm A}$\;

\caption{Compute Aaronson score, exact null $p$-value, and ANLPPT for keyed PFR watermark}
\label{alg:pfr_aaronson_anlppt}
\end{algorithm}



% \begin{algorithm}[tbp]
\begin{algorithm}[H]
\SetAlgoLined
\DontPrintSemicolon
\SetKwIF{If}{ElseIf}{Else}{if}{:}{else if}{else}{end}
\SetKwInOut{Input}{Input}\SetKwInOut{Output}{Output}

\Input{lookahead $K$, output length $N$, target model family $P$, draft model family $Q$,
initial prefix $w_{1:n}$, common reference measure $\mu$ on $\mathcal{Z}$ such that
$P(\cdot | c),Q(\cdot | c)\ll\mu$ for all contexts $c$,
labeler $\Lambda$, secret key $\mathsf{k}$}

\While{$n < N$}{
    $h \gets w_{1:n}$ \tcp*[r]{freeze block-start prefix}
    $L \gets \min\{K,\,N-n\}$\;

    \For{$s=1$ \KwTo $L$ \tcp*[r]{draft $L$ tokens from shared PFR sources}}{
        $c_s \gets h \,\|\, \tilde w_{1:s-1}$\;
        $\ell_s \gets \Lambda(c_s)$ \tcp*[r]{includes repeated-context masking if used}
        $\Pi_s \gets G(\mathsf{k}, \ell_s)$\;
        $\tilde w_s \gets \textsc{PFR}(Q(\cdot | c_s), \Pi_s, \mu)$\;
    }

    \If{$n+L < N$}{
        $c_{L+1} \gets h \,\|\, \tilde w_{1:L}$\;
    }

    \tcp{Prepare target-side evaluations for $c_1,\ldots,c_L$ and, if needed, $c_{L+1}$ in parallel}
    $accepted \gets \textbf{true}$\;

    \tcp{Verify each draft token using the same shared Poisson source}
    \For{$s=1$ \KwTo $L$}{
        $y_s \gets \textsc{PFR}(P(\cdot | c_s), \Pi_s, \mu)$\;
        $w_{n+1} \gets y_s$; $n \gets n+1$ \tcp*[r]{emit the target PFR winner}
        \If{$y_s \neq \tilde w_s$}{
            $accepted \gets \textbf{false}$\;
            \textbf{break}\; \tcp*[r]{correction: later draft contexts are invalid}
        }
    }

    \If{$accepted$ \textbf{and} $n < N$}{
        $\ell_{L+1} \gets \Lambda(c_{L+1})$\;
        $\Pi_{L+1} \gets G(\mathsf{k}, \ell_{L+1})$\;
        $w_{n+1} \gets \textsc{PFR}(P(\cdot | c_{L+1}), \Pi_{L+1}, \mu)$\;
        $n \gets n+1$ \tcp*[r]{bonus token}
    }
}
\caption{Watermarkable Speculative Sampling via Poisson Processes}
\label{alg:fwss_pfr}
\end{algorithm}



\begin{algorithm}[H]
\SetAlgoLined
\DontPrintSemicolon
\SetAlgoNoEnd
\SetKwIF{If}{ElseIf}{Else}{if}{:}{else if}{else}{end}
\SetKwInOut{Input}{Given}

\Input{lookahead $L$, number of drafts $B$, output length $N$, target model $P$,
draft model $Q$, initial prompt $w_{1:n}$, keyed Poisson source generator $G$,
and secret key $\mathsf{k}$}

\While{$n < N$}{
    $h \gets w_{1:n}$; $\ell \gets \min\{L,\,N-n\}$\;
    $\mathcal{C}_0 \gets \{h\}$; $\nu(h) \gets B$
    \tcp*[r]{root context; all $B$ homogeneous drafts start here}

    \For{$s=1$ \KwTo $\ell$}{
        $\mathcal{C}_s \gets \emptyset$\;
        \For{$c \in \mathcal{C}_{s-1}$}{
            $\Pi(c) \gets G(\mathsf{k}, c)$
            \tcp*[r]{context-indexed keyed source shared by draft and target at context $c$}
            \For{$i=1$ \KwTo $\nu(c)$}{
                $\tilde w(c,i) \gets \textsc{MPFR}(Q(\cdot | c), \Pi(c), i)$\;
            }
            $\mathcal{D}(c) \gets \{\tilde w(c,1),\ldots,\tilde w(c,\nu(c))\}$
            \tcp*[r]{distinct next-token proposals from context $c$}
            \For{$u \in \mathcal{D}(c)$}{
                $\nu(c \,\|\, u) \gets
                \big|\{\,i\in\{1,\ldots,\nu(c)\}: \tilde w(c,i)=u\,\}\big|$\;
                \tcp*[r]{how many drafts move to child context $c\|u$}
                $\mathcal{C}_s \gets \mathcal{C}_s \cup \{\,c \,\|\, u\,\}$\;
            }
        }
    }

    \For(\tcp*[f]{compute target-side winner at every internal context}){$c \in \mathcal{C}_0 \cup \cdots \cup \mathcal{C}_{\ell-1}$}{
        $y(c) \gets \textsc{MPFR}(P(\cdot | c), \Pi(c), 1)$\;
        \tcp*[r]{target-side mapped-Poisson winner at context $c$}
    }

    $c^\star \gets h$; $accepted \gets \textbf{true}$\;
    \tcp*[r]{start following the unique realized target path from the root}

    \For{$s=1$ \KwTo $\ell$}{
        $a_s \gets \mathds{1}\!\left\{y(c^\star)\in \mathcal{D}(c^\star)\right\}$\;
        \tcp*[r]{accept if the target winner is among the drafts at the current context}
        $w_{n+1} \gets y(c^\star)$; $n \gets n+1$\;
        \tcp*[r]{emit the target-side token, not a draft token}
        \If{$a_s = 0$}{
            $accepted \gets \textbf{false}$ and \textbf{break}
            \tcp*[r]{first rejection ends the speculative block}
        }
        \If{$n = N$}{
            \textbf{break}
        }
        $c^\star \gets c^\star \,\|\, y(c^\star)$
        \tcp*[r]{move to the next realized context}
    }

    \If(\tcp*[f]{if all $\ell$ steps were accepted, emit one extra target token}){$accepted$ \textbf{and} $n < N$}{
        $\Pi^\star \gets G(\mathsf{k}, c^\star)$\;
        $w_{n+1} \gets \textsc{MPFR}(P(\cdot | c^\star), \Pi^\star, 1)$\;
        $n \gets n+1$ \tcp*[r]{bonus token}
    }
}
\caption{Watermarkable Multi-Draft Speculative Sampling (Formal)}
\label{alg:multi_fwss_pfr_formal}
\end{algorithm}




% {\color{orange}
% \section{Questions to be considered} 

% \begin{itemize}
%     \item Is it possible, and if possible why, that an imporved version of Gumbel watermarking improves our speculative sampling? Related to~\citep{huangwatermarking}? 

%     \item Research question: if advance to multi-draft speculative sampling, what is the watermark rule, detection rule and detectability? Any theoretic promise, generalizing Aaronson's single-draft version? 

%     \item \citep[Theorem 3.2]{he2026improving} gives $\mathsf{WS}(P_\zeta) = H(P)$ if and only if $H(P_\zeta) = 0$, is this \emph{golden formula}? 

%     \item \textcolor{green}{How do we justify that we break the inevitable boundary????}

%     \item \textcolor{green}{How do use theory/experiment to discuss draft-invariance and watermark strength????}

%     \item \textcolor{green}{New paper on Aaronson's score?????}

%     \item \textcolor{green}{Emphasize more on the idea comes from Gumbel Watermark?????}

%     \item \textcolor{green}{Robustness via PML?????}

%     \item Why our sampling and watermarking do NOT conflict? Either \citep{hu2024inevitable}'s proof is wrong in our setup or we use the KL score in \citep{he2026improving}. 

%     \item If we use the KL score, we need to plot our fundamental tradeoff???? 
% \end{itemize}
% }

\newpage
\input{checklist.tex}

\end{document}